% ===================================================================
%  Preamble for the research report.  All section files assume these macros.
% ===================================================================
\documentclass[11pt,a4paper]{ctexart}

\usepackage{amsmath,amssymb,amsthm,mathtools}
\usepackage[margin=2.3cm]{geometry}
\usepackage{booktabs,longtable,array}
\usepackage{enumitem}
\usepackage{xcolor}
\usepackage{fancyhdr}
\usepackage[colorlinks=true,linkcolor=blue!55!black,urlcolor=blue!55!black]{hyperref}

% ---------- shorthands used by every section ----------
\newcommand{\wt}{\widetilde}
\newcommand{\wh}{\widehat}
\newcommand{\R}{\mathbb{R}}
\newcommand{\N}{\mathbb{N}}
\newcommand{\Om}{\Omega}
\newcommand{\dist}{\operatorname{dist}}
\newcommand{\diam}{\operatorname{diam}}
\newcommand{\supp}{\operatorname{supp}}
\newcommand{\esssup}{\operatorname*{ess\,sup}}
\newcommand{\codimA}{\operatorname{codim}_{A}}
\newcommand{\codimAl}{\operatorname{codim}_{A}^{\mathrm{lower}}}
\newcommand{\codimAu}{\operatorname{codim}_{A}^{\mathrm{upper}}}
\newcommand{\dimM}{\dim_{M}}
\newcommand{\Fsp}{F^{s}_{p,q}(\Om,d^{-aq})}
\newcommand{\tiF}{\wt F^{s}_{p,q}(\Om,d^{-aq})}
\newcommand{\hoF}{\wh F^{s}_{p,q}(\Om,d^{-aq})}
\newcommand{\nuA}{\nu_{a}}
\newcommand{\Omt}{\Om_{t}}
\newcommand{\Omr}{\Om_{r}}
\newcommand{\fid}[1]{\textnormal{\footnotesize\ttfamily #1}}
\newcommand{\defterm}[1]{\textbf{#1}}

% provenance note printed under each theorem
\newcommand{\factnote}[3]{%
  \par\noindent{\footnotesize\color{gray}%
  事实 \fid{#1}\ \textbullet\ 作者 \texttt{#2}\ \textbullet\ 前置 #3}\par\vspace{2pt}}

\theoremstyle{plain}
\newtheorem{theorem}{定理}[section]
\newtheorem{lemma}[theorem]{引理}
\newtheorem{proposition}[theorem]{命题}
\newtheorem{corollary}[theorem]{推论}
\theoremstyle{definition}
\newtheorem{definition}[theorem]{定义}
\newtheorem{example}[theorem]{例}
\newtheorem{remark}[theorem]{评注}
\newtheorem{notation}[theorem]{记号}

\setlist[itemize]{leftmargin=1.6em,itemsep=1pt,topsep=2pt}
\setlist[enumerate]{leftmargin=1.9em,itemsep=2pt,topsep=2pt}

\setlength{\emergencystretch}{1.5em}
\pagestyle{fancy}
\fancyhf{}
\lhead{\small 加权 Triebel--Lizorkin 空间中 $C^\infty$ 稠密性}
\rhead{\small 研究报告}
\cfoot{\small\thepage}


\title{\bfseries 加权 Triebel--Lizorkin 空间中 $C^\infty$ 稠密性的若干精确结论\\[4pt]
\large —— 基于 32 条机器验证事实，附第三方复核与修订记录}
\author{Danus 多智能体研究项目 \texttt{cx-density-weighted-tl}}
\date{}

\begin{document}
\maketitle

\begin{abstract}
\noindent
本文研究如下问题：设 $\Om\subset\R^n$ 为有界区域，$d(x)=\dist(x,\partial\Om)$，$a>0$，
$w(y)=d(y)^{-aq}$，问 $C^\infty(\Om)$ 与 $C_c^\infty(\Om)$ 是否在加权内在
Triebel--Lizorkin 空间 $F^s_{p,q}(\Om,d^{-aq})$ 中稠密。

本文的核心发现是一条\defterm{二分现象}：边界距离测度 $\nuA$ 的质量要么有限（情形 (I)），
要么在每个边界层上都已经无穷（情形 (II)），二者必居其一。

在情形 (II) 下，$\Fsp$ 退化为次数不超过 $k=\lfloor s\rfloor$ 的多项式空间，
于是 $C^\infty$ 平凡稠密而 $C_c^\infty\cap\Fsp=\{0\}$；
在情形 (I) 下，$C^\infty$ 稠密（$p=q$、$s\in(0,1)$），
且对 Hardy 加强空间 $\tiF$ 有完整的强化结论
\[
  \wt F^{s,1}_{p,p}(\Om,d^{-ap})=\wt F^{s,2}_{p,p}(\Om,d^{-ap})=\wt F^s_{p,p}(\Om,d^{-ap}).
\]

本文进一步证明该强化\defterm{不能}推广为对 $F^s_{p,q}$ 本身的陈述：
即使在权完全可积的情形 (I)，$\Om=(0,1)$、$s=\tfrac12$、$p=q=2$、$a=\tfrac14$ 时
仍有 $F^{s,1}\neq F^s=F^{s,2}$。
决定 $C_c^\infty$ 稠密性的临界量是 $sp+aq$ 与边界 Assouad 余维数
$\codimA(\partial\Om)$ 的比较，而\defterm{不是} $\int_\Om d^{-aq}<\infty$ 是否成立。
此外，本文证明整数阶修正空间中本质紧支的元素类稠密（截断引理），为 Poincar\'e 障碍所指出的「非整数 $s>1$ 时未加强空间不可稠密」提供一组正面对照。

这一点同时解释了为何 CLWZ 型 Hardy 加强项不是冗余的：
常数函数 $1\in\Fsp$ 但 $1\notin\tiF$，即 $\Fsp\not\subseteq\tiF$ 对一切非整数 $s$、一切 $a>0$、一切 $p,q$ 成立。

\medskip
\noindent\textbf{验证状态。}本文收录的 32 条结论，最初均由 LLM 验证器（\textsf{Verifier}）
判定 \texttt{verdict=correct} 后写入事实图，其中两条支柱事实另经一次\defterm{独立重推}复核。
此后一次\defterm{第三方（另一模型）审核}发现：第 \ref{sec:part2} 节中关于整数阶修正空间非退化的命题，
其原证明含一处发散积分被误判为有限，且其所选 witness $d^k$ 在若干情形下不成立。
本文已将该命题\defterm{改用 $C_c^\infty$ 作为 witness 重新证明}（结论仍然成立），
并把原证明的失效分析与反例作为评注保留（见 \S\ref{sec:part2} 相应评注与 \S\ref{sec:revision}）。
\textbf{这些结论未经形式化验证（无 Lean/Coq），亦未经人类同行评审}；所有渐近常数均未显式化。引用前须人工复核。
\end{abstract}

\tableofcontents
\bigskip

%=====================================================================
\section{引言、设定与基本定义}
%=====================================================================

\subsection{研究对象}

设 $n\ge1$，$\Om\subset\R^n$ 为非空有界域（连通、有界、开），
\[
  d(x):=\dist(x,\partial\Om),\qquad x\in\Om,
\]
为边界距离函数。固定参数
\[
  a>0,\qquad p,q\in[1,\infty),\qquad s\in(0,\infty),
\]
并记权函数
\begin{equation}\label{eq:weight}
  w(y):=d(y)^{-aq}.
\end{equation}

\begin{definition}[边界距离测度与边界层]\label{def:nu}
定义 $\Om$ 上的 Borel 测度
\begin{equation}\label{eq:nu}
  \nuA(E):=\int_E d(y)^{-aq}\,dy=\int_E w(y)\,dy,
\end{equation}
并对 $r>0$ 记边界层
\begin{equation}\label{eq:Omt}
  \Omr:=\{x\in\Om:\ d(x)<r\}.
\end{equation}
\end{definition}

\subsection{加权内在 Triebel--Lizorkin 范数}

沿用 CLWZ 的定义。写 $s=k+\sigma$，其中 $k=\lfloor s\rfloor\in\N_0$、$\sigma\in(0,1)$。
对 $f\in L^p_{\mathrm{loc}}(\Om)$，令
\begin{equation}\label{eq:intrinsic}
  \|f\|_{\Fsp}
  :=\|f\|_{W^{k,p}(\Om)}
  +\sum_{|\beta|=k}
  \Bigg(\int_\Om
    \Big(\int_\Om \big|D^\beta f(x)-D^\beta f(y)\big|^q\,|x-y|^{-n-\sigma q}
          \,d(y)^{-aq}\,dy\Big)^{p/q}dx
  \Bigg)^{1/p},
\end{equation}
各项取 $+\infty$ 时表示该量不有限。相应的空间为
\[
  \Fsp:=\Big\{f \text{ 在 } \Om \text{ 上 Lebesgue 可测}:\ \|f\|_{\Fsp}<\infty\Big\}.
\]

\begin{remark}[内层积分对全域取]\label{rem:fullintegral}
\eqref{eq:intrinsic} 中的内层积分\textbf{对全部 $y\in\Om$} 取，而非截断到球
$|x-y|<c\,d(x)$ 的局部化版本。这一结构是第 \ref{sec:part1} 节刚性引理
（见 \S\ref{sec:part1} 中的 $\sigma$-引理）成立的前提，也是本文多处论证的关键。
\end{remark}

当 $k=0$（即 $s=\sigma\in(0,1)$）时 \eqref{eq:intrinsic} 化为
\begin{equation}\label{eq:intrinsic-k0}
  \|f\|_{\Fsp}
  =\|f\|_{L^p(\Om)}
  +\Big(\int_\Om\Big(\int_\Om |f(x)-f(y)|^q|x-y|^{-n-\sigma q}\,d(y)^{-aq}dy\Big)^{p/q}dx\Big)^{1/p}.
\end{equation}

\subsection{两个闭包与 Hardy 加强空间}

\begin{definition}[$C_c^\infty$ 与 $C^\infty$ 的闭包]\label{def:closures}
在范数 $\|\cdot\|_{\Fsp}$ 下定义
\begin{align}
  F^{s,1}_{p,q}(\Om,d^{-aq})&:=\overline{\,C_c^\infty(\Om)\cap \Fsp\,}^{\;\|\cdot\|_{\Fsp}},\label{eq:F1}\\
  F^{s,2}_{p,q}(\Om,d^{-aq})&:=\overline{\,C^\infty(\Om)\cap \Fsp\,}^{\;\|\cdot\|_{\Fsp}}.\label{eq:F2}
\end{align}
称 $F^{s,1}$ 为\defterm{$C_c^\infty$ 闭包}、$F^{s,2}$ 为\defterm{$C^\infty$ 闭包}。
\end{definition}

\begin{definition}[Hardy 加强范数]\label{def:hardy}
定义两个加强范数
\begin{align}
  \|f\|_{\tiF}&:=\|f\|_{\Fsp}
    +\big\|f\,d^{-(s+a)}\big\|_{L^p(\Om)}
    +\big\|f\,d^{-(s+a+n/q)}\big\|_{L^q(\Om)},\label{eq:tilde}\\
  \|f\|_{\hoF}&:=\|f\|_{\Fsp}
    +\big\|f\,d^{-(s+a)}\big\|_{L^p(\Om)}.\label{eq:hat}
\end{align}
相应空间记作 $\tiF$ 与 $\hoF$，其闭包
$\wt F^{s,1},\wt F^{s,2},\wh F^{s,1},\wh F^{s,2}$
按 \eqref{eq:F1}--\eqref{eq:F2} 以对应范数定义。
\end{definition}

\begin{notation}
\eqref{eq:tilde}--\eqref{eq:hat} 中出现的两个权指数
\[
  s+a \qquad\text{与}\qquad s+a+\tfrac{n}{q}
\]
在本文中反复出现；后者因带额外因子 $d^{-n/q}$ 而严格更强，其作用见第
\ref{sec:part6} 节。
\end{notation}

\subsection{边界几何量}

本文多次用到边界的 Assouad 余维数与 Minkowski 维数。

\begin{definition}[Assouad 余维数]\label{def:codim}
设 $E\subset\R^d$ 非空，$B(x,R)$ 为欧氏球。称
\[
  \codimAl(E):=\sup\Big\{q\ge0:\ \exists C=C(q)\ge1,\ 
  |E_r\cap B(x,R)|\le C|B(x,R)|\Big(\tfrac{r}{R}\Big)^{q}\ 
  \Big\},
\]
\[
  \codimAu(E):=\inf\Big\{s\ge0:\ \exists c=c(s)>0,\ 
  |E_r\cap B(x,R)|\ge c|B(x,R)|\Big(\tfrac{r}{R}\Big)^{s}\ 
  \Big\},
\]
两式中上确界与下确界均取遍使相应不等式对\defterm{一切} $x\in E$ 与 $0<r<R$ 成立的 $q$ 或 $s$；
其中 $E_r:=\{t:\ \dist(t,E)\le r\}$ 为管状邻域。恒有 $\codimAl(E)\le\codimAu(E)$。
\end{definition}

\begin{definition}[Lipschitz 域与 uniform 域]\label{def:domains}
称 $\Om$ 为 \defterm{Lipschitz 域}，若对每个 $x_0\in\partial\Om$ 存在 $r>0$、
$\R^n$ 的一个旋转以及 Lipschitz 函数 $\varphi:\R^{n-1}\to\R$，使
\[
  \Om\cap B(x_0,r)=\{x=(x',x_n)\in B(x_0,r):\ x_n>\varphi(x')\},\quad
  \partial\Om\cap B(x_0,r)=\{x\in B(x_0,r):\ x_n=\varphi(x')\}.
\]
称 $\Om$ 为 \defterm{uniform 域}，若存在 $c_1,c_2>0$，使任意 $x,y\in\Om$ 可用
$\Om$ 内的可求长曲线 $\Gamma$ 连接，且
$L(\Gamma)\le c_1|x-y|$、$\min\{|x-z|,|y-z|\}\le c_2 d(z)$ 对一切 $z\in\Gamma$ 成立。
\end{definition}

\subsection{本文的结论结构}

全部 32 条事实按论证作用分为六组，分别对应第 \ref{sec:part1}--\ref{sec:part6} 节。其中 28 条在正文中以独立的定理、引理或命题形式陈述（即下表所计）；另 4 条仅作为前置或交叉引用出现，其标识同样列于附录 A：

\begin{center}
\small
\begin{tabular}{@{}cll@{}}
\toprule
节 & 主题 & 条数 \\
\midrule
\S\ref{sec:part1} & 基础引理：二分、刚性引理、多项式引理与边界几何 & 6 \\
\S\ref{sec:part2} & 情形 (II)：不可积权下的退化 & 5 \\
\S\ref{sec:part3} & 情形 (I)：可积权下的稠密性 & 4 \\
\S\ref{sec:part4} & 主强化定理与反例 & 3 \\
\S\ref{sec:part5} & 临界条件：边界余维数，而非权可积性 & 4 \\
\S\ref{sec:part6} & Hardy 项的非冗余性、延拓路线的终结与技术工具 & 6 \\
\midrule
 & 小计（独立陈述） & 28 \\
\bottomrule
\end{tabular}
\end{center}

\noindent
第 \ref{sec:concl} 节汇总主结论、给出事实依赖图并列出尚未解决的问题；
第 \ref{sec:limits} 节说明验证状态与局限。

%=====================================================================
\section{基础引理：二分、刚性引理、多项式引理与边界几何}\label{sec:part1}
%=====================================================================

本节汇总论证所依赖的六条基础引理。前三条围绕边界距离权重 $d^{-aq}$ 的二分类现象展开：该权重或者整体可积，或者在每个边界层上都已经发散，而后者正是迫使函数为常数的刚性机制；第四条与第五条分别提供"高阶弱导数恒定蕴含多项式"的代数工具与加权积分散射性的单调性；第六条则给出 Lipschitz 区域边界层的线性下界，从而保证所有 $d^{-c}$（$c\ge 1$）都不可积。这些结论彼此独立而又相互配合，是后续加权 Triebel--Lizorkin 空间估计的出发点。

%==================================================================
\subsection{边界距离权重的二分性}

\begin{lemma}[边界距离权重的二分性]
设 $n\ge 1$，$\Om$ 为 $\R^n$ 中的非空有界域（非空、连通、有界、开），对 $x\in\Om$ 令 $d(x)=\dist(x,\partial\Om)$，设 $a>0$、$q\in[1,\infty)$，并设 $\nuA$ 为 $\Om$ 上由
\[
  \nuA(E)=\int_E d(y)^{-aq}\,dy .
\]
定义的 Borel 测度。对 $r>0$ 令 $\Omr=\{x\in\Om : d(x)<r\}$。则下述两种情形恰有一种成立：
\begin{enumerate}[label=(\Roman*)]
\item $\nuA(\Om)<\infty$；
\item $\nuA(\Om)=\infty$，而在此情形下自动对每个 $r>0$ 都有 $\nuA(\Omr)=\infty$。
\end{enumerate}
\end{lemma}

\begin{proof}
由于 $d$ 在 $\Om$ 上连续且为正，且 $aq\ge 0$，函数 $d^{-aq}$ 在 $\Om$ 上 Borel 可测且非负，故对每个 Borel 集 $E\subseteq\Om$，$\nuA(E)$ 有定义（可能为 $+\infty$）；$\Omr$ 是开集，从而是 Borel 集。$\nuA(\Om)<\infty$ 与 $\nuA(\Om)=\infty$ 恰有一个成立，因此只需证明第二种情形蕴含对每个 $r>0$ 都有 $\nuA(\Omr)=\infty$。

设 $\nuA(\Om)=\infty$ 并固定 $r>0$。则 $\Om$ 是 $\Omr$ 与 $\Om\setminus\Omr$ 的不交并，故
\[
  \nuA(\Om)=\nuA(\Omr)+\int_{\Om\setminus\Omr} d(y)^{-aq}\,dy .
\]
对 $y\in\Om\setminus\Omr$ 有 $d(y)\ge r$，故 $d(y)^{-aq}\le r^{-aq}$；因此
\[
  \int_{\Om\setminus\Omr} d(y)^{-aq}\,dy\le r^{-aq}\,|\Om|<\infty ,
\]
因为 $\Om$ 有界。由于 $\nuA(\Om)=\infty$ 而第二项有限，故 $\nuA(\Omr)=\infty$。又 $r>0$ 任意，结论成立。
\end{proof}

\begin{remark}[直觉]
密度为 $d^{-aq}$ 的测度要么总质量有限；要么一旦总质量无限，就已经在每个边界层上无限。原因在于边界层之外密度被常数 $r^{-aq}$ 控制，而 $\Om$ 的体积有限。
\end{remark}

%==================================================================
\subsection{\texorpdfstring{$\sigma$}{sigma}-引理：权重不可积时的刚性}

\begin{lemma}[$\sigma$-引理：权重不可积时最高阶导数为常数]
设 $n\ge 1$，$\Om$ 为 $\R^n$ 中的非空有界域，$d(x)=\dist(x,\partial\Om)$，$a>0$，$q\in[1,\infty)$，$\sigma\in(0,1)$，$p\in[1,\infty)$，并在 $\Om$ 上令 $\nuA(dy)=d(y)^{-aq}\,dy$。对 $r>0$ 令 $\Omr=\{y\in\Om : d(y)<r\}$。假设 $\nuA(\Om)=\infty$。设 $g:\Om\to\R$ 为 Lebesgue 可测，并假设
\[
  \int_\Om \Big(\int_\Om |g(x)-g(y)|^q\,|x-y|^{-n-\sigma q}\,\nuA(dy)\Big)^{p/q} dx<\infty ,
\]
其中内层积分遍历全部 $y\in\Om$。则 $g$ 在 $\Om$ 上几乎处处为常数。
\end{lemma}

\begin{proof}
设 $\Om$ 为 $\R^n$ 中的非空有界域，$d(x)=\dist(x,\partial\Om)$，$a>0$，$q\in[1,\infty)$，$\sigma\in(0,1)$，$p\in[1,\infty)$，且 $\nuA(dy)=d(y)^{-aq}\,dy$。对 $r>0$ 令 $\Omr:=\{y\in\Om : d(y)<r\}$。假设 $\nuA(\Om)=\infty$。设 $g:\Om\to\R$ 为 Lebesgue 可测，并假设
\[
  W:=\int_\Om \big(J(x)\big)^{p/q}\,dx<\infty ,
\]
其中
\[
  J(x):=\int_\Om |g(x)-g(y)|^q\,|x-y|^{-n-\sigma q}\,\nuA(dy)\in[0,\infty]
\]
为 Borel 可测（被积函数联合可测）。由已验证的附录 A（边界距离权重的二分性），将其应用于 $\Om$、$a$、$q$，得到对每个 $r>0$ 都有 $\nuA(\Omr)=\infty$。

\emph{第一步（好点）.} 由于 $J\ge 0$ 且 $\int_\Om J^{p/q}<\infty$，集合
\[
  G:=\{x\in\Om : J(x)<\infty\}
\]
在 $\Om$ 中具有满 Lebesgue 测度。固定 $x\in G$ 并令 $r_x:=d(x)/2>0$。对每个 $y\in\Om$（特别对 $y\in\Om_{r_x}$）都有 $|x-y|\le\diam(\Om)<\infty$，因为 $\Om$ 有界。故 $|x-y|^{-n-\sigma q}\ge(\diam\Om)^{-n-\sigma q}$，从而
\[
  J(x)\ge(\diam\Om)^{-n-\sigma q}\int_{\Om_{r_x}} |g(x)-g(y)|^q\,\nuA(dy) .
\]
于是
\[
  \int_{\Om_{r_x}} |g(x)-g(y)|^q\,\nuA(dy)\le(\diam\Om)^{n+\sigma q}J(x)<\infty .
\]

\emph{第二步（集合 $B_\varepsilon(x)$）}. 对 $\varepsilon>0$ 令
\[
  A_\varepsilon(x):=\{y\in\Om_{r_x} : |g(x)-g(y)|>\varepsilon\},
  \qquad
  B_\varepsilon(x):=\{y\in\Om_{r_x} : |g(x)-g(y)|\le\varepsilon\} .
\]
两者均可测，且 $\Om_{r_x}$ 是它们的不交并。由第一步，
\[
  \varepsilon^q\,\nuA\big(A_\varepsilon(x)\big)
  \le\int_{A_\varepsilon(x)} |g(x)-g(y)|^q\,\nuA(dy)
  \le(\diam\Om)^{n+\sigma q}J(x)<\infty .
\]
由附录 A，$\nuA(\Om_{r_x})=\infty$，故
\[
  \nuA\big(B_\varepsilon(x)\big)=\nuA(\Om_{r_x})-\nuA\big(A_\varepsilon(x)\big)=\infty .
\]

\emph{第三步（配对两个好点）}. 设 $x_1,x_2\in G$，并令 $r:=\min(r_{x_1},r_{x_2})>0$。对 $i=1,2$ 有 $B_\varepsilon(x_i)\subseteq\Om_{r_{x_i}}$，且由 $r\le r_{x_i}$ 得 $\Omr\subseteq\Om_{r_{x_i}}$。因此
\[
  \Omr\setminus\big(A_\varepsilon(x_1)\cup A_\varepsilon(x_2)\big)\subseteq B_\varepsilon(x_1)\cap B_\varepsilon(x_2),
\]
并利用可加性以及第二步所得的 $\nuA(A_\varepsilon(x_1))<\infty$ 与 $\nuA(A_\varepsilon(x_2))<\infty$，
\[
  \nuA\big(B_\varepsilon(x_1)\cap B_\varepsilon(x_2)\big)
  \ge\nuA(\Omr)-\nuA\big(A_\varepsilon(x_1)\big)-\nuA\big(A_\varepsilon(x_2)\big)=\infty ,
\]
这里再次用了应用于 $r$ 的附录 A。特别地，$B_\varepsilon(x_1)\cap B_\varepsilon(x_2)$ 非空；在其中取 $y$。则
\[
  |g(x_1)-g(x_2)|\le|g(x_1)-g(y)|+|g(y)-g(x_2)|\le 2\varepsilon .
\]
由于 $\varepsilon>0$ 任意，得 $g(x_1)=g(x_2)$。

\emph{第四步（结论）}. 第三步表明 $g$ 在满测度集 $G$ 上为常数。由于 $\Om$ 非空且开，$|\Om|>0$，故 $G$ 非空；取 $x_0\in G$ 并令 $c:=g(x_0)$。则 $g=c$ 在 $\Om$ 上几乎处处成立。命题得证。
\end{proof}

\begin{remark}[直觉]
若权重不可积，则每个边界层 $\Om_{d(x)/2}$ 都携带无限权重；对任意两点 $x_1,x_2$，使 $g$ 保持在 $g(x_i)$ 的 $\varepsilon$ 邻域内的两个集合在一个公共内层 $\Omr$ 中都具有无限权重，因此必然相交，从而迫使 $|g(x_1)-g(x_2)|\le 2\varepsilon$。
\end{remark}

%==================================================================
\subsection{非可积刚性（\texorpdfstring{$\sigma$}{sigma}-引理，第二版本）}

\begin{lemma}[不可积刚性]
设 $n\ge 1$，$\Om$ 为 $\R^n$ 中的非空有界域，对 $x\in\Om$ 令 $d(x)=\dist(x,\partial\Om)$，设 $a>0$，$q\in[1,\infty)$，$\sigma\in(0,1)$，$p\in[1,\infty)$。在 $\Om$ 上定义 Borel 测度 $\nuA$：$\nuA(E)=\int_E d(y)^{-aq}\,dy$，并假设 $\nuA(\Om)=\infty$。设 $g:\Om\to\R$ 为 Lebesgue 可测，且假设
\[
  \int_\Om \Big(\int_\Om |g(x)-g(y)|^q\,|x-y|^{-n-\sigma q}\,d(y)^{-aq}\,dy\Big)^{p/q} dx<\infty .
\]
则 $g$ 在 $\Om$ 上 Lebesgue 几乎处处为常数。
\end{lemma}

\begin{proof}
由于 $d$ 在开集 $\Om$ 上连续且严格为正，且 $aq\ge 0$，函数 $y\mapsto d(y)^{-aq}$ 是 $\Om$ 上的非负 Borel 函数；故 $\nuA$ 是 $\Om$ 上取值于 $[0,+\infty]$ 的、定义良好的 Borel 测度，并且对每个 $r>0$，$\Omr=\{y\in\Om : d(y)<r\}$ 是开集，从而是 Borel 集。令 $D:=\diam(\Om)=\sup\{|u-v| : u,v\in\Om\}$；由于 $\Om$ 有界且非空，$0<D<\infty$。

对 $x\in\Om$ 令
\[
  J(x):=\int_\Om |g(x)-g(y)|^q\,|x-y|^{-n-\sigma q}\,d(y)^{-aq}\,dy\in[0,+\infty] .
\]
被积函数 $(x,y)\mapsto|g(x)-g(y)|^q|x-y|^{-n-\sigma q}d(y)^{-aq}$ 在对角线之外（一个 Lebesgue 零集）Borel 可测且非负，故由 Tonelli 定理，$J$ 是 Borel 可测的。假设给出 $\int_\Om J(x)^{p/q}\,dx<\infty$；由于 $p/q>0$，这迫使对几乎每个 $x\in\Om$（Lebesgue 意义下）都有 $J(x)<\infty$。又对每个 $x\in\Om$ 有 $d(x)>0$：$d$ 在开集 $\Om$ 上连续，而 $d(x)=0$ 会使 $x$ 落在 $\partial\Om$ 中，它与 $\Om$ 不交。令
\[
  E:=\{x\in\Om : J(x)<\infty\} ;
\]
则 $E$ 是 Lebesgue 可测的且 $|\Om\setminus E|=0$。

\emph{断言.} 若 $x$ 与 $z$ 属于 $E$，则 $g(x)=g(z)$。

固定 $x\in E$ 并令 $r_x:=d(x)/2>0$。对 $y\in\Om_{r_x}$ 有 $d(y)<r_x=d(x)/2<d(x)$，故 $y$ 不是 $x$ 的最近边界点，而距离函数的反向三角不等式给出
\[
  |x-y|\ge d(x)-d(y)>d(x)/2=r_x>0 .
\]
另一方面，由 $y\in\Om$ 与 $x\in\Om$ 得 $|x-y|\le D$。因此对每个 $y\in\Om_{r_x}$，
\[
  |x-y|^{-n-\sigma q}\ge D^{-n-\sigma q},
\]
并且由于 $\Om_{r_x}$ 含于 $\Om$，
\begin{equation}
  \int_{\Om_{r_x}} |g(x)-g(y)|^q\,d(y)^{-aq}\,dy\le D^{n+\sigma q}J(x)<\infty . \tag{1}
\end{equation}
对 $\varepsilon>0$ 令 $S_x:=\{y\in\Om_{r_x} : |g(x)-g(y)|>\varepsilon\}$。则利用 (1)，
\[
  \varepsilon^q\,\nuA(S_x)
  =\int_{S_x}\varepsilon^q\,d(y)^{-aq}\,dy
  \le\int_{S_x}|g(x)-g(y)|^q\,d(y)^{-aq}\,dy
  \le D^{n+\sigma q}J(x)<\infty ,
\]
故 $\nuA(S_x)<\infty$。由于 $\nuA(\Om)=\infty$ 且 $r_x>0$，二分性引理（附录 A）给出 $\nuA(\Om_{r_x})=\infty$；因此 $\nuA(\Om_{r_x}\setminus S_x)=\infty$。

现令 $z$ 为 $E$ 中的第二个点，并令 $r_z:=d(z)/2>0$、$r:=\min(r_x,r_z)>0$。定义 $S_z:=\{y\in\Om_{r_z} : |g(z)-g(y)|>\varepsilon\}$，由同样的论证有 $\nuA(S_z)<\infty$。由于 $\Omr$ 同时含于 $\Om_{r_x}$ 与 $\Om_{r_z}$，有
\[
  \Omr\setminus(S_x\cup S_z)=\Omr\setminus\big((S_x\cap\Omr)\cup(S_z\cap\Omr)\big)
\]
从而
\[
\begin{aligned}
  \nuA\big(\Omr\setminus(S_x\cup S_z)\big)
  &\ge\nuA(\Omr)-\nuA(S_x\cap\Omr)-\nuA(S_z\cap\Omr)\\
  &\ge\nuA(\Omr)-\nuA(S_x)-\nuA(S_z)=\infty ,
\end{aligned}
\]
这是因为 $\nuA(\Omr)=\infty$（二分性引理，$r>0$），而 $\nuA(S_x)$、$\nuA(S_z)$ 有限。特别地，$\Omr\setminus(S_x\cup S_z)$ 非空；在其中取 $y$。则 $y\in\Omr$ 且 $y\notin S_x$，故 $|g(x)-g(y)|\le\varepsilon$；同理 $|g(z)-g(y)|\le\varepsilon$。因此
\[
  |g(x)-g(z)|\le|g(x)-g(y)|+|g(y)-g(z)|\le2\varepsilon .
\]
由于 $\varepsilon>0$ 任意，得 $g(x)=g(z)$，断言得证。

于是 $g$ 在 $E$ 上为常数，又因 $|\Om\setminus E|=0$，$g$ 在 $\Om$ 上 Lebesgue 几乎处处为常数。
\end{proof}

\begin{remark}[直觉]
若测度 $d^{-aq}dy$ 在每个边界层上都有无限质量，则内层积分迫使 $|g(x)-g(y)|$ 在一个 $\nuA$-测度无限的 $y$-集上很小；两个这样的集合（对应于两个好点 $x$ 与 $z$）必然相交，从而 $g(x)=g(z)$。$\Om$ 的有界性用初等估计 $|x-y|\le\diam(\Om)$ 替代了错误的边界层估计。
\end{remark}

%==================================================================
\subsection{多项式引理}

\begin{lemma}[多项式引理：最高阶弱导数为常数蕴含 $f$ 为多项式]
设 $n\ge 1$，$\Om$ 为 $\R^n$ 中的非空有界域，设 $k\in\N_0=\{0,1,2,\dots\}$，$p\in[1,\infty)$，并设 $f\in W^{k,p}(\Om)$，即 $L^p(\Om)$ 中直到 $k$ 阶弱导数都属于 $L^p(\Om)$ 的函数构成的 Sobolev 空间。假设对每个满足 $|\beta|=k$ 的多重指标 $\beta\in\N_0^n$，都存在常数 $c_\beta\in\R$，使得弱导数 $D^\beta f$ 在 $\Om$ 上几乎处处等于 $c_\beta$。则存在 $\R^n$ 上次数至多 $k$ 的多项式 $P$，使得 $f=P$ 在 $\Om$ 上几乎处处成立。
\end{lemma}

\begin{proof}
设 $\Om$ 为 $\R^n$ 中的非空有界域，$k\in\N_0$，$p\in[1,\infty)$，$f\in W^{k,p}(\Om)$，并假设对每个满足 $|\beta|=k$ 的多重指标 $\beta$ 都存在 $c_\beta\in\R$，使得弱导数 $D^\beta f$ 在 $\Om$ 上几乎处处等于常数 $c_\beta$。

\emph{情形 $k=0$.} 此时满足 $|\beta|=0$ 的多重指标只有 $\beta=0$，且 $D^0f=f$，故 $f$ 几乎处处等于常数 $c_0$；这是次数 $0\le k$ 的多项式。

\emph{情形 $k\ge 1$.} 在 $\R^n$ 上定义多项式 $Q$：
\[
  Q(x):=\sum_{|\beta|=k} c_\beta\,\frac{x^\beta}{\beta!} .
\]
其 $k$ 阶弱导数就是常数 $c_\beta$，故 $g:=f-Q$ 属于 $W^{k,p}(\Om)$，且对每个满足 $|\beta|=k$ 的 $\beta$，$D^\beta g=0$ 在 $\Om$ 上几乎处处成立。

\emph{断言.} 设 $B=B(a,\rho)$ 是闭包含于 $\Om$ 的开球。则存在 $\R^n$ 上次数至多 $k-1$ 的多项式 $p_B$，使得 $g=p_B$ 在 $B$ 上几乎处处成立。

\emph{断言的证明.} 设 $\rho_\varepsilon$ 为标准磨光核（非负、光滑、支集在闭单位球内、积分为 $1$），并令 $\rho_\varepsilon(z)=\varepsilon^{-n}\rho(z/\varepsilon)$。取 $\varepsilon$ 满足 $0<\varepsilon<\dist(B,\partial\Om)$。对 $x\in B$ 及每个满足 $|\beta|=k$ 的 $\beta$，关于磨光函数导数的标准恒等式给出
\[
  D^\beta(g*\rho_\varepsilon)(x)=(D^\beta g*\rho_\varepsilon)(x)=0 ,
\]
因为 $D^\beta g=0$ 几乎处处且 $\rho_\varepsilon$ 是磨光核。因此光滑函数 $g*\rho_\varepsilon$ 的所有 $k$ 阶偏导数在 $B$ 上为零。由于 $B$ 连通且所有 $k$ 阶偏导数在 $B$ 上为零，$g*\rho_\varepsilon$ 是某个次数至多 $k-1$ 的多项式 $p_\varepsilon$ 在 $B$ 上的限制（这是"连通开集上 $k$ 阶导数全为零"对多项式的标准刻画：对每个坐标方向作限制并积分 $k$ 次表明 $g*\rho_\varepsilon$ 是每个变量中次数至多 $k-1$ 的多项式，从而是多项式）。

现令 $\varepsilon_j:=1/j$，其中 $j$ 充分大使得 $\varepsilon_j<\dist(B,\partial\Om)$。当 $j\to\infty$ 时，函数 $g*\rho_{\varepsilon_j}$ 在 $L^p(B)$ 中收敛于 $g$（内部磨光的标准性质）。每个 $g*\rho_{\varepsilon_j}$ 在 $B$ 上等于 $p_{\varepsilon_j}$，而 $p_{\varepsilon_j}$ 属于次数至多 $k-1$ 的多项式构成的有限维空间 $P_{k-1}$。多项式到开集 $B$ 的限制是单射，且在有限维空间 $P_{k-1}$ 上所有范数等价；于是 $p_{\varepsilon_j}$ 的系数向量构成 Cauchy 列，故存在 $p_B\in P_{k-1}$ 使得 $p_{\varepsilon_j}\to p_B$ 在 $L^p(B)$ 中成立（并且沿某个子列还逐点收敛）。由于 $p_{\varepsilon_j}=g*\rho_{\varepsilon_j}\to g$ 在 $L^p(B)$ 中，得 $g=p_B$ 在 $B$ 上几乎处处成立。断言得证。

现用开球 $B_j=B(x_j,r_j)$ 覆盖 $\Om$，并要求 $B_j$ 的闭包含于 $\Om$（例如取所有中心与半径均为有理数、闭包含于 $\Om$ 的球；由于 $\Om$ 开，$\Om$ 的每一点都落在这样的球中）。对每个 $j$，设 $p_j$ 为断言所提供的、次数至多 $k-1$ 且满足 $g=p_j$ 在 $B_j$ 上几乎处处成立的多项式。若 $B_i$ 与 $B_j$ 相交，则其交为非空开集，而 $g=p_i=p_j$ 在其上几乎处处成立，故作为多项式有 $p_i=p_j$（在非空开集上相等的两个多项式相等）。由于 $\Om$ 连通，对每个 $i,j$ 都存在覆盖中的有限球链 $B_{j_0}=B_i,\dots,B_{j_m}=B_j$，使相邻的球相交：事实上，由 $B_i$ 经这样的链可达的全部球之并 $U$ 是开集，且它在 $\Om$ 中的补集也是开集（若 $x\in\Om\setminus U$ 落在 $B_k$ 中且 $B_k$ 与 $U$ 相交，则 $B_k$ 将是可达的），故由连通性 $U=\Om$。沿这样的链所有 $p$ 都相同，于是 $p_i=p_j$。因此所有 $p_j$ 都等于一个固定的次数至多 $k-1$ 的多项式 $P$，且 $g=P$ 在 $\Om$ 上几乎处处成立，即 $f=Q+P$ 在 $\Om$ 上几乎处处成立。多项式 $Q+P$ 的次数至多 $k$。命题得证。
\end{proof}

\begin{remark}[直觉]
若全部 $k$ 阶弱导数都是常数，先减去以这些常数为 $k$ 阶导数的多项式，则所有 $k$ 阶弱导数都为零。磨光把一个 $k$ 阶导数全为零的函数在每个球上变成次数 $<k$ 的多项式，而多项式空间的有限维性与 $\Om$ 的连通性使这些局部多项式拼成一个整体多项式。
\end{remark}

%==================================================================
\subsection{距离权重积分发散性的单调性}

\begin{lemma}[距离权重积分发散性的单调性]
设 $n\ge 1$，$\Om$ 为 $\R^n$ 中的非空有界域，对 $x\in\Om$ 令 $d(x)=\dist(x,\partial\Om)$。设 $0\le c_1\le c_2$ 为实数。若 $\int_\Om d(y)^{-c_1}\,dy=\infty$，则 $\int_\Om d(y)^{-c_2}\,dy=\infty$。特别地，若 $a>0$、$q\in[1,\infty)$ 且 $\int_\Om d(y)^{-aq}\,dy=\infty$，则对每个实数 $c\ge aq$ 都有 $\int_\Om d(y)^{-c}\,dy=\infty$。
\end{lemma}

\begin{proof}
设 $\Om$ 为 $\R^n$ 中的非空有界域，$d(x)=\dist(x,\partial\Om)$。令 $R:=\diam(\Om)$，则 $0<R<\infty$ 且对每个 $y\in\Om$ 有 $0<d(y)\le R$。设 $0\le c_1\le c_2$。对每个 $y\in\Om$ 有
\[
  d(y)^{-c_2}=d(y)^{-c_1}\,d(y)^{-(c_2-c_1)}\ge d(y)^{-c_1}\,R^{-(c_2-c_1)} ,
\]
因为 $d(y)^{-(c_2-c_1)}\ge R^{-(c_2-c_1)}$。在 $\Om$ 上积分（两端都是非负可测函数，故积分不等式在 $[0,\infty]$ 中成立）得
\[
  \int_\Om d(y)^{-c_2}\,dy\ge R^{-(c_2-c_1)}\int_\Om d(y)^{-c_1}\,dy .
\]
因此，若 $\int_\Om d(y)^{-c_1}\,dy=\infty$，则 $\int_\Om d(y)^{-c_2}\,dy=\infty$。特别地，取 $c_1=aq$ 与 $c_2=c$（$c\ge aq$），$\int_\Om d^{-aq}$ 的发散性蕴含 $\int_\Om d^{-c}$ 的发散性。
\end{proof}

\begin{remark}[直觉]
在有界区域上距离函数不超过直径，因此 $d$ 的较高次幂至少是较低次幂的常数倍；于是较低次幂积分的发散性迫使所有更高次幂的积分发散。
\end{remark}

%==================================================================
\subsection{Lipschitz 区域边界层的线性下界}

\begin{theorem}[边界层的线性下界，以及所有幂 $d^{-c}$（$c\ge 1$）的不可积性]
设 $n\ge 1$，$\Om$ 为 $\R^n$ 中的有界 Lipschitz 域（即非空有界开连通集，使得对每个 $x_0\in\partial\Om$ 都存在 $r>0$、$\R^n$ 的一个旋转以及 Lipschitz 函数 $\phi:\R^{n-1}\to\R$，满足
\[
  \Om\cap B(x_0,r)=\{x=(x',x_n)\in B(x_0,r) : x_n>\phi(x')\}
\]
以及
\[
  \partial\Om\cap B(x_0,r)=\{x=(x',x_n)\in B(x_0,r) : x_n=\phi(x')\}\,\big).
\]
令 $d(x)=\dist(x,\partial\Om)$，$\Omt=\{x\in\Om : d(x)<t\}$。则存在常数 $\kappa>0$ 与 $t_0>0$，使得对每个 $t\in(0,t_0)$ 都有 $|\Omt|\ge\kappa t$。因此，对每个实数 $c\ge1$，
\[
  \int_\Om d(x)^{-c}\,dx=\infty .
\]
\end{theorem}

\begin{proof}
设 $\Om$ 为 $\R^n$ 中的有界 Lipschitz 域，$d(x)=\dist(x,\partial\Om)$，$\Omt=\{x\in\Om : d(x)<t\}$。

\emph{第一步（局部带）}. 由 Lipschitz 域的定义，对每个 $x_0\in\partial\Om$ 都存在 $r(x_0)>0$、一个旋转以及 Lipschitz 函数 $\phi$，使得
\[
  \Om\cap B(x_0,r(x_0))=\{(x',x_n)\in B(x_0,r(x_0)) : x_n>\phi(x')\}
\]
以及
\[
  \partial\Om\cap B(x_0,r(x_0))=\{(x',x_n)\in B(x_0,r(x_0)) : x_n=\phi(x')\} .
\]
由于 $\partial\Om$ 紧，这些球 $B(x_0,r(x_0))$ 中的有限多个（设为 $x_0=y_1,\dots,y_m$ 对应的那些）覆盖 $\partial\Om$。对每个 $j$，取 $\R^{n-1}$ 中的一个开球 $B'_j$（在第 $j$ 个坐标卡的旋转坐标下），使其闭包含于 $\partial\Om\cap B(y_j,r(y_j))$ 到 $\R^{n-1}$ 的投影之内，且它与该投影的补集的距离为正。由于图映射 $x'\mapsto(x',\phi_j(x'))$ 连续而该投影是相对开集，这样的球存在。由 $B'_j$ 的闭包的紧性及在坐标卡中的严格包含，存在 $t_0>0$，使得对每个 $j$、每个 $t\in(0,t_0)$ 与每个 $x'\in B'_j$，点 $(x',\phi_j(x')+t)$ 都落在 $B(y_j,r(y_j))$ 中。

\emph{第二步（带落在 $\Omt$ 内）}. 对 $j$ 与 $t\in(0,t_0)$ 令
\[
  W_j(t):=\{(x',x_n) : x'\in B'_j,\ \phi_j(x')<x_n<\phi_j(x')+t\}
\]
（在第 $j$ 个旋转坐标下）。则 $W_j(t)$ 含于 $\Om\cap B(y_j,r(y_j))$，从而含于 $\Om$。对 $x=(x',x_n)\in W_j(t)$，点 $(x',\phi_j(x'))$ 属于 $\partial\Om$，故 $d(x)\le|x_n-\phi_j(x')|<t$。因此 $W_j(t)$ 含于 $\Omt$ 中。

\emph{第三步（带的测度）}. 对每个 $j$ 与 $t\in(0,t_0)$，由 Fubini 定理得
\[
  |W_j(t)|=\int_{B'_j}\big((\phi_j(x')+t)-\phi_j(x')\big)\,dx'=t\,|B'_j| .
\]
由于各 $B'_j$ 为非空开球，对每个 $j$ 有 $|B'_j|>0$。故对每个 $j$ 有 $|\Omt|\ge|W_j(t)|$，特别地
\[
  |\Omt|\ge t\max_j|B'_j|=:\kappa t
\]
对所有 $t\in(0,t_0)$ 成立，其中 $\kappa>0$。

\emph{第四步（发散性）}. 对 $c\ge1$，分层表示公式（作代换 $s=t^{-c}$）给出
\[
  \int_\Om d(x)^{-c}\,dx
  =\int_0^\infty |\{x\in\Om : d(x)^{-c}>s\}|\,ds
  =\int_0^\infty c\,t^{-c-1}|\Omt|\,dt ,
\]
因为 $\{d^{-c}>t^{-c}\}=\{d<t\}$。利用第三步在 $(0,t_0)$ 上的估计，
\[
  \int_\Om d(x)^{-c}\,dx
  \ge\int_0^{t_0} c\,t^{-c-1}\,\kappa t\,dt
  =c\kappa\int_0^{t_0} t^{-c}\,dt=\infty ,
\]
因为对每个 $c\ge1$ 都有 $\int_0^{t_0} t^{-c}\,dt=\infty$（当 $c=1$ 时它为 $\log(1/t_0)=\infty$；当 $c>1$ 时原函数 $t^{1-c}/(1-c)$ 在 $t\to0$ 时趋于无穷）。因此对每个 $c\ge1$ 都有 $\int_\Om d^{-c}\,dx=\infty$。
\end{proof}

\begin{remark}[直觉]
Lipschitz 边界局部是某个函数的图像，因此在图像的每一小块上方都存在厚度为 $t$ 的带完全落在 $\Om$ 内，且带上每一点到边界的距离都 $<t$；带的测度等于底面面积乘以 $t$，而有限多个这样的带覆盖了整个边界。
\end{remark}


%=====================================================================
\section{情形 (II)：不可积权下的退化}\label{sec:part2}
%=====================================================================
% ===================================================================
%  第 2 组：情形 (II) —— 不可积权下的退化
% ===================================================================

本节处理边界距离权重不可积的情形，即 $\nuA(\Om)=\int_\Om d(y)^{-aq}\,dy=\infty$。此时加权内在 Triebel--Lizorkin 空间 $\Fsp$ 退化为次数不超过 $k=\lfloor s\rfloor$ 的多项式空间 $P_k(\Om)$：于是 $C^\infty(\Om)$ 自动稠密，而紧支光滑函数被完全排除。随后的三条事实表明这一退化是权重设计的产物而非本质现象——CLWZ 定义 1.3 的修改版同样只剩常数，而整数阶修改空间与"对一切导数都取加权分数阶半范数"的修正空间都保留了非多项式元素。

%==================================================================
\subsection{非可积权下分数阶空间的退化（第一版本）}

\begin{theorem}[不可积情形的退化]
设 $\Om$ 为 $\R^n$ 中的非空有界域，$a>0$ 且 $q\in[1,\infty)$，$d(y)=\dist(y,\partial\Om)$，$\nuA(E)=\int_E d(y)^{-aq}\,dy$。假设 $\nuA(\Om)=\infty$。设 $s\in(0,\infty)\setminus\N$，写 $s=k+\sigma$，其中 $k=\lfloor s\rfloor\in\N_0$、$\sigma\in(0,1)$，并设 $p\in[1,\infty)$。则由 CLWZ 的范数 (1.2) 定义的加权内蕴 Triebel--Lizorkin 空间 $\Fsp$ 与 $P_k(\Om)$ 重合，后者是总次数不超过 $k$ 的多项式在 $\Om$ 上限制所成的集合。于是：
\begin{enumerate}[label=(\alph*)]
\item $C^\infty(\Om)\cap\Fsp=\Fsp$;
\item $C_c^\infty(\Om)\cap\Fsp=\{0\}$;
\item $C^\infty(\Om)\cap F^s$ 在 $F^s$ 中的闭包 $F^{s,2}_{p,q}$ 等于 $F^s$，而 $C_c^\infty(\Om)\cap F^s$ 在 $F^s$ 中的闭包 $F^{s,1}_{p,q}$ 等于 $\{0\}$。
\end{enumerate}
\end{theorem}

\begin{proof}
以下始终固定 $\Om$ 上的可测函数 $f$ 与满足 $|\beta|=k$ 的多重指标 $\beta$，并记
\[
  J_{\beta,f}(x)=\int_\Om |D^\beta f(x)-D^\beta f(y)|^q\,|x-y|^{-n-\sigma q}\,d(y)^{-aq}\,dy .
\]
由 CLWZ (1.2)，$f$ 在 $\Fsp$ 中的范数为
\[
  \|f\|=\sum_{|\alpha|\le k}\|D^\alpha f\|_{L^p(\Om)}
  +\sum_{|\beta|=k}\Big(\int_\Om J_{\beta,f}(x)^{p/q}\,dx\Big)^{1/p},
\]
允许取值为 $+\infty$；$f$ 属于该空间当且仅当此量有限。

\emph{第一步（边界层具有无限测度）}. 由已验证的前置附录 A（关于 $\nuA$ 的二分性），对每个 $r>0$ 都有 $\nuA(\Omr)=\infty$，其中 $\Omr=\{x\in\Om : d(x)<r\}$。为完整起见，给出一行证明：
\[
  \nuA(\Om)=\nuA(\Omr)+\int_{\Om\setminus\Omr} d(y)^{-aq}\,dy
\]
而在 $\Om\setminus\Omr$ 上 $d(y)\ge r$，故第二项不超过 $r^{-aq}|\Om|<\infty$；于是 $\nuA(\Omr)=\infty$。

\emph{第二步（$\sigma$-引理）}. 设 $\sigma\in(0,1)$，$g$ 为 $\Om$ 上的可测函数，且
\[
  \int_\Om\Big(\int_\Om |g(x)-g(y)|^q\,|x-y|^{-n-\sigma q}\,d(y)^{-aq}\,dy\Big)^{p/q}dx<\infty .
\]
则 $g$ 在 $\Om$ 上几乎处处等于常数。

\emph{证明}. 外层积分有限，故
\[
  J(x):=\int_\Om |g(x)-g(y)|^q\,|x-y|^{-n-\sigma q}\,d(y)^{-aq}\,dy
\]
对几乎处处的 $x\in\Om$ 有限。固定这样的 $x$，令 $r_x=d(x)/2>0$（因 $x$ 属于开集 $\Om$，故 $d(x)>0$）。对每个 $y\in\Om_{r_x}$ 有 $|x-y|\le\diam(\Om)$，从而 $|x-y|^{-n-\sigma q}\ge(\diam\Om)^{-n-\sigma q}$；于是
\begin{equation}
  \int_{\Om_{r_x}} |g(x)-g(y)|^q\,d(y)^{-aq}\,dy
  \le(\diam\Om)^{n+\sigma q}\,J(x)<\infty . \tag{2.1}
\end{equation}
对 $\varepsilon>0$，令 $A_\varepsilon(x)=\{y\in\Om_{r_x} : |g(x)-g(y)|>\varepsilon\}$。则
\begin{equation}
  \varepsilon^q\nuA\big(A_\varepsilon(x)\big)
  \le\int_{\Om_{r_x}} |g(x)-g(y)|^q\,d(y)^{-aq}\,dy
  \le(\diam\Om)^{n+\sigma q}\,J(x)<\infty , \tag{2.2}
\end{equation}
故 $\nuA(A_\varepsilon(x))<\infty$。由第一步知 $\nuA(\Om_{r_x})=\infty$，因此集合 $B_\varepsilon(x):=\Om_{r_x}\setminus A_\varepsilon(x)$ 满足 $\nuA(B_\varepsilon(x))=\infty$。现设 $x_1,x_2$ 为两个使 $J$ 有限之点，令 $r=\min(r_{x_1},r_{x_2})>0$，并令 $A=A_\varepsilon(x_1)\cup A_\varepsilon(x_2)$。两个 $A_\varepsilon(x_i)$ 的 $\nuA$-测度均有限，故 $\nuA(A)<\infty$；而 $\Om_r$ 同时含于 $\Om_{r_{x_1}}$ 与 $\Om_{r_{x_2}}$。于是 $B_\varepsilon(x_1)\cap B_\varepsilon(x_2)$ 包含 $\Om_r\setminus A$，其 $\nuA$-测度为 $\nuA(\Om_r)-\nuA(A)=\infty$，故非空。取该交集中的一点 $y$。则
\[
  |g(x_1)-g(x_2)|\le|g(x_1)-g(y)|+|g(y)-g(x_2)|\le 2\varepsilon .
\]
由 $\varepsilon>0$ 的任意性得 $g(x_1)=g(x_2)$。使 $J(x)<\infty$ 的点 $x$ 之集在 $\Om$ 中为满测集，故 $g$ 几乎处处为常数。

\emph{第三步（空间的刻画）}. 设 $f\in\Fsp$。则 $f\in W^{k,p}(\Om)$，且对每个满足 $|\beta|=k$ 的多重指标 $\beta$，函数 $D^\beta f$（任一代表）属于 $L^p(\Om)$，而范数有限给出 $\int_\Om J_{\beta,f}(x)^{p/q}\,dx<\infty$。于是对几乎处处的 $x$ 有 $J_{\beta,f}(x)<\infty$，对 $g=D^\beta f$ 应用第二步即知 $D^\beta f$ 在 $\Om$ 上几乎处处为常数。因此 $f$ 的所有 $k$ 阶分布导数均为常数。

下面证明：若 $f\in L^1_{\mathrm{loc}}(\Om)$ 的所有 $k$ 阶分布导数皆为常数，则该函数几乎处处等于一个次数不超过 $k$ 的多项式。取常数 $c_\beta$（$|\beta|=k$）与多项式 $P(x)=\sum_{|\beta|=k}c_\beta x^\beta/\beta!$（低阶项无关紧要）。则 $f-P$ 的所有 $k$ 阶分布导数均为零。设 $\rho$ 为标准磨光核，$\rho_\varepsilon(x)=\varepsilon^{-n}\rho(x/\varepsilon)$。对充分小的 $\varepsilon>0$，函数 $u_\varepsilon:=(f-P)*\rho_\varepsilon$ 在开集 $\Om_\varepsilon=\{x\in\Om : d(x)>\varepsilon\}$ 上光滑，且 $u_\varepsilon$ 的每个 $k$ 阶导数在其上为零；故 $u_\varepsilon$ 在 $\Om_\varepsilon$ 的每个连通分支上是次数不超过 $k-1$ 的多项式。固定紧含于 $\Om$ 的球 $B$。对所有充分小的 $\varepsilon$，$u_\varepsilon$ 在 $B$ 上是次数不超过 $k-1$ 的多项式，且当 $\varepsilon$ 趋于 $0$ 时 $u_\varepsilon$ 在 $L^1(B)$ 中趋于 $f-P$。在 $B$ 上写 $u_\varepsilon(y)=\sum_{|\gamma|\le k-1}a_\gamma(\varepsilon)y^\gamma$。由于次数不超过 $k-1$ 的多项式空间是有限维的，且 $L^1(B)$ 中的收敛蕴含所有系数收敛（与该有限维空间的对偶基作积分即可），系数 $a_\gamma(\varepsilon)$ 当 $\varepsilon$ 趋于 $0$ 时收敛到极限 $a_\gamma$。令 $Q_B(y)=\sum_{|\gamma|\le k-1}a_\gamma y^\gamma$。则 $u_\varepsilon$ 在 $L^1(B)$ 中趋于 $Q_B$，故 $f-P=Q_B$ 在 $B$ 上几乎处处成立。现用紧含于 $\Om$ 的球 $B$ 覆盖 $\Om$。若两个这样的球 $B$ 与 $B'$ 相交，则多项式 $Q_B$ 与 $Q_{B'}$ 在非空开集 $B\cap B'$ 上几乎处处相等，因而恒等。由于 $\Om$ 连通，一串相交的球可连接覆盖中任意两个球，故所有多项式 $Q_B$ 都等于同一个多项式 $Q$。因此 $f=P+Q$ 在 $\Om$ 上几乎处处成立。这就证明了 $\Fsp$ 含于 $P_k(\Om)$。

反之，若 $f$ 是总次数不超过 $k$ 的多项式，则对每个 $|\beta|=k$，导数 $D^\beta f$ 为常数，故 $J_{\beta,f}=0$ 恒成立，半范数项消失；又因 $\Om$ 有界，$f$ 属于 $L^p(\Om)$，且对一切 $|\alpha|\le k$ 有 $D^\alpha f$ 属于 $L^p(\Om)$。于是 $f\in\Fsp$。因此 $\Fsp=P_k(\Om)$。

\emph{第四步（推论）}. (a) 每个多项式都属于 $C^\infty(\Om)$，故 $C^\infty(\Om)\cap F^s=F^s$，且该集合在 $F^s$ 中的闭包 $F^{s,2}$ 为 $F^s$。(b) 设 $g\in C_c^\infty(\Om)\cap F^s$。由 $F^s=P_k$，存在次数不超过 $k$ 的多项式 $Q$ 使 $g=Q$ 在 $\Om$ 上几乎处处成立。因 $g$ 在 $\Om$ 中紧支，存在 $\delta>0$ 使 $g=0$ 在非空开集 $\Om_\delta$ 上成立。于是 $Q=0$ 在 $\Om_\delta$ 上几乎处处成立；在非空开集上几乎处处为零的多项式恒为零，故 $Q=0$，从而 $g=0$ 几乎处处成立。于是 $C_c^\infty(\Om)\cap F^s=\{0\}$，进而 $F^{s,1}=\{0\}$。证毕。
\end{proof}

\begin{remark}[直觉]
当权重 $d^{-aq}$ 不可积时，分数阶半范数中的内层 $y$-积分面对的是边界点的无限质量。除非 $D^\beta f$ 为常数，对几乎每个 $x$ 差值 $|D^\beta f(x)-D^\beta f(y)|$ 都在一个具有无限加权测度的集合上为正，从而使半范数发散。因此只有次数不超过 $k$ 的多项式得以存活；光滑函数于是自动稠密，而紧支光滑函数全部被排除，因为次数不超过 $k$ 的紧支多项式必恒为零。
\end{remark}

%==================================================================
\subsection{非可积权下的退化定理（第二版本）}

\begin{theorem}[不可积权情形下的退化定理]
设 $n\ge 1$，$\Om$ 为 $\R^n$ 中的非空有界域，$d(x)=\dist(x,\partial\Om)$，$a>0$，$p,q\in[1,\infty)$，并设 $s\in(0,\infty)\setminus\N$，其中 $k=\lfloor s\rfloor$、$\sigma=s-k\in(0,1)$。对可测函数 $f:\Om\to\R$ 定义
\[
  \|f\|_{F^s}:=\|f\|_{W^{k,p}(\Om)}
  +\sum_{|\beta|=k}\Big(\int_\Om\Big(\int_\Om |D^\beta f(x)-D^\beta f(y)|^q\,|x-y|^{-n-\sigma q}\,d(y)^{-aq}\,dy\Big)^{p/q}dx\Big)^{1/p},
\]
若 $f\notin W^{k,p}(\Om)$ 或某项为无穷则取值为 $+\infty$，其中 $\|f\|_{W^{k,p}(\Om)}$ 是 $f$ 的不超过 $k$ 阶弱导数的 $L^p(\Om)$ 范数之和。令 $\Fsp:=\{f\in L^p(\Om) : \|f\|_{F^s}<\infty\}$。假设 $\int_\Om d(y)^{-aq}\,dy=\infty$（二分性附录 A 的替代情形 (II)）。则：
\begin{enumerate}[label=(\roman*)]
\item $\Fsp=P_k$，其中 $P_k$ 是 $\R^n$ 上次数不超过 $k$ 的多项式在 $\Om$ 上限制所成的集合。
\item $C^\infty(\Om)\cap\Fsp=\Fsp$，故 $C^\infty(\Om)\cap\Fsp$ 在 $\Fsp$ 中的闭包 $F^{s,2}_{p,q}(\Om,d^{-aq})$ 等于 $\Fsp$。
\item $C_c^\infty(\Om)\cap\Fsp=\{0\}$，故 $C_c^\infty(\Om)\cap\Fsp$ 在 $\Fsp$ 中的闭包 $F^{s,1}_{p,q}(\Om,d^{-aq})$ 等于 $\{0\}$。
\end{enumerate}
特别地，只要 $P_k$ 不是 $\{0\}$，$C^\infty(\Om)$ 就在 $\Fsp$ 中稠密，而 $C_c^\infty(\Om)$ 不稠密。
\end{theorem}

\begin{proof}
设 $n\ge1$，$\Om$ 为 $\R^n$ 中的非空有界域，$d(x)=\dist(x,\partial\Om)$，$a>0$，$p,q\in[1,\infty)$。固定 $s\in(0,\infty)\setminus\N$，令 $k:=\lfloor s\rfloor$、$\sigma:=s-k\in(0,1)$。对可测函数 $f:\Om\to\R$ 令
\[
  \|f\|_{F^s}:=\|f\|_{W^{k,p}(\Om)}
  +\sum_{|\beta|=k}\Big(\int_\Om\Big(\int_\Om |D^\beta f(x)-D^\beta f(y)|^q\,|x-y|^{-n-\sigma q}\,d(y)^{-aq}\,dy\Big)^{p/q}dx\Big)^{1/p},
\]
若 $f\notin W^{k,p}(\Om)$ 或某项为无穷则取值为 $+\infty$，其中 $\|f\|_{W^{k,p}(\Om)}:=\sum_{|\alpha|\le k}\|D^\alpha f\|_{L^p(\Om)}$。令 $F^s:=\Fsp:=\{f\in L^p(\Om) : \|f\|_{F^s}<\infty\}$。假设 $\nuA(\Om)=\int_\Om d(y)^{-aq}\,dy=\infty$（附录 A 的替代情形 (II)）。令 $D:=\diam(\Om)\in(0,\infty)$，于是对一切 $x,y\in\Om$ 有 $|x-y|\le D$，从而对一切 $x,y\in\Om$ 有 $|x-y|^{-n-\sigma q}\ge D^{-n-\sigma q}$。

\emph{第一步（多项式属于空间）}. 设 $P\in P_k$。因 $\Om$ 有界，$P$ 在 $\Om$ 上有界，故 $P\in W^{k,p}(\Om)$ 且范数有限。对每个满足 $|\beta|=k$ 的 $\beta$，弱导数 $D^\beta P$ 为常数，故对一切 $x,y$ 有 $D^\beta P(x)-D^\beta P(y)=0$，相应的内层积分恒为零。于是 $\|P\|_{F^s}<\infty$ 且 $P\in F^s$；从而 $P_k$ 含于 $F^s$。

\emph{第二步（只有多项式属于空间）}. 设 $f\in F^s$。则 $f\in W^{k,p}(\Om)$，且对每个满足 $|\beta|=k$ 的 $\beta$，可测函数 $g_\beta:=D^\beta f$ 满足
\[
  \int_\Om\Big(\int_\Om |g_\beta(x)-g_\beta(y)|^q\,|x-y|^{-n-\sigma q}\,\nuA(dy)\Big)^{p/q}dx<\infty ,
\]
因为这是 $\|f\|_{F^s}$ 中相应于 $\beta$ 的求和项。附录 A（SIGMA-LEMMA）的假设均成立（$\Om$ 为非空有界域，$a>0$，$q\in[1,\infty)$，$\sigma\in(0,1)$，$p\in[1,\infty)$，$\nuA(\Om)=\infty$，$g_\beta$ 可测），故 $g_\beta=D^\beta f$ 在 $\Om$ 上几乎处处为常数。由于这对每个满足 $|\beta|=k$ 的 $\beta$ 都成立，将附录 A（POLYNOMIAL LEMMA）应用于 $f\in W^{k,p}(\Om)$，即得 $\R^n$ 上次数不超过 $k$ 的多项式 $P$ 使 $f=P$ 在 $\Om$ 上几乎处处成立，即 $f\in P_k$。于是 $F^s$ 含于 $P_k$，结合第一步得 $F^s=P_k$。

\emph{第三步（推论 (a)）}. 每个多项式都在 $\Om$ 上无穷次可微，故 $P_k$ 含于 $C^\infty(\Om)\cap F^s$；由第二步，$F^s=P_k$ 含于 $C^\infty(\Om)\cap F^s$。于是 $C^\infty(\Om)\cap F^s=F^s$，且 $C^\infty(\Om)\cap F^s$ 在 $F^s$ 中的闭包 $F^{s,2}$ 等于 $F^s$。

\emph{第四步（推论 (b)）}. 设 $g\in C_c^\infty(\Om)\cap F^s$，并设 $g$ 不恒为零。令 $\delta:=\dist(\supp g,\partial\Om)>0$；则 $g=0$ 在边界层 $\Om_\delta:=\{y\in\Om : d(y)<\delta\}$ 上成立，且每个 $x\in\supp g$ 满足 $d(x)\ge\delta$。回顾初等估计：对一切 $x,y\in\Om$ 有 $|x-y|^{-n-\sigma q}\ge D^{-n-\sigma q}$。

\emph{情形 $k=0$}. 集合 $U:=\{x\in\Om : g(x)\ne 0\}$ 非空（因 $g$ 不恒为零）且为开集（因 $g$ 连续），故其 Lebesgue 测度为正；又 $U$ 含于 $\supp g$，故对 $x\in U$ 有 $d(x)\ge\delta$。固定 $x\in U$。因 $g=0$ 在 $\Om_\delta$ 上成立，
\begin{align*}
  \int_\Om |g(x)-g(y)|^q\,|x-y|^{-n-\sigma q}\,\nuA(dy)
  &\ge D^{-n-\sigma q}\int_{\Om_\delta}|g(x)-g(y)|^q\,\nuA(dy)\\
  &=D^{-n-\sigma q}|g(x)|^q\nuA(\Om_\delta)=\infty ,
\end{align*}
其中 $\nuA(\Om_\delta)=\infty$ 由附录 A 给出。故对每个 $x\in U$ 内层积分均为无穷，于是
\[
  \int_\Om\big(\text{inner integral}\big)^{p/q}dx\ge|U|\cdot(+\infty)=+\infty ,
\]
这与 $g\in F^s$ 相矛盾。
\emph{情形 $k\ge1$}. 若每个满足 $|\beta|=k$ 的弱导数 $D^\beta g$ 都恒为零，则由附录 A，$g$ 几乎处处与某多项式一致；因 $g$ 连续且 $\Om$ 为连通开集，$g$ 在 $\Om$ 上将等于该多项式，而在有界域 $\Om$ 中具紧支的多项式恒为零，这与 $g\ne0$ 矛盾。故存在满足 $|\beta|=k$ 的 $\beta$ 使 $h:=D^\beta g$ 不恒为零。则 $h\in C_c^\infty(\Om)$，$\supp h$ 含于 $\supp g$，且 $h=0$ 在 $\Om_\delta$ 上成立。集合 $U:=\{x\in\Om : h(x)\ne0\}$ 非空且为开集，故测度为正，且含于 $\supp h$，进而含于 $\supp g$，故对 $x\in U$ 有 $d(x)\ge\delta$。固定 $x\in U$。因 $h=0$ 在 $\Om_\delta$ 上成立且 $|x-y|^{-n-\sigma q}\ge D^{-n-\sigma q}$，
\[
  \int_\Om |h(x)-h(y)|^q\,|x-y|^{-n-\sigma q}\,\nuA(dy)
  \ge D^{-n-\sigma q}|h(x)|^q\nuA(\Om_\delta)=\infty ,
\]
此处再次用到附录 A。故对每个 $x\in U$ 内层积分均为无穷，并且
\[
  \int_\Om\big(\text{inner integral}\big)^{p/q}dx\ge|U|\cdot(+\infty)=+\infty ,
\]
这与 $\|g\|_{F^s}$ 中第 $\beta$ 项有限相矛盾。

两种情形都得到 $g=0$，即 $C_c^\infty(\Om)\cap F^s=\{0\}$，且 $C_c^\infty(\Om)\cap F^s$ 在 $F^s$ 中的闭包 $F^{s,1}$ 为 $\{0\}$。
\end{proof}

\begin{remark}[直觉]
若边界权重不可积，则对每个使某个最高阶导数非零的 $x$，内层积分都在一薄边界层上看到无限多个使该导数为零的点 $y$，而核函数有下界 $\diam(\Om)^{-n-\sigma q}$；于是内层积分在该导数非零的非空开集上处处为无穷，外层积分随之发散。
\end{remark}

%==================================================================
\subsection{修改后的高阶内在空间的退化}

\begin{theorem}[修改后的高阶内蕴空间的退化]
设 $n\ge1$，$\Om$ 为 $\R^n$ 中的非空有界域，对 $x\in\Om$ 令 $d(x)=\dist(x,\partial\Om)$，设 $k\ge1$ 为整数，$\sigma\in(0,1)$，$p,q\in[1,\infty)$，并令 $s:=k+\sigma$。假设 $\int_\Om d(y)^{-kq}\,dy=\infty$。设 $f$ 为 $\Om$ 上的 Lebesgue 可测函数，其所有满足 $|\beta|\le k$ 的弱导数 $D^\beta f$ 均存在且 Lebesgue 可测，并令
\[
  S(f):=\sum_{|\beta|\le k}\Big(\int_\Om\Big[\int_\Om |D^\beta f(x)-D^\beta f(y)|^q\,|x-y|^{-n-\sigma q}\,d(y)^{-(k-|\beta|)q}\,dy\Big]^{p/q}dx\Big)^{1/p},
\]
其中若相应积分非有限，则该求和项规定为 $+\infty$。若 $S(f)<\infty$，则 $f$ 在 $\Om$ 上 Lebesgue 几乎处处为常数。于是：
\begin{enumerate}[label=(\roman*)]
\item 参考论文定义 1.3 中修改后的加权内蕴 Triebel--Lizorkin 空间（其范数为 $\|f\|_{F^s_{p,q}(\Om)}+S(f)$，其中 $\|\cdot\|_{F^s_{p,q}(\Om)}$ 为无权 Triebel--Lizorkin 范数）恰为常数函数构成的一维空间 $\R\cdot\mathbf{1}$；
\item 使 Hardy 加强总量
\[
  T(f):=S(f)+\sum_{|\beta|\le k}\Big(\|D^\beta f/d^{s-|\beta|}\|_{L^p(\Om)}+\|D^\beta f/d^{s-|\beta|+n/q}\|_{L^q(\Om)}\Big)
\]
有限的可测函数构成的空间为零空间 $\{0\}$。
\end{enumerate}
\end{theorem}

\begin{proof}
以下 $\nuA$ 表示 $\Om$ 上以 $d(y)^{-aq}$ 为密度的 Borel 测度，故 $\nuA(\Om)=\int_\Om d(y)^{-aq}\,dy$。命题中的假设是
\[
  \nu_k(\Om)=\int_\Om d(y)^{-kq}\,dy=\infty .
\]

\emph{第一步（$\beta=0$ 的项与 $\sigma$-引理）}. $S(f)$ 中相应于 $\beta=0$ 的求和项为
\[
  \Big(\int_\Om\Big[\int_\Om |f(x)-f(y)|^q\,|x-y|^{-n-\sigma q}\,d(y)^{-kq}\,dy\Big]^{p/q}dx\Big)^{1/p},
\]
它不超过 $S(f)<\infty$。故非负可测函数
\[
  x\mapsto J(x):=\int_\Om |f(x)-f(y)|^q\,|x-y|^{-n-\sigma q}\,d(y)^{-kq}\,dy
\]
满足 $\int_\Om J(x)^{p/q}\,dx<\infty$。取 $a:=k$（于是 $\sigma$-引理中的权 $d(y)^{-aq}$ 即 $d(y)^{-kq}$），$n$、$q$、$\sigma$、$p$、$\Om$ 均不变，$g:=f$，此时 $\sigma$-引理（附录 A）还要求假设 $\nuA(\Om)=\infty$，而这正是我们的假设 $\nu_k(\Om)=\infty$。$\sigma$-引理的其余假设（$n\ge1$，$a=k>0$，$q\in[1,\infty)$，$\sigma\in(0,1)$，$p\in[1,\infty)$，$\Om$ 为非空有界域，$g=f$ 为 Lebesgue 可测）均成立。其结论给出 $f$ 在 $\Om$ 上 Lebesgue 几乎处处为常数。

\emph{第二步（(i)）}. 设 $X$ 为范数 $\|f\|_{F^s_{p,q}(\Om)}+S(f)$ 有限的可测函数 $f$ 构成的空间；这就是参考论文定义 1.3 中修改后的加权内蕴空间（无权范数 $\|\cdot\|_{F^s_{p,q}(\Om)}$ 即公式 (1.3) 中取 $a=0$ 的范数）。若 $f\in X$，则 $S(f)<\infty$，故由第一步 $f$ 几乎处处为常数；于是 $X$ 含于常数函数构成的空间 $\R\cdot\mathbf{1}$（模去几乎处处相等）。反之，设 $c$ 为实常数，在 $\Om$ 上 $f=c$。$f$ 的所有正阶弱导数均为零，且对每个多重指标 $\beta$，$D^\beta f$ 或为 $0$（若 $|\beta|\ge1$）或为常数 $c$（若 $\beta=0$）；两种情形下 $D^\beta f$ 在连通集 $\Om$ 上均为常数，故对一切 $x,y\in\Om$ 有 $|D^\beta f(x)-D^\beta f(y)|=0$，$S(f)$ 的每个求和项均为零，即 $S(f)=0$。此外，在无权范数 $\|\cdot\|_{F^s_{p,q}(\Om)}$ 中一切正阶导数均为零，一切最高阶差分半范数均为零，故
\[
  \|f\|_{F^s_{p,q}(\Om)}=\|c\|_{L^p(\Om)}=|c|\,|\Om|^{1/p}<\infty .
\]
于是 $f\in X$，且 $X=\R\cdot\mathbf{1}$。

\emph{第三步（(ii)）}. 设 $Y$ 为满足 $T(f)<\infty$ 的可测函数 $f$ 构成的空间，其中 $T$ 为命题中的 Hardy 加强总量。若 $f\in Y$，则 $S(f)\le T(f)<\infty$，故由第一步，$f=c$ 在 $\Om$ 上几乎处处为常数。对该常数，$T(f)$ 中满足 $|\beta|\ge1$ 的求和项均为零，剩余的求和项为
\[
  \|c/d^{s}\|_{L^p(\Om)}=|c|\Big(\int_\Om d(y)^{-sp}\,dy\Big)^{1/p}
\]
以及
\[
  \|c/d^{s+n/q}\|_{L^q(\Om)}
  =|c|\Big(\int_\Om d(y)^{-(s+n/q)q}\,dy\Big)^{1/q}
  =|c|\Big(\int_\Om d(y)^{-sq-n}\,dy\Big)^{1/q}.
\]
我们断言 $\int_\Om d(y)^{-sq-n}\,dy=\infty$。事实上，对 $y\in\Om$ 有 $d(y)\le\diam(\Om)=:D$，其中 $0<D<\infty$（$\Om$ 有界），且因 $s=k+\sigma$ 而 $\sigma>0$，故 $sq+n>kq$；于是
\[
  d(y)^{-sq-n}=d(y)^{-kq}\,d(y)^{-(sq+n-kq)}\ge D^{-(sq+n-kq)}\,d(y)^{-kq},
\]
故
\[
  \int_\Om d(y)^{-sq-n}\,dy\ge D^{-(sq+n-kq)}\int_\Om d(y)^{-kq}\,dy=\infty
\]
由假设。因此 $T(f)<\infty$ 迫使 $|c|=0$，即 $f=0$ 几乎处处成立，从而 $Y=\{0\}$。
\end{proof}

\begin{remark}[直觉]
修改后范数中 $\beta=0$ 的项就是 $f$ 本身带权 $d^{-kq}$ 的加权分数阶半范数。由于 $k\ge1$ 且 $q\ge1$，指数 $kq$ 至少为 $1$，故在 Lipschitz 型区域上该权在边界附近不可积；$\sigma$-引理于是迫使 $f$ 为常数。可见为吸收 Leibniz 法则的溢出而引入的额外距离权 $d^{-(k-|\beta|)q}$ 过度了：在 $\beta=0$ 处它除常数外一概消灭。
\end{remark}

%==================================================================
\subsection{整数阶修改空间的非退化性}

\begin{theorem}[整数阶修正空间的非退化（修正版）]
\label{thm:int-nondegen}
设 $\Om$ 为 $\R^n$ 中的非空有界域，$k\in\N$，$p,q\in[1,\infty)$。设 $F^k_{p,q}(\Om)$ 为 CLWZ 定义 1.6 中修正后的加权 Triebel--Lizorkin 空间，其范数为
\[
  \|f\|_{F^k_{p,q}}=\|f\|_{L^p(\Om)}+\sum_{r=0}^k[f]_{r,k}+\|f/d^k\|_{L^p(\Om)},
\]
其中
\[
  [f]_{r,k}=\Big(\int_\Om\Big(\int_{|h|<d(x)/(2(k+1))}|\Delta_h^{k+1-r}f(x)|^q\,|h|^{-n-(k-r)q}\,d(x)^{-rq}\,dh\Big)^{p/q}dx\Big)^{1/p},
\]
$\Delta_h^m f(x)=\sum_{\nu=0}^m(-1)^{m-\nu}\binom{m}{\nu}f(x+\nu h)$，且 $d(x)=\dist(x,\partial\Om)$。
则
\[
  C_c^\infty(\Om)\subset F^k_{p,q}(\Om),
\]
从而 $F^k_{p,q}(\Om)\not\subset P_{k-1}(\Om)$，后者是总次数不超过 $k-1$ 的多项式在 $\Om$ 上限制所成的集合。

特别地，若对某个 $a>0$ 有
\[
  \nuA(\Om)=\int_\Om d(x)^{-aq}\,dx=\infty ,
\]
则对每个非整数 $s$，分数阶空间 $\Fsp$ 都退化为多项式（附录 A），而整数阶修正空间 $F^k_{p,q}(\Om)$ 仍含有非多项式函数。
\end{theorem}

\begin{proof}
\textbf{第一步（witness 的选择）}。原事实取 $f_0(x)=d(x)^k$ 作 witness，该选择不成立（原因见评注 \ref{rem:int-why-fail} 与 \ref{rem:int-counter}）。本证明改取 $\varphi\in C_c^\infty(\Om)$（该论证已由独立复核确认，附录 A）。差别在于：$\varphi$ 光滑且支集紧，其 $m$ 阶差分具有\emph{全阶}消失量 $|h|^m$，恰好抵消核 $|h|^{-n-(k-r)q}$ 的奇性；而 $d$ 仅有 Lipschitz 正则性，$d^k$ 的差分至多只有一阶消失量，无法抵消。

\textbf{第二步（$L^p$ 项与 Hardy 权项）}。取 $\delta>0$ 使 $\supp\varphi\subset\{x\in\Om:\ d(x)\ge\delta\}$。则
\[
  \|\varphi\|_{L^p(\Om)}<\infty,\qquad
  \Big\|\frac{\varphi}{d^k}\Big\|_{L^p(\Om)}\le\delta^{-k}\|\varphi\|_{L^p(\Om)}<\infty .
\]

\textbf{第三步（半范数：抵消核的奇性）}。固定 $r\in\{0,\dots,k\}$，记 $m=k+1-r$、$c=1/(2(k+1))$，并令
\[
  M_m:=\max_{|\alpha|=m}\|D^\alpha\varphi\|_{L^\infty(\Om)}<\infty .
\]
由 $m$ 重积分表示（或对 $\Delta_h^m$ 作 Taylor 展开）得
\begin{equation}\label{eq:int-diff}
  |\Delta_h^m\varphi(x)|\le C_m\,M_m\,|h|^m\qquad (x\in\Om,\ |h|\le1).
\end{equation}
于是在 $|h|<c\,d(x)$ 上
\[
  |\Delta_h^m\varphi(x)|^q\,|h|^{-n-(k-r)q}
  \le (C_mM_m)^q\,|h|^{\,mq-n-(k-r)q}
  =(C_mM_m)^q\,|h|^{\,q-n},
\]
这里用到 $mq-(k-r)q=(k+1-r)q-(k-r)q=q$。又因 $q\ge1$，
\[
  \int_{|h|<c\,d(x)}|h|^{\,q-n}\,dh=\frac{\omega_{n-1}}{q}\big(c\,d(x)\big)^q<\infty .
\]

\textbf{第四步（支集传播：$d(x)$ 的下界）}。第三步中的估计不能直接在 $\supp\varphi$ 上使用，因为当
$\Delta_h^m\varphi(x)\ne0$ 时，点 $x$ 本身未必落在 $\supp\varphi$ 内。注意到
\[
  \Delta_h^m\varphi(x)=\sum_{\nu=0}^{m}(-1)^{m-\nu}\binom{m}{\nu}\varphi(x+\nu h),
\]
故若它非零，则存在 $\nu\in\{0,\dots,m\}$ 使 $x+\nu h\in\supp\varphi$，从而 $d(x+\nu h)\ge\delta$。又 $\nu\le m=k+1-r\le k+1$ 且 $|h|<c\,d(x)=d(x)/(2(k+1))$，故
\[
  \nu|h|\le(k+1)\cdot\frac{d(x)}{2(k+1)}=\frac{d(x)}{2}.
\]
由 $d$ 的 $1$-Lipschitz 性得 $d(x+\nu h)\le d(x)+|x-(x+\nu h)|=d(x)+\nu|h|\le\tfrac32d(x)$，于是
\begin{equation}\label{eq:int-support}
  d(x)\ \ge\ \tfrac23\,\delta\qquad\text{只要}\ \ \Delta_h^m\varphi(x)\ne0 .
\end{equation}
特别地，$d(x)^{-rq}\le(2\delta/3)^{-rq}$ 在 $\{\Delta_h^m\varphi(x)\ne0\}$ 上一致有界（$r\ge0$）。

\textbf{第五步（把 $h$ 的积分分成两段）}。$|h|$ 可能大于 $1$（当 $d(x)$ 较大时），此时不必再靠 $|h|^m$ 的小尺度消失，而用平凡界即可。由 \eqref{eq:int-diff}（该式要求 $|h|\le1$）与差分的平凡估计，
\[
  |\Delta_h^m\varphi(x)|^q\le
  \begin{cases}
    (C_mM_m)^q\,|h|^{mq}, & |h|\le1,\\[2pt]
    \big(2^m\|\varphi\|_{L^\infty(\Om)}\big)^q=:(C_m')^q, & |h|>1 .
  \end{cases}
\]
在 $|h|\le1$ 段，第三步已给出核 $|h|^{-n-(k-r)q}$ 与 $|h|^{mq}$ 相消为 $|h|^{q-n}$，其在 $|h|\le1$ 上可积。
在 $|h|>1$ 段，
\[
  \int_{1<|h|<c\,d(x)}|h|^{-n-(k-r)q}\,dh\ \le\ \frac{\omega_{n-1}}{(k-r)q}<\infty\quad(r<k),
\]
而当 $r=k$（此时指数为 $-n$）有 $\int_{1<|h|<c\,d(x)}|h|^{-n}dh\lesssim\log\big(1+c\,d(x)\big)$，又 $d(x)\le\diam\Om<\infty$，故亦有限。

\textbf{第六步（对 $x$ 积分）}。由第四步，只需在集合 $E_x:=\{x:\Delta_h^m\varphi(x)\ne0\}$ 上估计，而在其上 $d(x)\ge\tfrac23\delta$ 且 $d(x)\le\diam\Om$。合并第三、五步，存在仅依赖于 $\Om,m,q,\delta,M_m,\|\varphi\|_{L^\infty}$ 的常数 $C''$ 使
\[
  \int_{|h|<c\,d(x)}|\Delta_h^{k+1-r}\varphi(x)|^q\,|h|^{-n-(k-r)q}\,d(x)^{-rq}\,dh\ \le\ C''
\qquad(x\in\Om),
\]
于是
\[
  [\varphi]_{r,k}\le (C'')^{1/q}|\Om|^{1/p}<\infty .
\]
对有限多个 $r=0,\dots,k$ 合并，得 $\|\varphi\|_{F^k_{p,q}}<\infty$，故 $\varphi\in F^k_{p,q}(\Om)$。

\textbf{第七步（非退化）}。取非零的 $\varphi\in C_c^\infty(\Om)$。若 $\varphi$ 在 $\Om$ 上几乎处处等于某个多项式 $P$，则 $P$ 必在非空开集 $\Om\setminus\supp\varphi$ 上为零，从而 $P\equiv0$，与 $\varphi\ne0$ 矛盾。故 $\varphi\notin P_{k-1}(\Om)$，于是 $F^k_{p,q}(\Om)\not\subset P_{k-1}(\Om)$。

最后一项断言：附录 A 给出对每个非整数 $s$ 有 $\Fsp=P_{\lfloor s\rfloor}$；而 $F^k_{p,q}(\Om)$ 的定义不涉及参数 $a$，且如上已证它含有非多项式函数。
\end{proof}

\begin{remark}[原证明为何失效：一个发散的积分]\label{rem:int-why-fail}
原事实的证明在 $r<k$（此时 $(k-r)q>0$）时断言
\[
  \int_{|h|<c\,d(x)}|h|^{-n-(k-r)q}\,dh=C_2(n,q,k)\,d(x)^{-(k-r)q}.
\]
该等式\emph{不成立}。令 $\alpha=(k-r)q>0$，由极坐标公式，对任意 $R>0$
\[
  \int_{|h|<R}|h|^{-n-\alpha}\,dh
  =\omega_{n-1}\int_0^R\rho^{-1-\alpha}\,d\rho=+\infty .
\]
故左端恒为 $+\infty$，不可能等于有限量 $C_2d(x)^{-(k-r)q}$。

错误的来源是：原证明只得到与 $h$ 无关的粗糙界 $|\Delta_h^{k+1-r}f_0(x)|\le C_1d(x)^k$，其中\emph{没有任何随 $h\to0$ 消失的因子}。面对 $|h|^{-n-(k-r)q}$ 这样的奇性核，仅靠与 $h$ 无关的上界无法保证可积性。要使内层积分有限，必须从差分中提取出足够高阶的 $|h|$ 消失量；\eqref{eq:int-diff} 中 $m$ 阶全阶消失的做法正是如此。
\end{remark}

\begin{remark}[原 witness $d^k$ 的适用范围：$k=1$ 可用，$k\ge2$ 失效]\label{rem:int-counter}
原事实以 $f_0(x)=d(x)^k$ 作 witness。经独立复核（附录 A 中相应的两条记录），
其适用范围如下。本节初稿曾断言 $k=1$ 时 $[d]_{0,1}=+\infty$，该断言\defterm{有误}、现已撤回，理由见 \S\ref{sec:revision}。

\emph{(a) $k=1$ 时 $d$ 确实属于 $F^1_{p,q}((0,1))$。}
取 $n=1$、$\Om=(0,1)$、$d(x)=\min(x,1-x)$，并记 $\lambda(x):=|x-\tfrac12|$。对 $r=0$ 项（此时 $k+1-r=2$，核指数为 $-1-q$）有
\[
  |\Delta_h^2d(x)|\le 2\lambda(x),\qquad
  \Delta_h^2d(x)=0\ \ \text{只要}\ |h|\le\tfrac12\lambda(x).
\]
其关键是：$d$ 在 $(0,\tfrac12)$ 与 $(\tfrac12,1)$ 上分别是恒等映射与仿射映射，故只要三点 $x,x+h,x+2h$ 不跨越中点 $\tfrac12$，二阶差分\emph{恒为零}。于是对每个固定的 $x$，集合 $\{h:\Delta_h^2d(x)\ne0\}$ 与 $0$ 有正距离，内层积分在 $|h|\to0$ 处\emph{不}发散，并可得显式界
\[
  I(x):=\int_{|h|<d(x)/4}|\Delta_h^2d(x)|^q\,|h|^{-1-q}\,dh\le\frac{2^{2q+1}}{q}.
\]
对 $r=1$ 项，由 $d$ 为 $1$-Lipschitz 得 $|\Delta_h^1d(x)|\le|h|$，从而该半范数不超过 $(2/(4^qq))^{1/q}$。故 $[d]_{0,1},[d]_{1,1}<\infty$，即 $d\in F^1_{p,q}((0,1))$。
换言之，原 witness 在 $k=1$ 时\emph{是}可用的。

\emph{(b) $k\ge2$ 时 $d^k$ 不属于 $F^k_{p,q}((0,1))$。}
令 $\lambda_0:=1/(8(k+1))$，取 $\lambda\in(0,\lambda_0)$、$x:=\tfrac12-\lambda$、$T:=d(x)/(2(k+1))$；则 $d(x)=x$ 且 $T\asymp\lambda$。
对 $h\in(\lambda,T)$ 与 $\nu=1,\dots,k+1$ 有 $x+\nu h\in(\tfrac12,1)$，故 $d(x+\nu h)=1-x-\nu h$。令 $S(t):=(1-t)^k$，这是次数为 $k$ 的多项式，于是
\[
  \Delta_h^{k+1}f_0(x)
  =(-1)^{k+1}\big(x^k-(1-x)^k\big)+\Delta_h^{k+1}S(x)
  =(-1)^{k+1}\big(x^k-(1-x)^k\big),
\]
因为 $(k+1)$ 阶差分把次数为 $k$ 的多项式打为零，即 $\Delta_h^{k+1}S\equiv0$。
由中值定理，$\big|x^k-(1-x)^k\big|\asymp\lambda$。又因 $T\asymp\lambda$，
\[
  \int_{\lambda<|h|<2\lambda}|h|^{-1-kq}\,dh\ \asymp\ \lambda^{-kq},
\]
故（把 $h$ 限制在 $(\lambda,2\lambda)$ 上即可）
\[
  \int_{|h|<T}|\Delta_h^{k+1}f_0(x)|^q\,|h|^{-1-kq}\,dh
  \ \gtrsim\ \lambda^{q}\cdot\lambda^{-kq}
  =\lambda^{-(k-1)q}.
\]
由于 $(k-1)q\ge1$（由 $k\ge2$ 与 $q\ge1$），从而
\[
  \int_0^{\lambda_0}\lambda^{-(k-1)q}\,d\lambda=+\infty .
\]
故 $[f_0]_{0,k}=+\infty$，即 $f_0\notin F^k_{p,q}((0,1))$。

\emph{结论}：原 witness 仅在 $k=1$ 时有效，$k\ge2$ 时必须更换。本文统一改用 $C_c^\infty$ 作 witness（见本节定理 \ref{thm:int-nondegen}），该证明已由独立复核（附录 A）确认。
\end{remark}

\begin{remark}[与 Poincar\'e 障碍并不矛盾]
附录 A 指出：当 $s>1$ 时，$C_c^\infty(\Om)$ 在\emph{未加强}的 $F^s$ 中绝不稠密。本定理断言的是 $C_c^\infty(\Om)\subset F^k_{p,q}(\Om)$，即这些函数\emph{属于}该空间，而非\emph{稠密}于其中。二者并不冲突：前者说闭包是 $F^s$ 的真子空间，后者说 $F^k_{p,q}$ 含有非平凡元素。
\end{remark}

\begin{remark}[直觉：为何 $k=1$ 可以，$k\ge2$ 不行]
范数用 $d(x)^{-rq}$ 给第 $r$ 个差分项加权、用 $d(x)^{-kq}$ 给零阶项加权。在半径与 $d(x)$ 相当的球上，
$d^k$ 被常数乘 $d(x)^k$ 控制，故\emph{每个}差分都被 $Cd(x)^k$ 控制；这一步对一切 $k$ 都成立，
但它\defterm{不足以}保证内层积分收敛，因为核 $|h|^{-n-(k-r)q}$ 在 $h=0$ 处是奇异的。

真正起决定作用的是差分随 $h\to0$ 的\emph{消失阶}。
当 $k=1$ 时 $r=0$ 项只需要二阶差分：$d$ 在中点两侧分别是恒等与仿射映射，故只要三点不跨越中点，二阶差分\emph{恒为零}；
非零集在 $h$ 方向与 $0$ 有正距离，配合 $r=k$ 项中 Lipschitz 性给出的额外因子 $|h|$，全部内层积分收敛。

当 $k\ge2$ 时 $r=0$ 项需要 $(k+1)$ 阶差分。取 $x=\tfrac12-\lambda$、$h\in(\lambda,T)$（$T\asymp\lambda$），
则 $x+\nu h\in(\tfrac12,1)$ 对所有 $1\le\nu\le k+1$ 成立，而 $t\mapsto(1-t)^k$ 是次数恰为 $k$ 的多项式，
其 $(k+1)$ 阶差分为零；于是 $(k+1)$ 阶差分只剩中点两侧的不平衡项 $(-1)^{k+1}(x^k-(1-x)^k)\asymp\lambda$，
它\defterm{不随 $h\to0$ 消失}。将 $h$ 限制在 $(\lambda,2\lambda)$ 上，则
\[
  \int_{\lambda}^{2\lambda}h^{-1-kq}\,dh=\frac{(2\lambda)^{-kq}-\lambda^{-kq}}{-kq}\asymp\lambda^{-kq},
\]
故内层积分的下界为 $\lambda^{q}\cdot\lambda^{-kq}=\lambda^{-(k-1)q}$；而 $(k-1)q\ge1$（因 $k\ge2$、$q\ge1$）使
$\int_0^{\lambda_0}\lambda^{-(k-1)q}d\lambda=+\infty$，即该积分在 $\lambda\to0$ 时发散。

因此 $d^k$ 只在 $k=1$ 时是有效的 witness；一般 $k$ 必须改用 $C_c^\infty$，那里差分的全阶消失量 $|h|^{k+1-r}$ 恰好抵消核的奇性。
\end{remark}

\subsection{整数阶空间的稠密性：截断引理}

\begin{theorem}[截断引理]
设 $n \ge 1$，设 $\Om$ 是 $\R^n$ 中的非空有界区域，对 $x \in \Om$ 令 $d(x) = \dist(x, \partial\Om)$；设 $k \in \N = \{1,2,3,\dots\}$，$p, q \in [1,\infty)$。设 $F^k_{p,q}(\Om)$ 是 CLWZ 定义 1.6 中的\defterm{修正加权 Triebel--Lizorkin 空间}，其范数定义为
\[
\|f\|_{F^k_{p,q}} := \|f\|_{L^p(\Om)} + \sum_{r=0}^{k} [f]_{r,k} + \|f/d^k\|_{L^p(\Om)},
\]
其中
\[
[f]_{r,k} := ( \int_\Om ( \int_{|h| < T(x)} |\Delta_h^{k+1-r} f(x)|^q |h|^{-n-(k-r)q} d(x)^{-r q} dh )^{p/q} dx )^{1/p},
\]
这里 $T(x) := d(x)/(2(k+1))$，且 $\Delta_h^m f(x) := \sum_{\nu=0}^{m} (-1)^{m-\nu} \binom{m}{\nu} f(x+\nu h)$。设 $\varepsilon > 0$，并设 $\chi_\varepsilon : \Om \to \R$ 为 $C^\infty$ 的，满足 $0 \le \chi_\varepsilon \le 1$，$\chi_\varepsilon = 1$ 于 $\{x \in \Om : d(x) \ge 2 \varepsilon\}$ 上，$\chi_\varepsilon = 0$ 于 $\{x \in \Om : d(x) < \varepsilon\}$ 上，且对每个多重指标 $\alpha$ 都有 $|D^\alpha \chi_\varepsilon| \le C_{|\alpha|} \varepsilon^{-|\alpha|}$ 于 $\Om$ 上（$C_{|\alpha|}$ 与 $\varepsilon$ 无关）；这样的函数族是存在的（显式构造见第一步）。令 $g_\varepsilon := 1 - \chi_\varepsilon$。则对每个 $f \in F^k_{p,q}(\Om)$：(i) $f \chi_\varepsilon$ 属于 $F^k_{p,q}(\Om)$；(ii) $\|f - f \chi_\varepsilon\|_{F^k_{p,q}(\Om)}^p$ 不超过 $C(n,k,p,q)$ 乘以如下 $k+3$ 个量之和，这些量都统一在同一个边界层 $\{x \in \Om : d(x) < 4 \varepsilon\}$ 上积分：$\int_{d<4\varepsilon}|f|^p dx$，$\int_{d<4\varepsilon}|f/d^k|^p dx$，以及对 $r = 0,\dots,k$ 的 $\int_{d<4\varepsilon} I_r(x)^{p/q} dx$，其中 $I_r(x) := \int_{|h|<T(x)}|\Delta_h^{k+1-r}f(x)|^q|h|^{-n-(k-r)q}d(x)^{-rq}dh$；特别地，当 $\varepsilon \to 0$ 时 $\|f - f \chi_\varepsilon\|_{F^k_{p,q}(\Om)} \to 0$。从而线性子空间 $F^k_{p,q,c}(\Om) := \{ f \in F^k_{p,q}(\Om) : f \mbox{ 本性支撑于 } \Om \mbox{ 的某个紧子集} \}$ 在 $F^k_{p,q}(\Om)$ 中稠密。
\end{theorem}

\begin{proof}
全文记 $T(x) = d(x)/(2(k+1))$，并对 $r = 0,\dots,k$ 记
\[
I_r(x) := \int_{|h|<T(x)} |\Delta_h^{k+1-r} f(x)|^q |h|^{-n-(k-r)q} d(x)^{-rq} dh,
\]
于是 $[f]_{r,k} = ( \int_\Om I_r(x)^{p/q} dx )^{1/p}$。由假设 $f \in F^k_{p,q}(\Om)$，故对每个 $r$ 都有 $I_r^{p/q} \in L^1(\Om)$，且 $f/d^k \in L^p(\Om)$。记 $D := \diam(\Om) \in (0,\infty)$。

\textbf{第一步（截断函数的构造）} 取 $\psi \in C_c^\infty(\R^n)$ 满足 $\psi \ge 0$，$\supp \psi$ 含于闭单位球内，$\int \psi = 1$。对 $\delta > 0$ 令 $\psi_\delta(z) := \delta^{-n} \psi(z/\delta)$，并令 $G_\delta(x) := 1$（当 $x \in \Om$ 且 $d(x) \ge 3 \delta$ 时），$G_\delta(x) := 0$（其它情形）。对 $x \in \Om$ 定义
\[
\chi_\varepsilon(x) := (G_{\varepsilon/2} * \psi_{\varepsilon/2})(x) = \int_{\R^n} \psi_{\varepsilon/2}(z) G_{\varepsilon/2}(x-z) dz.
\]
则 $\chi_\varepsilon$ 在 $\Om$ 上是 $C^\infty$ 的，且 $0 \le \chi_\varepsilon \le 1$。若 $d(x) \ge 2 \varepsilon$ 且 $|z| < \varepsilon/2$，则 $d(x-z) \ge d(x) - |z| > 2\varepsilon - \varepsilon/2 > 3\varepsilon/2$ 且 $x-z \in \Om$，故 $G_{\varepsilon/2}(x-z) = 1$，从而 $\chi_\varepsilon(x) = 1$。若 $d(x) < \varepsilon$ 且 $|z| < \varepsilon/2$，则 $d(x-z) \le d(x) + |z| < 3\varepsilon/2$，故 $G_{\varepsilon/2}(x-z) = 0$，从而 $\chi_\varepsilon(x) = 0$。最后，在 $\Om$ 上 $D^\alpha \chi_\varepsilon = G_{\varepsilon/2} * D^\alpha \psi_{\varepsilon/2}$，由此
\[
|D^\alpha \chi_\varepsilon| \le \|G_{\varepsilon/2}\|_\infty \|D^\alpha \psi_{\varepsilon/2}\|_{L^1} = 2^{|\alpha|} \|D^\alpha \psi\|_{L^1} \varepsilon^{-|\alpha|}.
\]
于是 $\chi_\varepsilon$ 具有所需的诸性质，其中 $C_{|\alpha|} := 2^{|\alpha|}\|D^\alpha \psi\|_{L^1}$。（下文中只用到上面列出的这些性质。）

\textbf{第二步（有限差分的 Leibniz 公式）} 对一切 $m \ge 0$、一切定义在一个（对相应平移封闭的）集合上的函数 $u, v$ 以及一切 $x, h$，有
\[
\Delta_h^m(uv)(x) = \sum_{a=0}^{m} \binom{m}{a} \Delta_h^a v(x+(m-a)h) \Delta_h^{m-a} u(x). \qquad (L)
\]
\emph{证明。} 对每个整数 $j \ge 0$ 有 $u(x+jh) = \tau_h^j u(x) = (I + \Delta_h)^j u(x) = \sum_{a=0}^{j} \binom{j}{a} \Delta_h^a u(x)$，其中 $\tau_h u = u(\cdot +h)$ 而 $I$ 是恒等算子（二项式定理）。于是
\[
\begin{aligned}
\Delta_h^m(uv)(x) &= \sum_{j=0}^{m} (-1)^{m-j} \binom{m}{j} u(x+jh) v(x+jh) \\
&= \sum_{j=0}^{m} (-1)^{m-j} \binom{m}{j} v(x+jh) \sum_{a=0}^{j} \binom{j}{a} \Delta_h^a u(x) \\
&= \sum_{a=0}^{m} \Delta_h^a u(x) \sum_{j=a}^{m} (-1)^{m-j} \binom{m}{j} \binom{j}{a} v(x+jh).
\end{aligned}
\]
由于 $\binom{m}{j} \binom{j}{a} = \binom{m}{a} \binom{m-a}{j-a}$，内层和等于 $\binom{m}{a} \sum_{b=0}^{m-a} (-1)^{m-a-b} \binom{m-a}{b} v(x+(a+b)h) = \binom{m}{a} \Delta_h^{m-a} v(x+ah)$。因此 $\Delta_h^m(uv)(x) = \sum_{a=0}^{m} \binom{m}{a} \Delta_h^a u(x) \Delta_h^{m-a}v(x+ah)$；把求和指标 $a$ 换成 $b := m-a$，并利用 $\binom{m}{m-b} = \binom{m}{b}$，即得取 $u = f$、$v = g_\varepsilon$ 的 (L)。这就证明了 (L)。

\textbf{第三步（截断函数差分的尺度与支集）} 设 $1 \le a \le m$，并设 $y \in \Om$ 使得点 $y, y+h, \dots, y+a h$ 全都落在 $\Om$ 中，且连接相邻两点的线段也都落在 $\Om$ 中。由于 $\chi_\varepsilon$ 在 $\Om$ 上是 $C^\infty$ 的，且当 $|\alpha| \le a$ 时 $|D^\alpha \chi_\varepsilon| \le C_{|\alpha|} \varepsilon^{-|\alpha|}$，对中值定理取迭代积分形式（在凸球 $B(y,a|h|)$ 上）可得
\[
|\Delta_h^a g_\varepsilon(y)| = |\Delta_h^a \chi_\varepsilon(y)| \le C'_a \varepsilon^{-a} |h|^a, \qquad (S1)
\]
其中 $C'_a$ 只依赖于 $a$ 与 $n$（把因子 $n^a$ 与常数 $C_{|\alpha|}$ 一并吸收进去）。此外，若 $\Delta_h^a g_\varepsilon(y) \ne 0$，则某个点 $y + \nu h$（$0 \le \nu \le a$）满足 $g_\varepsilon(y+\nu h) \ne 0$，从而 $d(y+\nu h) < 2 \varepsilon$（因为 $g_\varepsilon = 0$ 于 $\{d \ge 2 \varepsilon\}$ 上）。特别地，若再有 $y = x+(m-a)h$，其中 $x \in \Om$ 且 $|h| < T(x)$，则
\[
\begin{aligned}
d(x) &\le d(y+\nu h) + |x - (y+\nu h)| < 2 \varepsilon + (m-a+\nu)|h| \\
&\le 2 \varepsilon + m |h| < 2 \varepsilon + d(x)/2, \qquad (S2)
\end{aligned}
\]
故 $d(x) < 4 \varepsilon$。

\textbf{第四步（$L^p$ 项）} 由于 $|g_\varepsilon| \le 1$ 且 $g_\varepsilon = 0$ 于 $\{d \ge 2 \varepsilon\}$ 上，
\[
\int_\Om |f - f \chi_\varepsilon|^p dx = \int_\Om |f|^p |g_\varepsilon|^p dx \le \int_{d < 2 \varepsilon} |f|^p dx \le \int_{d<4\varepsilon} |f|^p dx,
\]
\[
\int_\Om |(f - f \chi_\varepsilon)/d^k|^p dx \le \int_{d < 2 \varepsilon} |f/d^k|^p dx \le \int_{d<4\varepsilon} |f/d^k|^p dx,
\]
且两者当 $\varepsilon \to 0$ 时都趋于 $0$，因为 $|f|^p, |f/d^k|^p \in L^1(\Om)$ 且 $|\{d < 4 \varepsilon\}| \to 0$。它们正是条款 (ii) 中那 $k+3$ 个量的前两个。

\textbf{第五步（半范数项）} 固定 $r \in \{0,\dots,k\}$，令 $m := k+1-r$（于是 $1 \le m \le k+1$），并简记 $w := f g_\varepsilon = f - f \chi_\varepsilon$。对 $x \in \Om$ 与 $|h| < T(x)$，点 $x, x+h, \dots, x+(k+1)h$ 全都落在 $\Om$ 中，这是因为 $|x+\nu h - x| = \nu|h| \le (k+1)T(x) = d(x)/2 < d(x)$，而开球 $B(x,d(x))$ 含于 $\Om$；因此 (L) 与 (S1) 在该差分模板内均可用。由 (L) 与初等不等式 $(\sum_{a=0}^{m} c_a)^q \le C(m,q) \sum_{a=0}^{m} c_a^q$（当 $q \ge 1$ 时取 $C(m,q) := (m+1)^{q-1}$ 即成立；把二项式系数吸收进 $C(m,q)$），
\[
|\Delta_h^m(w)(x)|^q \le C(m,q) \sum_{a=0}^{m} |\Delta_h^a g_\varepsilon(x+(m-a)h)|^q |\Delta_h^{m-a}f(x)|^q. \qquad (4.0)
\]
对 $x \in \Om$ 与 $0 \le a \le m$ 定义
\[
J_r(x) := \int_{|h|<T(x)} |\Delta_h^m w(x)|^q |h|^{-n-(k-r)q} d(x)^{-rq} dh,
\]
\[
A_a(x) := \int_{|h|<T(x)} |\Delta_h^a g_\varepsilon(x+(m-a)h)|^q |\Delta_h^{m-a}f(x)|^q |h|^{-n-(k-r)q} d(x)^{-rq} dh.
\]
用非负核对方程 (4.0) 积分便得 $J_r(x) \le C(m,q) \sum_{a=0}^{m} A_a(x)$。又，对 $p, q > 0$ 与 $m \ge 0$，存在常数 $C_1(m,p,q)$，使得对所有 $c_a \ge 0$ 都有 $(\sum_{a=0}^{m} c_a)^{p/q} \le C_1(m,p,q) \sum_{a=0}^{m} c_a^{p/q}$（当 $p \ge q$ 时用 $(\sum c_a)^{p/q} \le (m+1)^{p/q-1} \sum c_a^{p/q}$；当 $p < q$ 时用 $t \mapsto t^{p/q}$ 的凹性，它给出 $(\sum c_a)^{p/q} \le \sum c_a^{p/q}$）。由于 $[w]_{r,k}^{p}$ 的内层积分正是 $J_r(x)$，故
\[
[w]_{r,k}^p = \int_\Om J_r(x)^{p/q} dx \le C_2(m,p,q) \sum_{a=0}^{m} \int_\Om A_a(x)^{p/q} dx. \qquad (4.1)
\]
下面就 $a$ 的三个范围估计 $\int_\Om A_a(x)^{p/q} dx$。

(a) $a = 0$。此时 $\Delta_h^0 g_\varepsilon = g_\varepsilon$，而 $\Delta_h^m f$ 是 $m$ 阶差分。由 (S2) 取 $a = 0$（这时条件只是 $g_\varepsilon(x+mh) \ne 0$，即 $d(x+mh) < 2 \varepsilon$），被积函数在 $d(x) \ge 4 \varepsilon$ 处为零，且 $|g_\varepsilon| \le 1$。故 $A_0(x) \le \mathbf{1}_{d(x)<4\varepsilon} I_r(x)$，从而
\[
\int_\Om A_0(x)^{p/q} dx \le \int_{d<4\varepsilon} I_r(x)^{p/q} dx \to 0,
\]
因为 $I_r^{p/q} \in L^1(\Om)$ 且 $|\{d<4\varepsilon\}| \to 0$。

(b) $1 \le a \le m-1$。令 $r' := r+a$，则 $r' \in \{r+1,\dots,k\}$ 且 $m-a = k+1-r'$。由 (S1)，$|\Delta_h^a g_\varepsilon(x+(m-a)h)|^q \le C_a^q \varepsilon^{-aq} |h|^{aq}$。于是在集合 $d(x) < 4 \varepsilon$ 上（由 (S2)，在其外被积函数为零），
\[
\begin{aligned}
A_a(x) &\le C_a^q \varepsilon^{-aq} \int_{|h|<T(x)} |h|^{aq} |\Delta_h^{m-a}f(x)|^q |h|^{-n-(k-r)q} d(x)^{-rq} dh \\
&= C_a^q \varepsilon^{-aq} d(x)^{aq} I_{r'}(x),
\end{aligned}
\]
这里用到 $aq - n - (k-r)q = -n - (k-r-a)q = -n-(k-r')q$ 以及 $d(x)^{-rq} = d(x)^{aq} d(x)^{-r'q}$（等价地 $d(x)^{-r'q} d(x)^{aq} = d(x)^{-rq}$，因为 $r' = r+a$）。在 $\{d(x) < 4 \varepsilon\}$ 上有 $\varepsilon^{-aq} d(x)^{aq} = (d(x)/\varepsilon)^{aq} < 4^{aq}$。因此
\[
A_a(x) \le C_a^q 4^{aq} \mathbf{1}_{d(x)<4\varepsilon} I_{r'}(x),
\]
\[
\int_\Om A_a(x)^{p/q} dx \le C_a^p 4^{ap} \int_{d<4\varepsilon} I_{r'}(x)^{p/q} dx \to 0,
\]
因为 $I_{r'}^{p/q} \in L^1(\Om)$（这里 $r' \le k$；$r'=k$ 那一项正是 $k$ 阶源半范数，由 $f \in F^k$ 而有限）。

(c) $a = m$。此时 $|\Delta_h^m g_\varepsilon| = |\Delta_h^m(1 - \chi_\varepsilon)| = |\Delta_h^m \chi_\varepsilon|$（因为常数 $1$ 的 $m$ 阶差分为零，$m \ge 1$），$\Delta_h^0 f = f$，且由 (S1) $|\Delta_h^m g_\varepsilon(x)|^q \le C_m^q \varepsilon^{-mq} |h|^{mq}$。于是在 $d(x) < 4 \varepsilon$ 上（由 (S2)），
\[
A_m(x) \le C_m^q \varepsilon^{-mq} |f(x)|^q d(x)^{-rq} \int_{|h|<T(x)} |h|^{mq} |h|^{-n-(k-r)q} dh.
\]
现在 $mq - n - (k-r)q = q - n$，故径向积分为 $\int_{|h|<T(x)} |h|^{q-n} dh = \omega_{n-1} T(x)^q / q$（指数 $q-1 > -1$ 使它在 $h = 0$ 处可积），其中 $T(x) = d(x)/(2(k+1))$。因此
\[
A_m(x) \le C \varepsilon^{-mq} |f(x)|^q d(x)^{-rq} d(x)^q = C \varepsilon^{-mq} |f(x)|^q d(x)^{(1-r)q},
\]
再取 $p/q$ 次幂，
\[
\begin{aligned}
A_m(x)^{p/q} &\le C \varepsilon^{-mp} |f(x)|^p d(x)^{(1-r)p} = C \varepsilon^{-mp} (|f(x)|/d(x)^k)^p d(x)^{(k+1-r)p} \\
&= C \varepsilon^{-mp} (|f(x)|/d(x)^k)^p d(x)^{mp},
\end{aligned}
\]
这里用到 $|f|^p d^{(1-r)p} = (|f|/d^k)^p d^{kp} d^{(1-r)p} = (|f|/d^k)^p d^{(k+1-r)p}$ 以及 $m = k+1-r$。在 $\{d(x) < 4 \varepsilon\}$ 上有 $\varepsilon^{-mp} d(x)^{mp} = (d(x)/\varepsilon)^{mp} < 4^{mp}$；因此
\[
\int_\Om A_m(x)^{p/q} dx \le C 4^{mp} \int_{d<4\varepsilon} |f/d^k|^p dx \to 0.
\]
把 (a)--(c) 与 (4.1) 结合，得到对每个 $r = 0,\dots,k$ 都有 $[f - f \chi_\varepsilon]_{r,k} \to 0$，并且其 $p$ 次幂 $[f-f \chi_\varepsilon]_{r,k}^p$ 被 $C(n,k,p,q)$ 乘以相应的 $\{d<4\varepsilon\}$ 上的积分之和所界。再与第四步结合，即得 $\|f - f \chi_\varepsilon\|_{F^k_{p,q}} \to 0$ 以及条款 (ii)。特别地，对每个 $\varepsilon$ 都有 $\|w\|_{F^k_{p,q}} < \infty$，故 $w = f - f \chi_\varepsilon \in F^k_{p,q}$；又因 $f \in F^k_{p,q}$ 且 $F^k_{p,q}$ 是线性空间，故 $f \chi_\varepsilon = f - w \in F^k_{p,q}$。这就证明了 (i)。

\textbf{第六步（紧支与稠密性）} 由于 $\chi_\varepsilon = 0$ 于 $\{d < \varepsilon\}$ 上，函数 $f \chi_\varepsilon$ 本性支撑于 $\{x \in \Om : d(x) \ge \varepsilon\}$，而后者是 $\R^n$ 的闭子集、含于有界集 $\Om$ 中，故为紧集且含于 $\Om$。因此 $f \chi_\varepsilon \in F^k_{p,q,c}(\Om)$，且 $\|f - f \chi_\varepsilon\|_{F^k_{p,q}} \to 0$。这就证明了 $F^k_{p,q,c}(\Om)$ 在 $F^k_{p,q}(\Om)$ 中稠密。
\end{proof}

\begin{remark}[直觉]
这是整数阶修正空间（CLWZ 定义 1.6）稠密性问题中的截断那一半。截断误差 $w = f(1-\chi_\varepsilon)$ 支撑在边界层 $\{d<2\varepsilon\}$ 中。在 $\Delta_h^{k+1-r}(\chi_\varepsilon f)$ 的离散 Leibniz 公式中，每一项要么带有 $f$ 的差分在基点 $x$ 处的全阶消失性（$a = 0$ 的项与 $1 \le a \le m-1$ 的项），要么带有截断函数额外的消失性 $|\Delta_h^m \chi_\varepsilon| \le C(|h|/\varepsilon)^m$（$a = m$ 的项）。真正起作用的只是两个记账恒等式。在中间的项里，权重恒等式 $d(x)^{-rq} = d(x)^{aq} d(x)^{-r'q}$（其中 $r' = r+a > r$）把该项变成 $r'$ 阶的源半范数，而因子 $(d(x)/\varepsilon)^{aq} \le 4^{aq}$ 是无害的，因为该层的宽度与 $\varepsilon$ 相若。在最顶端的项里，指数恒等式 $mq - n - (k-r)q = q-n$ 使径向积分成为 $\int_{|h|<T}|h|^{q-n}dh = \omega_{n-1}T^q/q$（于是贡献一个因子 $d(x)^q$），它与 $d(x)^{-rq}$ 结合给出 $d(x)^{(1-r)q}$，再在乘除 $d(x)^{kp}$ 之后给出因子 $(|f|/d^k)^p d(x)^{mp} \le 4^{mp}(|f|/d^k)^p$。关键的手段是 Leibniz 公式的那个版本 (L)：其中 $f$ 的差分取在\defterm{基点} $x$，而截断函数的差分取在平移点 $x+(m-a)h$；这就把权重 $d(x)^{-rq}$ 与源半范数指标 $r'$ 都附着在同一个点 $x$ 上，因而无需任何环域（annulus）或移位基点的比较。（另一个版本把 $f$ 的差分取在 $x+ah$，恰恰会产生此前诸次尝试都无法克服的那种环域困难。）注意这一结果与距离幂 $d^k$ 自身的失效（当 $k \ge 2$ 时 $d^k \notin F^k$，见附录 A）并不矛盾：截断必须使用光滑的截断函数与全阶消去，而不能使用边界距离恒等式。
\end{remark}

%==================================================================
\subsection{修正后的高阶加权空间的非退化性}

\begin{theorem}[修正后的高阶加权空间非退化]
设 $n\ge1$，$\Om$ 为 $\R^n$ 中的非空有界域，对 $x\in\Om$ 令 $d(x)=\dist(x,\partial\Om)$，设 $\sigma\in(0,1)$，$q\in[1,\infty)$，$k\ge1$ 为整数，并令 $s:=k+\sigma$。设 $a>0$ 使 $\int_\Om d(y)^{-aq}\,dy<\infty$（情形 (I)；例如对每个满足 $aq<1$ 的有界 Lipschitz 域成立）。对 $\Om$ 上的可测函数 $g$ 令
\[
  [g]_{\sigma,q,a}:=\Big(\int_\Om\int_\Om |g(x)-g(y)|^q\,|x-y|^{-n-\sigma q}\,d(y)^{-aq}\,dy\,dx\Big)^{1/q},
\]
并定义 $X$ 为这样的 Lebesgue 可测函数 $f$ 之集：其弱导数 $D^\alpha f$（$|\alpha|\le k$）存在，且
\[
  \|f\|_X:=\sum_{|\alpha|\le k}\|D^\alpha f\|_{L^q(\Om)}
  +\sum_{|\beta|=k}[D^\beta f]_{\sigma,q,a}
  +\sum_{|\alpha|\le k}[D^\alpha f]_{\sigma,q,a}
\]
有限（这就是内蕴的 CLWZ (1.2) 范数，只是额外要求每个导数——而不只是最高阶导数——都属于加权分数阶空间）。则每个 $f\in C_c^\infty(\Om)$ 都属于 $X$，故 $X$ 含有非常数函数；特别地，$X$ 既不是 $\{0\}$，也不是常数构成的一维空间 $\R\cdot\mathbf{1}$。这与 CLWZ 定义 1.3 的修改内蕴空间形成对照：只要 $\int_\Om d^{-kq}\,dy=\infty$，后者就等于 $\R\cdot\mathbf{1}$（在加入 Hardy 项后则为 $\{0\}$）。
\end{theorem}

\begin{proof}
设 $f\in C_c^\infty(\Om)$；则每个满足 $|\alpha|\le k$ 的 $D^\alpha f$ 都是支集含于 $K:=\supp f$ 的有界光滑函数，故对一切 $|\alpha|\le k$ 有 $\|D^\alpha f\|_{L^q(\Om)}<\infty$。对每个 $|\alpha|\le k$ 令
\[
  B_\alpha:=\sup\Big\{\frac{|D^\alpha f(u)-D^\alpha f(v)|}{\min(|u-v|,2)} : u,v\in\Om,\ u\ne v\Big\}<\infty ;
\]
它有限，因为 $D^\alpha f$ 是梯度有界的有界光滑函数，且当 $|u-v|\le1$ 时 $\min(|u-v|,2)\le|u-v|$，而当 $|u-v|\ge1$ 时 $\min(|u-v|,2)\ge1$ 且 $D^\alpha f$ 有界。于是对一切 $u,v\in\Om$，
\[
  |D^\alpha f(u)-D^\alpha f(v)|\le B_\alpha\min(|u-v|,2).
\]
因此对每个 $|\alpha|\le k$，
\begin{align*}
  [D^\alpha f]^q_{\sigma,q,a}
  &=\int_\Om\int_\Om |D^\alpha f(x)-D^\alpha f(y)|^q\,|x-y|^{-n-\sigma q}\,d(y)^{-aq}\,dy\,dx\\
  &\le B_\alpha^q\int_\Om d(y)^{-aq}\Big(\int_\Om \min(|x-y|,2)^q\,|x-y|^{-n-\sigma q}\,dx\Big)dy\\
  &\le B_\alpha^q\Big(\int_\Om d(y)^{-aq}\,dy\Big)\Big(\int_{\R^n}\min(|z|,2)^q\,|z|^{-n-\sigma q}\,dz\Big),
\end{align*}
其中最后一步在作代换 $z=x-y$ 后，将内层的 $x$-积分（在 $\Om$ 上）延拓为 $\R^n$ 上的积分。令
\[
  Z:=\int_{\R^n}\min(|z|,2)^q\,|z|^{-n-\sigma q}\,dz .
\]
在 $|z|=2$ 处分割：当 $|z|\le2$ 时 $\min(|z|,2)^q=|z|^q$，且 $|z|^{q-n-\sigma q}$ 在 $0$ 附近可积，因为 $q-n-\sigma q=-n+(1-\sigma)q>-n$（由于 $\sigma<1$）；当 $|z|\ge2$ 时 $\min(|z|,2)^q=2^q$，且 $2^q|z|^{-n-\sigma q}$ 在无穷远处可积，因为 $\sigma q>0$。故 $Z<\infty$。由假设 $\int_\Om d(y)^{-aq}\,dy<\infty$，故对每个 $|\alpha|\le k$ 有 $[D^\alpha f]_{\sigma,q,a}<\infty$，特别对每个 $|\beta|=k$ 成立。于是 $\|f\|_X<\infty$ 且 $f\in X$。

空间 $X$ 含有每个 $f\in C_c^\infty(\Om)$；取任一非零的此类 $f$（例如支集含于 $\Om$ 内某球中的非零光滑鼓包函数）即得 $X$ 中的非常数元素。故 $X$ 既不是 $\{0\}$，也不是常数空间。
\end{proof}

\begin{remark}[直觉]
CLWZ 定义 1.3 的退化来自\emph{加大}的权重 $d^{-(k-|\beta|)q}$。自然的修正改为：每个导数都保留原来的可积权重 $d^{-aq}$，但对每个导数（而不只是最高阶导数）都取加权分数阶半范数。在原来的可积权重下，紧支光滑函数属于该空间：它们的导数以 $\min(|x-y|,2)$ 为 Lipschitz 型界，于是半范数分解为 $d^{-aq}$ 的总质量与一个有限的径向积分之积。
\end{remark}

%=====================================================================
\section{情形 (I)：可积权下的稠密性}\label{sec:part3}
%=====================================================================
% ===================================================================
%  第 3 组：情形 (I) —— 可积权下的稠密性
% ===================================================================

本节处理权重可积的情形 (I)，即 $\nuA(\Om)=\int_\Om d(y)^{-aq}\,dy<\infty$。此时加权内在 Triebel--Lizorkin 空间 $\Fsp$ 含有足够多的函数，稠密性不再自动成立，必须借助内点半范数估计、磨光与 Whitney 局部化来证明。第一小节给出内点加权威半范数估计及可积权下的磨光收敛这一纯技术工具，其后三条事实依次建立 $p=q$、$s\in(0,1)$ 时的 $C^\infty$ 稠密性、Hardy 加强空间中的稠密性，以及区间上 $s>1$ 的高阶稠密性；与第 2 组情形 (II) 的退化现象相对照，这些结论说明可积权下光滑函数确实可以在完整范数意义下逼近空间中每个元素。

%==================================================================
\subsection{内点加权威半范数估计与可积权下的磨光收敛}

\begin{lemma}[可积情形下的内点加权威半范数估计与磨光收敛]
设 $n\ge 1$，$\Om$ 为 $\R^n$ 中的非空有界域，对 $x\in\Om$ 记 $d(x)=\dist(x,\partial\Om)$，设 $\sigma\in(0,1)$，$q\in[1,\infty)$。设 $w\colon\Om\to[0,\infty)$ 为 Borel 函数，满足
\[
  \|w\|_{L^1(\Om)}=\int_\Om w(y)\,dy<\infty .
\]
固定 $\delta>0$，并假设
\[
  W_\delta:=\esssup\,\{w(y) : y\in\Om,\ d(y)\ge\delta/2\}<\infty .
\]
设 $h\in L^q(\Om)$ 为 Lebesgue 可测函数，在边界层 $\{x\in\Om : d(x)<\delta\}$ 上 $h=0$ a.e.，且其无权 Gagliardo 半范数有限
\[
  [h]_{\sigma,q}:=\Big(\int_\Om\int_\Om |h(x)-h(y)|^q\,|x-y|^{-n-\sigma q}\,dy\,dx\Big)^{1/q}<\infty .
\]
则
\begin{equation}
  \int_\Om\int_\Om |h(x)-h(y)|^q\,|x-y|^{-n-\sigma q}\,w(y)\,dy\,dx
  \le W_\delta\,[h]^q_{\sigma,q}
  +2^{n+\sigma q}\delta^{-n-\sigma q}\,\|w\|_{L^1(\Om)}\,\|h\|^q_{L^q(\Om)} .
  \tag{3.1}
\end{equation}
此外，再假设 $w$ 在 $\Om$ 上连续。设 $\rho$ 为支于 $\R^n$ 闭单位球内的非负光滑函数，$\int\rho=1$，对 $\varepsilon>0$ 令 $\rho_\varepsilon(x)=\varepsilon^{-n}\rho(x/\varepsilon)$，设 $k$ 为 $h$ 到 $\R^n$ 的零延拓，并令 $h_\varepsilon:=(k*\rho_\varepsilon)|_\Om$。则对每个满足 $0<\varepsilon<\delta/2$ 的 $\varepsilon$，都有 $h_\varepsilon\in C^\infty(\Om)$、$h_\varepsilon=0$ 于 $\{x\in\Om : d(x)<\delta/2\}$ 上，且加权威半范数
\[
  \int_\Om\int_\Om |(h_\varepsilon-h)(x)-(h_\varepsilon-h)(y)|^q\,|x-y|^{-n-\sigma q}\,w(y)\,dy\,dx
\]
当 $\varepsilon\to 0$ 时趋于 $0$。
\end{lemma}

\begin{proof}
\emph{第一部分（估计）}. 由于在边界层 $\{z : d(z)<\delta\}$ 上 $h=0$ a.e.，只要 $d(y)<\delta$ 且 $d(x)<\delta$，被积函数 $|h(x)-h(y)|^q|x-y|^{-n-\sigma q}w(y)$ 即为零。将 $y$ 上的积分拆成两个集合 $\{d(y)<\delta/2\}$ 与 $\{d(y)\ge\delta/2\}$。

\emph{(a) 集合 $\{d(y)<\delta/2\}$}. 若 $d(x)<\delta$，则 $h(x)=h(y)=0$ a.e.，故仅有满足 $d(x)\ge\delta$ 的 $x$ 有贡献；对这种 $x$ 以及满足 $d(y)<\delta/2$ 的 $y$，距离函数的反向三角不等式给出
\[
  |x-y|\ge d(x)-d(y)>\delta-\delta/2=\delta/2 ,
\]
从而 $|x-y|^{-n-\sigma q}\le(\delta/2)^{-n-\sigma q}=2^{n+\sigma q}\delta^{-n-\sigma q}$。因此
\begin{equation}
\begin{split}
  &\int_\Om\int_{d(y)<\delta/2} |h(x)-h(y)|^q\,|x-y|^{-n-\sigma q}\,w(y)\,dy\,dx \\
  &\qquad\le 2^{n+\sigma q}\delta^{-n-\sigma q}
  \int_{d(x)\ge\delta}|h(x)|^q\,dx\cdot\int_{d(y)<\delta/2}w(y)\,dy \\
  &\qquad\le 2^{n+\sigma q}\delta^{-n-\sigma q}\,\|w\|_{L^1(\Om)}\,\|h\|^q_{L^q(\Om)} .
\end{split}
  \tag{3.2}
\end{equation}

\emph{(b) 集合 $\{d(y)\ge\delta/2\}$}. 在该集合上按定义 $w(y)\le W_\delta$，故
\begin{equation}
\begin{split}
  \int_\Om\int_{d(y)\ge\delta/2} |h(x)-h(y)|^q\,|x-y|^{-n-\sigma q}\,w(y)\,dy\,dx
  &\le W_\delta\int_\Om\int_\Om |h(x)-h(y)|^q\,|x-y|^{-n-\sigma q}\,dy\,dx \\
  &= W_\delta\,[h]^q_{\sigma,q} .
\end{split}
  \tag{3.3}
\end{equation}
将 (3.2) 与 (3.3) 相加即得所写出的估计 (3.1)。

\emph{第二部分（磨光）}. 再假设 $w$ 在 $\Om$ 上连续；此时 $w$ 在 $\Om$ 的任一紧子集上的上确界均有限，特别地
\[
  W':=\sup\{w(y) : d(y)\ge\delta/4\}<\infty ,
\]
因为 $\{d\ge\delta/4\}$ 是 $\Om$ 的紧子集。设 $k$ 为 $h$ 的零延拓。则 $k\in L^q(\R^n)$，且差一个零测集，$\supp k$ 含于 $\{z\in\Om : d(z)\ge\delta\}$ 之中，这是因为 $h$ 在 $\{d<\delta\}$ 上 a.e. 为零。

\emph{断言 1（$k$ 在 $\R^n$ 上的无权 Gagliardo 半范数有限）}. 事实上有
\begin{equation}
\begin{split}
  [k]^q_{\sigma,q,\R^n}
  &=\int_{\R^n}\int_{\R^n}|k(x)-k(y)|^q\,|x-y|^{-n-\sigma q}\,dy\,dx \\
  &=[h]^q_{\sigma,q,\Om}
  +2\int_{x\in\supp k}\int_{y\notin\Om}|h(x)|^q\,|x-y|^{-n-\sigma q}\,dy\,dx ,
\end{split}
  \tag{3.4}
\end{equation}
而对 $x\in\supp k$ 与 $y\notin\Om$ 有 $|x-y|\ge d(x)\ge\delta$；因此
\[
  \int_{y\notin\Om}|x-y|^{-n-\sigma q}\,dy
  \le\int_{|x-y|\ge\delta}|x-y|^{-n-\sigma q}\,dy
  =\frac{\omega_{n-1}\delta^{-\sigma q}}{\sigma q}<\infty
\]
（径向积分：$\omega_{n-1}\int_\delta^\infty r^{-n-\sigma q}r^{n-1}dr=\omega_{n-1}\int_\delta^\infty r^{-1-\sigma q}dr=\omega_{n-1}\delta^{-\sigma q}/(\sigma q)$，其中 $\omega_{n-1}$ 是 $\R^n$ 中单位球面 $S^{n-1}$ 的 $(n-1)$ 维测度）。故 (3.4) 中的交叉项至多为
\[
  \frac{2\omega_{n-1}\delta^{-\sigma q}}{\sigma q}\,\|h\|^q_{L^q(\Om)}<\infty ,
\]
从而
\[
  [k]^q_{\sigma,q,\R^n}\le[h]^q_{\sigma,q}
  +\frac{2\omega_{n-1}\delta^{-\sigma q}}{\sigma q}\,\|h\|^q_{L^q(\Om)}<\infty .
\]

\emph{断言 2（$h_\varepsilon$ 光滑且在边界附近消失）}. 函数 $h_\varepsilon:=(k*\rho_\varepsilon)|_\Om$ 在 $\Om$ 上无穷次可微；且若 $0<\varepsilon<\delta/2$，则 $h_\varepsilon$ 在 $\{x\in\Om : d(x)<\delta/2\}$ 上为零。事实上，若 $x\in\Om$ 且 $h_\varepsilon(x)\neq0$，则存在 $z\in\supp k$ 使 $|x-z|\le\varepsilon$，于是
\[
  d(x)\ge d(z)-|x-z|\ge\delta-\varepsilon>\delta/2 .
\]

\emph{断言 3（在 $L^q$ 与无权半范数中的收敛）}. 当 $\varepsilon\to0$ 时 $\|h_\varepsilon-h\|_{L^q(\Om)}\to0$ 且 $[h_\varepsilon-h]_{\sigma,q,\Om}\to0$。对 $L^q$ 部分，
\[
  \|h_\varepsilon-h\|_{L^q(\Om)}\le\|k*\rho_\varepsilon-k\|_{L^q(\R^n)}\to0 ,
\]
这是磨光的标准 $L^q$ 连续性：由 Minkowski 不等式
\[
  \|k*\rho_\varepsilon-k\|_{L^q}\le\int\rho(z)\,\|\tau_z k-k\|_{L^q}\,dz\to0
\]
由控制收敛可得，其中用到 $\|\tau_z k-k\|_{L^q}\le2\|k\|_{L^q}$ 以及平移在 $L^q$ 中的连续性。对半范数部分使用平移公式：对 $g\in L^q(\R^n)$，
\begin{equation}
  [g]^q_{\sigma,q,\R^n}=\int_{\R^n}\|\tau_z g-g\|^q_{L^q(\R^n)}\,|z|^{-n-\sigma q}\,dz ,
  \tag{3.5}
\end{equation}
这由代换 $y=x-z$ 得到。由于与 $\rho_\varepsilon$ 的卷积与平移交换，
\[
  \tau_z(k*\rho_\varepsilon)-(k*\rho_\varepsilon)=(\tau_z k-k)*\rho_\varepsilon ,
\]
故
\[
  \tau_z(k*\rho_\varepsilon)-(k*\rho_\varepsilon)-(\tau_z k-k)=(\tau_z k-k)*\rho_\varepsilon-(\tau_z k-k),
\]
从而由 (3.5)
\begin{equation}
  [k*\rho_\varepsilon-k]^q_{\sigma,q,\R^n}
  =\int_{\R^n}\big\|(\tau_z k-k)*\rho_\varepsilon-(\tau_z k-k)\big\|^q_{L^q(\R^n)}\,|z|^{-n-\sigma q}\,dz .
  \tag{3.6}
\end{equation}
对每个固定的 $z$，函数 $\tau_z k-k$ 属于 $L^q(\R^n)$ 且
\[
  \big\|(\tau_z k-k)*\rho_\varepsilon-(\tau_z k-k)\big\|_{L^q}\to0
  \qquad(\varepsilon\to0)
\]
（磨光在 $L^q$ 中收敛）；此外该范数至多为 $2\|\tau_z k-k\|_{L^q}$，因为由 $\int\rho=1$ 与 $\rho\ge0$（Jensen 不等式）有 $\|g*\rho_\varepsilon\|_{L^q}\le\|g\|_{L^q}$。控制函数
\[
  2^q\,\|\tau_z k-k\|^q_{L^q}\,|z|^{-n-\sigma q}
\]
由断言 1 可积，其积分为 $2^q[k]^q_{\sigma,q,\R^n}<\infty$。故由控制收敛，$[k*\rho_\varepsilon-k]^q_{\sigma,q,\R^n}\to0$；又 $h_\varepsilon-h$ 是 $k*\rho_\varepsilon-k$ 在 $\Om$ 上的限制，所以
\[
  [h_\varepsilon-h]^q_{\sigma,q,\Om}\le[k*\rho_\varepsilon-k]^q_{\sigma,q,\R^n}\to0 .
\]

\emph{结论}. 最后，由断言 2 以及 $h$ 在 $\{d<\delta\}$ 上为零可知，$h_\varepsilon-h$ 在 $\{d<\delta/2\}$ 上为零。于是把第一部分应用于函数 $h_\varepsilon-h$，将 $\delta$ 换为 $\delta/2$，并注意权 $w$ 在更大的集合 $\{d\ge\delta/4\}$ 上被 $W'<\infty$ 控制，即得
\begin{equation}
\begin{split}
  &\int_\Om\int_\Om |(h_\varepsilon-h)(x)-(h_\varepsilon-h)(y)|^q\,|x-y|^{-n-\sigma q}\,w(y)\,dy\,dx \\
  &\qquad\le W'\,[h_\varepsilon-h]^q_{\sigma,q}
  +2^{n+\sigma q}(\delta/2)^{-n-\sigma q}\,\|w\|_{L^1(\Om)}\,\|h_\varepsilon-h\|^q_{L^q(\Om)}\longrightarrow0
\end{split}
  \tag{3.7}
\end{equation}
当 $\varepsilon\to0$ 时成立，此处用到断言 3。这正是所断言的、$h_\varepsilon-h$ 的加权威半范数收敛到 $0$。
\end{proof}

\begin{remark}[直觉]
在 $h$ 的支撑集上，到边界的距离至少是 $\delta$，因此权重的奇性部分与支撑集相互分离：在那里核函数是可积的，具有径向界 $\omega_{n-1}\delta^{-\sigma q}/(\sigma q)$，而只有 $w$ 的总质量进入估计。在权重有界的地方，加权威半范数被无权半范数控制。对一个在边界距离 $\ge\delta$ 处被支撑的函数作磨光，其磨光仍保持在距离 $\ge\delta/2$ 之外，并在无权 Gagliardo 半范数下收敛，因而也在加权威半范数下收敛。
\end{remark}

%==================================================================
\subsection{情形 (I) 中 \texorpdfstring{$C^\infty$}{C-infinity} 在 \texorpdfstring{$p=q$}{p=q}、\texorpdfstring{$s\in(0,1)$}{s in (0,1)} 时的稠密性}

\begin{theorem}[情形 (I)：$p=q$、$s\in(0,1)$ 时 $C^\infty$ 的稠密性]
设 $n\ge1$，$\Om$ 为 $\R^n$ 中的非空有界域，$d(x)=\dist(x,\partial\Om)$，$a>0$，$s\in(0,1)$，$p\in[1,\infty)$，并假设
\[
  \int_\Om d(y)^{-ap}\,dy<\infty
\]
（附录 A 的情形 (I)）。设 $F^s_{p,p}(\Om,d^{-aq})$ 为加权内蕴空间，其范数为
\[
  \|f\|_{F^s}:=\|f\|_{L^p(\Om)}
  +\Big(\int_\Om\int_\Om |f(x)-f(y)|^p\,|x-y|^{-n-sp}\,d(y)^{-ap}\,dy\,dx\Big)^{1/p}.
\]
则 $C^\infty(\Om)\cap F^s_{p,p}(\Om,d^{-aq})$ 在 $F^s_{p,p}(\Om,d^{-aq})$ 中稠密，即 $C^\infty(\Om)\cap F^s_{p,p}(\Om,d^{-aq})$ 在 $F^s_{p,p}(\Om,d^{-aq})$ 中的闭包 $F^{s,2}_{p,p}(\Om,d^{-aq})$ 等于 $F^s_{p,p}(\Om,d^{-aq})$。
\end{theorem}

\begin{proof}
设 $\Om$ 为 $\R^n$ 中的非空有界域，$d(x)=\dist(x,\partial\Om)$，$a>0$，$s\in(0,1)$，$p\in[1,\infty)$，并假设
\[
  \nu(\Om):=\int_\Om d(y)^{-ap}\,dy<\infty .
\]
令 $v(y):=d(y)^{-ap}$，
\[
  [f]:=\Big(\int_\Om\int_\Om |f(x)-f(y)|^p\,|x-y|^{-n-sp}\,v(y)\,dy\,dx\Big)^{1/p};
\]
则 $F^s_{p,p}(\Om,d^{-aq})=\{f\in L^p : \|f\|_{L^p}+[f]<\infty\}$。

我们使用 Kijaczko 的结果（\texttt{arXiv:2110.04012}，定理 19），其陈述为：对 $\R^d$ 中的非空开集 $\Om$，$s\in(0,1)$，$p\in[1,\infty)$，以及非负、局部有界且满足
\[
  \int_\Om w(x)(1+|x|)^{-d-sp}\,dx<\infty
  \quad\text{与}\quad
  \int_\Om v(x)(1+|x|)^{-d-sp}\,dx<\infty ,
\]
的权 $w,v$，记
\[
  \wt W^{s,p;w,v}(\Om):=\Big\{f\ \text{measurable} :
  \int_\Om\int_\Om |f(x)-f(y)|^p|x-y|^{-d-sp}w(x)v(y)\,dx\,dy<\infty\Big\}.
\]
则集合 $C^\infty(\Om)\cap\wt W^{s,p;w,v}(\Om)$ 在该半范空间中稠密。其证明通过 Whitney 型局部磨光
\[
  P^\eta f=\sum_n (f\psi_n)*h_{\eta(Q_n)} ,
\]
得到稠密性，并且尺度可取为对所有 $n$ 满足 $\eta(Q_n)<1/k$。

我们还使用 Dyda--Kijaczko 的结果（\texttt{arXiv:2009.00077}，定理 3）：对 $f\in L^p(\Om)$，只要对每个 $n$ 有 $\eta(Q_n)\to0$，$P^\eta f$ 就在 $L^p(\Om)$ 中收敛到 $f$。

\emph{第一步（验证外部定理的假设）}. 取 $w:=1$、$v:=d^{-ap}$，两个权都在 $\Om$ 上连续且为正，因而局部有界。由于 $\Om$ 有界，$(1+|x|)^{-n-sp}$ 在 $\Om$ 上被正常数从上、下界控制，故定理 19 的两个条件归结为 $w\in L^1(\Om)$ 与 $v\in L^1(\Om)$；前者因为 $|\Om|<\infty$ 成立，后者即 $\nu(\Om)<\infty$。于是定理 19 适用，且我们的 $F^s$ 恰为 $L^p(\Om)\cap\wt W$，其范数为 $\|f\|_{L^p}+[f]$。

\emph{第二步（选择尺度）}. 设 $f\in F^s$。对每个 $k\in\N$ 选取尺度 $\eta_k$，使得对所有 $n$ 有 $\eta_k(Q_n)<1/k$ 且 $[P^{\eta_k}f-f]<1/k$（定理 19）。则对每个固定的 $n$ 有 $\eta_k(Q_n)\to0$，故由 Dyda--Kijaczko 定理 3 得 $\|P^{\eta_k}f-f\|_{L^p(\Om)}\to0$。因此
\[
  \|P^{\eta_k}f-f\|_{F^s}
  =\|P^{\eta_k}f-f\|_{L^p}+[P^{\eta_k}f-f]\longrightarrow0 .
\]

\emph{第三步（逼近序列属于交集）}. 每个 $P^{\eta_k}f$ 既属于 $C^\infty(\Om)$ 又属于 $\wt W$（定理 19）。它是局部 $L^p$ 平均的有限重叠和，故
\[
  \|P^{\eta_k}f\|_{L^p}\le C\,\|f\|_{L^p}<\infty ;
\]
从而 $P^{\eta_k}f\in C^\infty(\Om)\cap F^s$。

\emph{第四步（结论）}. 于是 $C^\infty(\Om)\cap F^s$ 在 $F^s$ 中稠密，即
\[
  F^{s,2}_{p,p}(\Om,d^{-aq})=F^s_{p,p}(\Om,d^{-aq}) .
\]
\end{proof}

\begin{remark}[直觉]
Kijaczko 的单侧稠密性定理在 $w=1$、$v=d^{-ap}$ 可积时适用，而它给出的 Whitney 磨光逼近序列还在 $L^p$ 中收敛（Dyda--Kijaczko 定理 3），于是得到完整范数意义下的稠密性。
\end{remark}

%==================================================================
\subsection{情形 (I) 中 Hardy 加强空间内的稠密性}

\begin{theorem}[情形 (I)：$p=q$、$s\in(0,1)$ 时 Hardy 加强空间中的稠密性]
设 $n\ge1$，$\Om$ 为 $\R^n$ 中的非空有界域，$d(x)=\dist(x,\partial\Om)$，$a>0$，$s\in(0,1)$，$p\in[1,\infty)$，且 $\int_\Om d(y)^{-ap}\,dy<\infty$（情形 (I)）。取附录 A 中的 $F^s_{p,p}(\Om,d^{-aq})$，令
\[
  \wt F^s_{p,p}(\Om,d^{-aq})
  :=\Big\{f\in F^s_{p,p}(\Om,d^{-aq}) :
  \|f d^{-(s+a)}\|_{L^p}+\|f d^{-(s+a+n/p)}\|_{L^p}<\infty\Big\}
\]
（CLWZ (1.6) 的 Hardy 加强空间，$q=p$）。则 $C^\infty(\Om)\cap\wt F^s_{p,p}(\Om,d^{-aq})$ 在 $\wt F^s_{p,p}(\Om,d^{-aq})$ 中按如下范数稠密：
\[
  \|f\|_{\wt F^s}:=\|f\|_{F^s}+\|f d^{-(s+a)}\|_{L^p}+\|f d^{-(s+a+n/p)}\|_{L^p} ;
\]
等价地，$C^\infty(\Om)\cap\wt F^s_{p,p}(\Om,d^{-aq})$ 在 $\wt F^s_{p,p}(\Om,d^{-aq})$ 中的闭包等于 $\wt F^s_{p,p}(\Om,d^{-aq})$。
\end{theorem}

\begin{proof}
设 $\Om$ 为 $\R^n$ 中的非空有界域，$a>0$，$s\in(0,1)$，$p\in[1,\infty)$，且 $\int_\Om d^{-ap}<\infty$。记
\[
  \|f\|_{F^s}=\|f\|_{L^p}+[f],\qquad
  [f]=\Big(\int\int |f(x)-f(y)|^p\,|x-y|^{-n-sp}\,d(y)^{-ap}\,dy\,dx\Big)^{1/p},
\]
并令 $\wt F^s$ 为由满足如下条件的 $f\in F^s$ 组成的集合：
\[
  \|f\|_{\wt F^s}:=\|f\|_{F^s}+\|f d^{-(s+a)}\|_{L^p}+\|f d^{-(s+a+n/p)}\|_{L^p}<\infty .
\]
令
\[
  w_H:=d(y)^{-(s+a)p},\qquad w_Q:=d(y)^{-(s+a)p-n}.
\]

\emph{外部输入.} Kijaczko（\texttt{arXiv:2110.04012}，定理 19）：对 $\R^d$ 中的非空开集 $\Om$，$s\in(0,1)$，$p\in[1,\infty)$，以及满足
\[
  \int w(x)(1+|x|)^{-d-sp}\,dx<\infty
  \quad\text{与}\quad
  \int v(x)(1+|x|)^{-d-sp}\,dx<\infty ,
\]
的非负局部有界权 $w,v$，记
\[
  \wt W^{s,p;w,v}(\Om)
  =\Big\{f : \int\int |f(x)-f(y)|^p|x-y|^{-d-sp}w(x)v(y)\,dx\,dy<\infty\Big\};
\]
则集合 $C^\infty(\Om)\cap\wt W^{s,p;w,v}(\Om)$ 在该半范空间中稠密；其证明使用 Whitney 磨光 $P^\eta f=\sum_n(f\psi_n)*h_{\eta(Q_n)}$，其中对所有 $n$ 有 $\eta(Q_n)<1/k$，而由 Dyda--Kijaczko（\texttt{arXiv:2009.00077}，定理 3），只要对所有 $n$ 有 $\eta(Q_n)\to0$，这些磨光就满足 $P^\eta f\to f$ 于 $L^p$ 中。取 $w=1$、$v=d^{-ap}$ 时假设成立，故\defterm{存在}一组满足 $\eta(Q_n)<1/k$ 的逐立方体尺度，使 $P^{\eta}f\to f$ 于 $F^s$ 中。
这是\defterm{存在性}陈述，不能读作“任意满足该上界的尺度都收敛”；后者要靠下文由所引定理证明导出的误差预算 \eqref{eq:seminorm-budget}。固定 $f\in\wt F^s$。

\emph{Whitney 立方体上的比较}. 在放大的立方体 $Q_n^{**}$ 上 $d$ 与 $l(Q_n)$ 可比，故
\[
  w_H\le C\,l(Q_n)^{-(s+a)p}=:H_n ,
  \qquad
  w_Q\le C'\,l(Q_n)^{-(s+a)p-n}=:H'_n ,
\]
而在较小的立方体 $Q_n^*$ 上有 $H_n\le C''d^{-(s+a)p}$ 与 $H'_n\le C''d^{-(s+a)p-n}$。令
\[
  u_n:=(f\psi_n)*h_{\eta_n}-f\psi_n ,
\]
其支集在 $Q_n^{**}$ 中；于是 $P^\eta f-f=\sum_n u_n$，且每一点至多属于 $N$ 个 $Q_n^{**}$，故
\begin{equation}
  \int|P^\eta f-f|^p\,w_H\le C(N)\sum_n H_n\int|u_n|^p .
  \tag{3.8}
\end{equation}
对每个 $n$，
\begin{equation}
  \int|u_n|^p\le 2^p\int h(z)\,\|\tau_{\eta_n z}(f\psi_n)-f\psi_n\|^p_{L^p}\,dz\longrightarrow0
  \tag{3.9}
\end{equation}
（平移在 $L^p$ 中的连续性）。又有
\begin{equation}
  \sum_n H_n\|f\psi_n\|^p_{L^p}
  \le C''\sum_n\int_{Q_n^*}|f|^p\,d^{-(s+a)p}
  \le C'''\int|f|^p\,w_H<\infty
  \tag{3.10}
\end{equation}
由有限重叠得到。

\emph{统一的一次尺度选择.}
至此出现了\defterm{四项}互不相同的要求：
式 (3.9) 后的收敛性需要 $\eta_k(Q_n)<1/k$；
\defterm{原加权分数阶半范数}需要 $[P^{\eta}f-f]_{F^s}<1/k$；
$w_H$ 项需要 $H_n\int|u_n(\eta)|^p\le 2^{-n}/(C(N)k^{p})$；$w_Q$ 项需要 $H_n'\int|u_n(\eta)|^p\le 2^{-n}/(C(N)k^{p})$——其中 $H_n,H_n'$ 是下方 (3.8) 之前定义的\defterm{Hardy 权常数}（不是局部常数 $A_n$），$C(N)$ 是 (3.8) 中由有限重叠产生的常数。
必须说明存在\defterm{同一族}尺度同时满足这四项，否则前面的估计只是若干互不相干的“存在某尺度”陈述。

\emph{第四项为何不能由前三项自动得到.}
所引定理只断言\defterm{存在}一组尺度使半范数误差趋于零；
而本文为控制两个 Hardy 项必须在每个立方体上\defterm{重新}选取（更小的）尺度。
关键在于 $[P^{\eta}f-f]_{F^s}$ 关于各尺度坐标 $\eta(Q_n)$ \defterm{并不单调}——减小某个 $\eta(Q_n)$ 未必使该半范数变小，
故“存在一组”不能推广为“任意更小尺度仍收敛”，也不能靠取所有阈值的下确界绕过。因此必须另取一个\defterm{阈值}。

\emph{半范数项：恢复所引定理的二变量加权估计.}
这里必须小心：所引定理 19 的证明\defterm{不是}对单变量函数 $f\psi_n$ 使用普通 $L^p(\Om)$ 中的平移连续性。
该证明的核心对象是定义在 $\R^{2n}$ 上的\defterm{二变量差商}
\begin{equation}\label{eq:gn-def}
  g_n(x,y):=\frac{f(x)\psi_n(x)-f(y)\psi_n(y)}{|x-y|^{\,n/p+s}}\;\mathbf 1_{\Om\times\Om}(x,y),
\end{equation}
相应地其估计为
\begin{equation}\label{eq:seminorm-budget}
  [P^{\eta}f-f]_{W^{s,p;w,v}}^{p}
  \;\le\;M\sum_{n=1}^{\infty}\int_{\R^n}
  \big\|\tau_{\eta(Q_n)u}g_n-g_n\big\|_{L^p(\R^{2n},\,w\times v)}^{p}\,h(u)\,du ,
  \qquad M=12^{\,n(p-1)},
\end{equation}
其中 $h$ 为磨光核（$\supp h=B(0,1)$、$\int h=1$、$h\in C^\infty$）、$\tau_v$ 表示平移、
$w\times v$ 指测度 $w(x)v(y)\,dx\,dy$。算子本身为
$P^{\eta}f=\sum_{n}(f\psi_n)*h_{\eta(Q_n)}$，其中 $\eta:\mathcal W\to(0,\infty)$ 须满足
\begin{equation}\label{eq:eta-admissible}
  \eta(Q)<2\varepsilon\,l(Q)\qquad(\forall Q\in\mathcal W),\qquad 0<\varepsilon<\sqrt[4]{5/4}-1<\tfrac14 ;
\end{equation}
这是算子有定义的前提（它保证 $P^\eta f\in C^\infty(\Om)$）。常数 $M=12^{n(p-1)}$ 来源于每点至多属于
$12^n$ 个放大立方体 $Q_n^{**}$。
在本节取 $w=1$、$v=d^{-ap}$ 时，\eqref{eq:seminorm-budget} 左端正是 $[\,\cdot\,]_{F^s_{p,p}(\Om,d^{-ap})}^{p}$，
故两个空间记号在此参数下等价。

必须指出两点，否则把 \eqref{eq:seminorm-budget} 换成单变量形式是\defterm{无效}的：
\begin{enumerate}[label=(\roman*)]
  \item \eqref{eq:seminorm-budget} 中的范数含奇异核 $|x-y|^{-n/p-s}$、指示函数 $\mathbf 1_{\Om\times\Om}$
        与权 $w\times v$，\defterm{不是}普通 $L^p(\Om)$ 范数，普通平移连续性无法控制它；
  \item $\omega:=w\times v$ 一般\defterm{不具平移不变性}，故 $\|\tau_{tu}G-G\|_{L^p(\omega)}\to0$ 不能由
        “$C_c$ 在 $L^p(\omega)$ 中稠密”两句推出（稠密性论证中 $\|\tau_{tu}(G-G')\|_{L^p(\omega)}$
        并不等于 $\|G-G'\|_{L^p(\omega)}$）。故下面的论证改用\defterm{局部项 + 交叉项}的三分解（见式 \eqref{eq:three-term}），而不在单一区域上比较权。
\end{enumerate}

\emph{$g_n$ 的支集，以及为什么必须做三分解.}
必须注意 $g_n$ 的支集\defterm{并不}落在 $Q_n^{**}\times Q_n^{**}$ 内：取
$x\in\supp\psi_n$ 与 $y\in\Om\setminus Q_n^{**}$，则 $\psi_n(y)=0$，但一般仍有
\[
  g_n(x,y)=\frac{f(x)\psi_n(x)}{|x-y|^{n/p+s}}\ \ne\ 0 .
\]
正确的支集包含关系是
\begin{equation}\label{eq:gn-support}
  \supp g_n\ \subset\ \big(Q_n^{**}\times\Om\big)\ \cup\ \big(\Om\times Q_n^{**}\big).
\end{equation}
同样地，$g_n$ \defterm{不是}全局差商 $g$ 在 $Q_n^{**}\times Q_n^{**}$ 上限制后乘有界因子：由
\begin{equation}\label{eq:gn-split}
  f(x)\psi_n(x)-f(y)\psi_n(y)
  =\psi_n(x)\bigl(f(x)-f(y)\bigr)+f(y)\bigl(\psi_n(x)-\psi_n(y)\bigr)
\end{equation}
可见第二项不能被 $f(x)-f(y)$ 的差商控制。因此不能只在 $Q_n^{**}\times Q_n^{**}$ 上比较权函数，
必须按原文把 $\|\tau_tg_n-g_n\|^p_{L^p(w\times v)}$ 分成\defterm{局部项与交叉项}。

\emph{三分解.} 记 $\omega:=w\times v$。对 $t<\eta(Q_n)$，按 $(x,y)$ 是否越过 $Q_n^{**}$ 拆开积分，
得（与文献 [\ref{ref:DK}] 定理 12 的证明、以及它在 [\ref{ref:K}] 定理 19 中的加权版本同构；
注意 $w\times v$ 不对称，故比对称情形多一项）
\begin{equation}\label{eq:three-term}
  \big\|\tau_tg_n-g_n\big\|^p_{L^p(\omega)}\ \le\ I_1+I_2+I_3 ,
\end{equation}
其中
\begin{align*}
  I_1&:=\int_{Q_n^{*}}\!\int_{Q_n^{**}}
   \frac{\bigl|f(x-t)\psi_n(x-t)-f(y-t)\psi_n(y-t)-f(x)\psi_n(x)+f(y)\psi_n(y)\bigr|^p}
        {|x-y|^{n+sp}}\,w(x)v(y)\,dy\,dx,\\
  I_2&:=\int_{Q_n^{*}}\!\int_{\Om\setminus Q_n^{**}}
   \frac{\bigl|f(x)\psi_n(x)-f(x-t)\psi_n(x-t)\bigr|^p}
        {|x-y|^{n+sp}}\,w(x)v(y)\,dy\,dx,\\
  I_3&:=\int_{Q_n^{*}}\!\int_{\Om\setminus Q_n^{**}}
   \frac{\bigl|f(x)\psi_n(x)-f(x-t)\psi_n(x-t)\bigr|^p}
        {|x-y|^{n+sp}}\,w(y)v(x)\,dy\,dx .
\end{align*}

\emph{第 1 步（局部项）}. 在 $Q_n^{**}$ 上 $w$、$v$ 均有界，记
\[
  A_n:=\Big(\sup_{Q_n^{**}}w\Big)\Big(\sup_{Q_n^{**}}v\Big)<\infty
  \qquad\bigl(w=1,\ v=d^{-ap}\asymp l(Q_n)^{-ap}\ \text{于}\ Q_n^{**}\bigr).
\]
因 $Q_n^{*}\subset Q_n^{**}$，
\[
  I_1\ \le\ A_n\int_{Q_n^{*}}\!\int_{Q_n^{**}}
  \frac{\bigl|\,\cdot\,\bigr|^p}{|x-y|^{n+sp}}\,dy\,dx
  \ \le\ A_n\big\|\tau_tg_n-g_n\big\|^p_{L^p(\R^{2n},dx)} .
\]

\emph{第 2 步（交叉项）}. 用文献 [\ref{ref:DK}] 引理 11：若 $x\in Q_n^{*}$ 而 $y\notin Q_n^{**}$，则
\[
  |x-y|\ \ge\ c\,|y-x_n|,\qquad c=\frac{\varepsilon}{\varepsilon+\sqrt n},
\]
$x_n$ 为 $Q_n$ 的中心。于是
\[
  I_2\le c^{-n-sp}\Big(\sup_{Q_n^{*}}w\Big)D_n^{(v)}
  \big\|\tau_t(f\psi_n)-f\psi_n\big\|^p_{L^p(dx)},\quad
  I_3\le c^{-n-sp}\Big(\sup_{Q_n^{*}}v\Big)D_n^{(w)}
  \big\|\tau_t(f\psi_n)-f\psi_n\big\|^p_{L^p(dx)},
\]
其中
\[
  D_n^{(v)}:=\int_{\Om\setminus Q_n^{**}}\frac{v(y)}{|y-x_n|^{n+sp}}\,dy,\qquad
  D_n^{(w)}:=\int_{\Om\setminus Q_n^{**}}\frac{w(y)}{|y-x_n|^{n+sp}}\,dy .
\]
令 $D_n:=c^{-n-sp}\max\big\{\sup_{Q_n^{*}}w\cdot D_n^{(v)},\ \sup_{Q_n^{*}}v\cdot D_n^{(w)}\big\}$，则
$I_2+I_3\le 2D_n\|\tau_t(f\psi_n)-f\psi_n\|^p_{L^p(dx)}$。

\emph{$D_n$ 的显式界（在本文假设下有限）}.
当 $y\notin Q_n^{**}$ 时 $|y-x_n|\ge\tfrac12l(Q_n^{**})\ge\tfrac12l(Q_n)$，故
\begin{equation}\label{eq:Dn-bounds}
  \begin{aligned}
    D_n^{(v)}&=\int_{\Om\setminus Q_n^{**}}\frac{d(y)^{-ap}}{|y-x_n|^{n+sp}}\,dy
    \ \le\ \Big(\frac{l(Q_n)}{2}\Big)^{-n-sp}\int_\Om d(y)^{-ap}\,dy\\
    &=\ 2^{\,n+sp}\,l(Q_n)^{-n-sp}\int_\Om d(y)^{-ap}\,dy\ <\ \infty ,
  \end{aligned}
\end{equation}
这里用的正是本节假设 $\int_\Om d^{-ap}dy<\infty$（即 $v=d^{-ap}\in L^1(\Om)$，与 [\ref{ref:K}] 条件 (6)
在有界域上的等价形式一致）。同理，由 $w=1$ 得
\begin{equation}\label{eq:Dnw}
  D_n^{(w)}=\int_{\Om\setminus Q_n^{**}}\frac{dy}{|y-x_n|^{n+sp}}
  \ \le\ 2^{\,n+sp}l(Q_n)^{-n-sp}|\Om|\ <\ \infty .
\end{equation}
又 $\sup_{Q_n^{*}}w=1$、$\sup_{Q_n^{*}}v\asymp l(Q_n)^{-ap}$，故
\[
  D_n\ \le\ c^{-n-sp}\max\big\{D_n^{(v)},\ C\,l(Q_n)^{-ap}D_n^{(w)}\big\}\ <\ \infty ,
\]
其中 $D_n^{(v)},D_n^{(w)}$ 由 \eqref{eq:Dn-bounds}、\eqref{eq:Dnw} 估计。

\emph{第 3 步（两个误差预算，均由 Lebesgue 平移连续性导出）}.
所需的是 $g_n\in L^p(\R^{2n},dx)$ 与 $f\psi_n\in L^p(\R^n,dx)$。后者显然：$|\psi_n|\le1$ 给出
$\|f\psi_n\|_{L^p}\le\|f\|_{L^p}<\infty$。前者需要一点估计，我们把它写全。

\emph{引理. $g_n\in L^p(\R^{2n},dx)$。}
由 \eqref{eq:gn-split} 作两项估计，记
\[
  A:=\int_{\R^n}\!\!\int_{\R^n}\frac{|f(x)-f(y)|^p|\psi_n(x)|^p}{|x-y|^{n+sp}}\,dy\,dx,
  \qquad
  B:=\int_{\R^n}\!\!\int_{\R^n}\frac{|f(y)|^p|\psi_n(x)-\psi_n(y)|^p}{|x-y|^{n+sp}}\,dy\,dx .
\]
\emph{项 $A$.} 因 $\supp\psi_n\subset Q_n^{*}$ 而 $d\asymp l(Q_n)$ 于 $Q_n^{*}$ 上，把 $y$ 的积分分两段：
(I) $|x-y|<\tfrac12d(x)$ 时 $d(y)\asymp d(x)\asymp l(Q_n)$，故 $1\le C\,l(Q_n)^{ap}d(y)^{-ap}$，于是
\[
  A_{\mathrm{I}}\ \le\ C\,l(Q_n)^{ap}\int_\Om\!\!\int_\Om\frac{|f(x)-f(y)|^p}{|x-y|^{n+sp}}d(y)^{-ap}dy\,dx
  \ =\ C\,l(Q_n)^{ap}[f]^p_{F^s};
\]
(II) $|x-y|\ge\tfrac12d(x)$ 时 $|x-y|\ge c\,l(Q_n)$。此处\defterm{不能}把因子 $|Q_n^{*}|$ 同时加到第一项上：
由 $|f(x)-f(y)|^p\le2^{p-1}\big(|f(x)|^p+|f(y)|^p\big)$ 只能得到
\[
  \int_{Q_n^{*}}\!\!\int_\Om|f(x)-f(y)|^p\,dy\,dx
  \ \le\ 2^{p-1}\big(|\Om|+|Q_n^{*}|\big)\|f\|^p_{L^p} ,
\]
故
\[
  A_{\mathrm{II}}\ \le\ c^{-n-sp}l(Q_n)^{-n-sp}\cdot 2^{p-1}\big(|\Om|+|Q_n^{*}|\big)\|f\|^p_{L^p}\ <\ \infty .
\]
因 $|Q_n^{*}|\asymp l(Q_n)^n$，该界与 $l(Q_n)^{-sp}\|f\|^p_{L^p}$ 同为 $n$ 的确定函数；
一般情形下它\defterm{依赖}于 $n$（$f$ 的 $L^p$ 质量可以集中在很小的立方体中），
但对\defterm{每个固定的} $n$ 均有限，而这正是本节所需的全部。

\emph{项 $B$.} 用 $\psi_n\in C_c^\infty$ 的 Lipschitz 界与 $|\psi_n|\le1$：
$|\psi_n(x)-\psi_n(y)|\le C\min\{1,|x-y|/l(Q_n)\}$（同 [\ref{ref:DK}] 中分割函数的性质）。
于是 $|\psi_n(x)-\psi_n(y)|\ne0$ 迫使 $x$ 或 $y$ 属于 $\supp\psi_n\subset Q_n^{*}$。分两段：
(I\('\)) $|x-y|\le 2l(Q_n)$：$|\psi_n(x)-\psi_n(y)|^p\le C^p l(Q_n)^{-p}|x-y|^{p}$，于是
\[
  B_{\mathrm{I}'}\le C^p l(Q_n)^{-p}\int_\Om|f(y)|^p\!\!\int_{|x-y|\le 2l(Q_n)}\!\!|x-y|^{p-n-sp}dx\,dy
  \ \lesssim\ l(Q_n)^{-sp}\|f\|^p_{L^p},
\]
其中用到 $p-n-sp>-n$（即 $p>sp$，由 $0<s<1\le p$ 保证）故内层积分 $\asymp l(Q_n)^{p-sp}$；
(II\('\)) $|x-y|>2l(Q_n)$：此时 $|\psi_n(x)-\psi_n(y)|\le2$，且
$\int_{|x-y|>2l(Q_n)}|x-y|^{-n-sp}dx\lesssim l(Q_n)^{-sp}$，故
$B_{\mathrm{II}'}\lesssim l(Q_n)^{-sp}\|f\|^p_{L^p}$。

合并得
\[
  \|g_n\|^p_{L^p(\R^{2n},dx)}\ \le\ 2^{p-1}(A+B)
  \ \lesssim\ l(Q_n)^{ap}[f]^p_{F^s}+l(Q_n)^{-sp}\|f\|^p_{L^p}\ <\ \infty .
\]
（注意该界依赖于 $n$，这是允许的：它只用于确定阈值 $\theta^{(1a)}_{n,k}$，而该阈值本就允许依赖于 $n$ 与 $k$。）

于是\defterm{普通}（Lebesgue）平移连续性给出，并关于 $|t|$ 一致地，
\[
  \big\|\tau_tg_n-g_n\big\|_{L^p(\R^{2n},dx)}\longrightarrow0,
  \qquad
  \big\|\tau_t(f\psi_n)-f\psi_n\big\|_{L^p(\R^n,dx)}\longrightarrow0
  \qquad(t\to0).
\]
据此取阈值 $\theta^{(1a)}_{n,k},\theta^{(1b)}_{n,k}>0$，使得当
$0<|t|<\min\{\theta^{(1a)}_{n,k},\theta^{(1b)}_{n,k}\}$ 时
\[
  \big\|\tau_tg_n-g_n\big\|^p_{L^p(dx)}<\frac{1}{Mk^{p}2^{n+1}A_n},
  \qquad
  \big\|\tau_t(f\psi_n)-f\psi_n\big\|^p_{L^p(dx)}<\frac{1}{Mk^{p}2^{n+2}D_n}.
\]
（注意此处\defterm{仅}用到 Lebesgue 范数，不涉及非平移不变的权 $w\times v$；
（注意此处\defterm{仅}用到 Lebesgue 范数，不涉及非平移不变的权 $w\times v$；加权的作用完全体现在 $A_n$ 与 $D_n$ 这两个有限常数上。）

\emph{第 4 步（合并）}. 令 $\theta^{(1)}_{n,k}:=\min\{\theta^{(1a)}_{n,k},\theta^{(1b)}_{n,k}\}$。
由第 1--3 步及 \eqref{eq:three-term}，只要 $0<\eta(Q_n)<\theta^{(1)}_{n,k}$，
\[
  \big\|\tau_{\eta(Q_n)u}g_n-g_n\big\|^p_{L^p(\omega)}
  \le I_1+I_2+I_3
  \le \frac{1}{Mk^{p}2^{n+1}}+\frac{2}{Mk^{p}2^{n+2}}
  \le \frac{1}{Mk^{p}2^{n}} ,
\]
对一切 $|u|\le1$ 成立（因 $\supp h\subset B(0,1)$）。代入 \eqref{eq:seminorm-budget} 并注意 $\int h=1$，
\[
  [P^{\eta}f-f]_{F^s}^{p}
  \ \le\ M\sum_{n}\frac{1}{Mk^{p}2^{n}}\ \le\ \frac{1}{k^{p}},
  \qquad\text{即}\qquad
  [P^{\eta}f-f]_{F^s}\ \le\ \frac1k .
\]
此结论对\defterm{任意}满足该上界的尺度族成立，这正是下文把 $\theta^{(1)}_{n,k}$ 并入统一选择的原因。

为此固定 $k\in\N$ 与立方体下标 $n$，记
\[
  e_n(\eta):=\int|u_n(\eta)|^p,\qquad u_n(\eta):=(f\psi_n)*h_\eta-f\psi_n .
\]
由平移在 $L^p$ 中的连续性，对每个固定的 $n$ 有 $e_n(\eta)\to0$（$\eta\to0$）。于是可分别取阈值
$\theta^{(2)}_{n,k}>0$、$\theta^{(3)}_{n,k}>0$，使得
\[
  0<\eta<\theta^{(2)}_{n,k}\ \Longrightarrow\ H_n\,e_n(\eta)\le\frac{2^{-n}}{C(N)k^{p}},
  \qquad
  0<\eta<\theta^{(3)}_{n,k}\ \Longrightarrow\ H'_n\,e_n(\eta)\le\frac{2^{-n}}{C(N)k^{p}}.
\]
令
\[
  \eta_k(Q_n):=\frac12\min\Big\{\frac1k,\ 2\varepsilon\,l(Q_n),\ \theta^{(1)}_{n,k},\ \theta^{(2)}_{n,k},\ \theta^{(3)}_{n,k}\Big\}>0 .
\]
因 $\eta_k(Q_n)\le\varepsilon\,l(Q_n)<2\varepsilon\,l(Q_n)$，算子定义所需的条件 \eqref{eq:eta-admissible} 成立；
又因 $\eta_k(Q_n)\le 1/(2k)$，
对每个固定的 $n$ 有 $\eta_k(Q_n)\to0$（$k\to\infty$），故 Dyda--Kijaczko 定理 3 所需的
$\eta(Q_n)\to0$ 亦满足。

\emph{逐项收敛.}
\textbf{(半范数项)} 由 \eqref{eq:seminorm-budget} 及 $\eta_k(Q_n)<\theta^{(1)}_{n,k}$，
\[
  [P^{\eta_k}f-f]_{F^s}<\frac1k\longrightarrow0 .
\]
\textbf{($L^p$ 项)} 由 $\eta_k(Q_n)<1/k$ 与 Dyda--Kijaczko 定理 3，$\|P^{\eta_k}f-f\|_{L^p}\to0$。
\textbf{($w_H$ 项)} 由式 (3.8) 与 $\eta_k$ 的第三项要求，
\[
  \int|P^{\eta_k}f-f|^p\,w_H\le C(N)\sum_n H_n\int|u_n|^p\le\frac{1}{k^{p}}\longrightarrow0 .
\]
\textbf{($w_Q$ 项)} 同法，由第四项要求，并利用
\[
  \sum_n H'_n\|f\psi_n\|^p_{L^p}\le C'''\int|f|^p\,w_Q<\infty ,
\]
故 $\int|P^{\eta_k}f-f|^p\,w_Q\le 1/k^{p}\longrightarrow0$。

由以上四项，$\|P^{\eta_k}f-f\|_{\wt F^s}\le 1/k+1/k+1/k=3/k\longrightarrow0$，即\defterm{同一族} $\{P^{\eta_k}f\}_{k}$ 在 $\wt F^s$ 范数中收敛于 $f$；
求和的可控性来自 $Q_n^{**}$ 的有限重叠（每点至多属于 $N:=12^{n}$ 个），与逐立方体选取的尺度无关。

\emph{逼近序列属于交集}. 最后，
\[
  \int|P^\eta f|^p\,w_H\le C(N)\sum_n H_n\|f\psi_n\|^p_{L^p}<\infty
\]
对 $w_Q$ 同样成立，故每个 $P^{\eta_k}f$ 属于 $C^\infty\cap\wt F^s$。因此 $C^\infty\cap\wt F^s$ 在 $\wt F^s$ 中稠密。
\end{proof}

\begin{remark}[直觉]
在一个 Whitney 立方体上，两个 Hardy 权重都与常数可比，因此 Hardy 加权的误差是一个有限加权的局部误差之和；由于 $f$ 属于 Hardy 加强空间，这个和收敛，而尺度的对角选取使误差趋于零。
\end{remark}

%==================================================================
\subsection{情形 (I) 中区间上 \texorpdfstring{$s>1$}{s>1} 时的稠密性}

\begin{theorem}[情形 (I)：$s\in(1,\infty)\setminus\N$、$p=q$ 时区间上 $C^\infty$ 的稠密性]
设 $\Om:=(0,1)$，并对 $x\in(0,1)$ 令 $d(x):=\min(x,1-x)=\dist(x,\{0,1\})$。设 $s\in(1,\infty)\setminus\N$，令 $k:=\lfloor s\rfloor$，$\sigma:=s-k\in(0,1)$，设 $q\in[1,\infty)$，$a>0$。设 $F^s_{q,q}((0,1),d^{-aq})$ 为加权内蕴 Triebel--Lizorkin 空间，其范数为
\[
  \|f\|:=\sum_{m=0}^{k}\|f^{(m)}\|_{L^q(0,1)}+[f^{(k)}]_{\sigma,q,a},
\]
其中 $f^{(m)}$ 为弱导数，
\[
  [g]_{\sigma,q,a}:=\Big(\int_0^1\int_0^1|g(x)-g(y)|^q\,|x-y|^{-1-\sigma q}\,d(y)^{-aq}\,dy\,dx\Big)^{1/q}.
\]
则 $C^\infty(0,1)\cap F^s_{q,q}((0,1),d^{-aq})$ 在 $F^s_{q,q}((0,1),d^{-aq})$ 中稠密；等价地，$C^\infty(0,1)\cap F^s$ 在 $F^s$ 中的闭包等于 $F^s$。
\end{theorem}

\begin{proof}
设 $\Om:=(0,1)$，$d(x):=\min(x,1-x)$。$f$ 的内蕴加权范数为
\[
  \|f\|:=\sum_{m=0}^{k}\|f^{(m)}\|_{L^q(0,1)}+[f^{(k)}]_{\sigma,q,a},
\]
其中
\[
  [g]_{\sigma,q,a}:=\Big(\int_0^1\int_0^1|g(x)-g(y)|^q\,|x-y|^{-1-\sigma q}\,d(y)^{-aq}\,dy\,dx\Big)^{1/q};
\]
这是 CLWZ (1.2) 的一维情形（恰有一个多重指标 $\beta$ 满足 $|\beta|=k$，即 $\beta=(k)$，且 $D^\beta f=f^{(k)}$）。

\emph{情形 A：$aq<1$}. 此时
\[
  \int_0^1 d(y)^{-aq}\,dy=2\int_0^{1/2}y^{-aq}\,dy<\infty .
\]
设 $f\in F^s_{q,q}((0,1),d^{-aq})$，令 $g:=f^{(k)}$。则 $g\in L^q(0,1)$ 且 $[g]_{\sigma,q,a}<\infty$，即 $g$ 属于 CLWZ (1.2) 中 $s=\sigma\in(0,1)$、$p=q$ 的分数阶空间 $F^\sigma_{q,q}((0,1),d^{-aq})$。由附录 A（情形 (I) 中 $p=q$、$s\in(0,1)$ 的 $C^\infty$ 稠密性，应用于 $\Om=(0,1)$、$s=\sigma$、$p=q$ 以及总积分有限的同一权 $d^{-aq}$），集合 $C^\infty(0,1)\cap F^\sigma_{q,q}$ 在 $F^\sigma_{q,q}$ 中稠密。故存在 $C^\infty(0,1)$ 中的序列 $(g_j)$，使得
\begin{equation}
  \|g_j-g\|_{L^q(0,1)}\to0
  \quad\text{与}\quad
  [g_j-g]_{\sigma,q,a}\to0 .
  \tag{3.11}
\end{equation}
对 $x\in(0,1)$ 定义
\[
  f_j(x):=\sum_{m=0}^{k-1}\frac{f^{(m)}(0)}{m!}\,x^m
  +\frac{1}{(k-1)!}\int_0^x (x-t)^{k-1}g_j(t)\,dt .
\]
数 $f^{(m)}(0)$ 是定义良好的：$f\in W^{k,q}(0,1)$ 且 $q\ge1$ 时，由 $f^{(k-1)}$ 的绝对连续性与 $f^{(k)}$ 的可积性可知，极限
\[
  f^{(m)}(0):=\lim_{x\to0+}f^{(m)}(x)
\]
对 $0\le m\le k-1$ 存在且有限。每个 $g_j$ 在 $(0,1)$ 上光滑，故 $f_j$ 在 $(0,1)$ 上光滑。对积分项求导 $k$ 次得 $f_j^{(k)}=g_j$，且对 $0\le m\le k-1$ 有 $f_j^{(m)}(0)=f^{(m)}(0)$。

带积分余项的 Taylor 公式对 $f\in W^{k,q}(0,1)$ 成立，因为 $f^{(k-1)}$ 绝对连续且 $f^{(k)}\in L^1(0,1)$；它给出，对 a.e.\ $x\in(0,1)$，
\begin{equation}
  f(x)=\sum_{m=0}^{k-1}\frac{f^{(m)}(0)}{m!}\,x^m
  +\frac{1}{(k-1)!}\int_0^x (x-t)^{k-1}g(t)\,dt .
  \tag{3.12}
\end{equation}
于是，记 $h:=g_j-g$，$e_j:=f_j-f$，有
\begin{equation}
  e_j(x)=\frac{1}{(k-1)!}\int_0^x (x-t)^{k-1}h(t)\,dt ,
  \tag{3.13}
\end{equation}
且对 $0\le m\le k-1$，
\begin{equation}
  e_j^{(m)}(x)=\frac{1}{(k-1-m)!}\int_0^x (x-t)^{k-1-m}h(t)\,dt
  \tag{3.14}
\end{equation}
（因为 $f_j$ 与 $f$ 的 Taylor 多项式部分相互抵消）。对 $0\le m\le k-1$，H\"older 不等式给出
\begin{equation}
  |e_j^{(m)}(x)|
  \le\frac{1}{(k-1-m)!}\Big(\int_0^x (x-t)^{(k-1-m)q'}\,dt\Big)^{1/q'}\|h\|_{L^q(0,1)}
  \le C_{k,q}\,\|h\|_{L^q(0,1)} ,
  \tag{3.15}
\end{equation}
其中 $q'$ 为共轭指数（当 $q=1$ 时取 $q'=\infty$，估计立即成立），且 $t$ 积分有限是因为 $(k-1-m)q'\ge0$。因此
\begin{equation}
  \|e_j^{(m)}\|_{L^q(0,1)}\le C_{k,q}\,\|h\|_{L^q(0,1)}\longrightarrow0
  \quad\text{for } 0\le m\le k-1,
  \qquad
  e_j^{(k)}=h ,
  \tag{3.16}
\end{equation}
故由 (3.11) 得 $\|f_j-f\|_{W^{k,q}(0,1)}=\sum_{m=0}^{k}\|e_j^{(m)}\|_{L^q(0,1)}\to0$。此外由 (3.11)
\[
  [f_j^{(k)}-f^{(k)}]_{\sigma,q,a}=[h]_{\sigma,q,a}\longrightarrow0
\]
因此 $\|f_j-f\|\to0$，且 $f_j\in C^\infty(0,1)\cap F^s_{q,q}((0,1),d^{-aq})$，因为 $f_j$ 光滑，其阶数不超过 $k-1$ 的导数有界（$0\le m\le k-1$ 时 $|f_j^{(m)}|$ 被多项式部分加上 $C_{k,q}\|g_j\|_{L^q}$ 控制）；最高阶导数按构造满足 $f_j^{(k)}=g_j$，而 $g_j$ 只属于目标空间 $F^\sigma_{q,q}$，\defterm{一般并不有界}。附带说明：本步只用到 $m\le k-1$ 的有界性。这就证明了情形 A 中的稠密性。

\emph{情形 B：$aq\ge1$}. 此时
\[
  \int_0^1 d(y)^{-aq}\,dy=\infty ,
\]
故对权 $d^{-aq}$，$\Om=(0,1)$ 属于情形 (II)。由附录 A（非可积情形下的退化，取 $n=1$、$\Om=(0,1)$、$a$、$q$、$p=q$、$s\in(0,\infty)\setminus\N$），加权内蕴空间 $F^s_{q,q}((0,1),d^{-aq})$ 等于 $P_k$，即次数不超过 $k$ 的实多项式在 $(0,1)$ 上限制的集合。每个这样的多项式都在 $(0,1)$ 上无穷次可微，故 $C^\infty(0,1)\cap F^s=F^s$，稠密性陈述平凡成立。

两种情形合起来覆盖了所有 $a>0$，证明完成。
\end{proof}

\begin{remark}[直觉]
在一维情形下最高阶导数可以积分回去：每个 $f\in F^s$ 都能由 $g=f^{(k)}$ 通过带积分余项的 Taylor 公式还原。于是 $g$ 在 $\sigma$ 阶空间 $F^\sigma$ 中的逼近 $g_j$（由已验证的 $s\in(0,1)$ 可积权稠密性事实给出）可以再积分成光滑的 $f_j$，满足 $f_j^{(k)}=g_j$；Taylor 余项则控制所有低阶导数。这完全绕开了 Leibniz 障碍，因为并没有对 $f$ 作截断。
\end{remark}


%=====================================================================
\section{主强化定理与反例}\label{sec:part4}
%=====================================================================
% ===================================================================
%  第 4 组：主强化定理与反例
% ===================================================================

本节把前面两种情形合并为整个论证的主强化定理：当 $s\in(0,1)$ 且 $p=q$ 时，Hardy 加强空间 $\wt F^s_{p,p}(\Om,d^{-aq})$ 中 $C_c^\infty$ 与 $C^\infty$ 的闭包不但彼此相等，而且都等于整个空间。随后的小节处理权不可积的情形 (II)，说明此时 $\wt F^s$ 整体退化为零空间，而较弱的 $\wh F^s$ 只在 $L^p$ Hardy 项可积时保留常数函数，且 $C_c^\infty$ 的闭包恒为零；最后一小节给出区间 $(0,1)$ 上的具体反例，表明 $C^\infty$ 可以稠密而 $C_c^\infty$ 不稠密，因此形如 $F^{s,1}=F^s$ 的加强在一般情况下不成立，而仅仅是两个闭包相等这一"修正的稠密性"才是可证的。

%==================================================================
\subsection{主强化定理：\texorpdfstring{$C_c^\infty$}{Cc-infinity} 的闭包等于整个 Hardy 加强空间}

\begin{theorem}[关于 $s\in(0,1)$ 与 $p=q$ 的强化：$C_c^\infty$ 的闭包等于整个 Hardy 加强空间]
设 $n\ge 1$，$\Om$ 为 $\R^n$ 中的非空有界域，$d(x)=\dist(x,\partial\Om)$，$a>0$，$s\in(0,1)$，$p\in[1,\infty)$（本节记 $q:=p$，故 $d^{-aq}=d^{-ap}$，两种写法通用）。设 $\wt F^s_{p,p}(\Om,d^{-aq})$ 为附录 A 中的 Hardy 加强空间，即所有满足
\[
  \|f\|_{\wt F^s}:=\|f\|_{F^s}+\|f d^{-(s+a)}\|_{L^p}+\|f d^{-(s+a+n/p)}\|_{L^p}<\infty ,
\]
的 $f\in F^s_{p,p}(\Om,d^{-aq})$ 所构成的空间，并设 $\wt F^{s,1}_{p,p}(\Om,d^{-aq})$ 与 $\wt F^{s,2}_{p,p}(\Om,d^{-aq})$ 分别为
\[
  C_c^\infty(\Om)\cap\wt F^s_{p,p}(\Om,d^{-aq})
  \quad\text{与}\quad
  C^\infty(\Om)\cap\wt F^s_{p,p}(\Om,d^{-aq})
\]
在 $\wt F^s_{p,p}(\Om,d^{-aq})$ 中的闭包。则
\[
  \wt F^{s,1}_{p,p}(\Om,d^{-aq})
  =\wt F^{s,2}_{p,p}(\Om,d^{-aq})
  =\wt F^s_{p,p}(\Om,d^{-aq}) .
\]
特别地，$C_c^\infty(\Om)$ 在 $\wt F^s_{p,p}(\Om,d^{-aq})$ 中稠密。
\end{theorem}

\begin{proof}
设 $\Om$ 为 $\R^n$ 中的非空有界域，$a>0$，$s\in(0,1)$，$p\in[1,\infty)$，并令
\[
  \wt F^s:=\wt F^s_{p,p}(\Om,d^{-aq})
\]
为 Hardy 加强空间，其范数为
\[
  \|f\|_{\wt F^s}:=\|f\|_{F^s}+\|f d^{-(s+a)}\|_{L^p}+\|f d^{-(s+a+n/p)}\|_{L^p} .
\]
设 $\wt F^{s,1}$ 与 $\wt F^{s,2}$ 分别为 $C_c^\infty(\Om)\cap\wt F^s$ 与 $C^\infty(\Om)\cap\wt F^s$ 在 $\wt F^s$ 中的闭包。

\emph{外部输入}。我们使用如下完整陈述（\texttt{CLWZ\_2026}，Theorem~1.2(i)）：对 $\R^n$ 中的每个有界域 $\Om$、每个 $s\in(0,1)$、每个 $p,q\in[1,\infty)$ 以及每个 $a\in(0,\infty)$，均有
\[
  \wt F^{s,1}_{p,q}(\Om,d^{-aq})=\wt F^{s,2}_{p,q}(\Om,d^{-aq}) .
\]

\emph{第一步（应用二择一定理）}。由附录 A，下列两种情形恰有一种成立。

\emph{情形 (I)：$\int_\Om d(y)^{-ap}\,dy<\infty$}。此时由附录 A 可知 $C^\infty(\Om)\cap\wt F^s$ 在 $\wt F^s$ 中稠密，即
\[
  \wt F^{s,2}=\wt F^s ;
\]
结合外部 Theorem~1.2(i)（取 $q=p$）可得
\[
  \wt F^{s,1}=\wt F^{s,2}=\wt F^s .
\]

\emph{情形 (II)：$\int_\Om d(y)^{-ap}\,dy=\infty$}。此时由附录 A 得 $\wt F^s=\{0\}$，从而
\[
  \wt F^{s,1}=\wt F^{s,2}=\wt F^s=\{0\} .
\]

\emph{第二步（合并两种情形）}。由于情形 (I) 与 (II) 穷尽所有可能，故在任何情形下都有
\[
  \wt F^{s,1}_{p,p}(\Om,d^{-aq})
  =\wt F^{s,2}_{p,p}(\Om,d^{-aq})
  =\wt F^s_{p,p}(\Om,d^{-aq}) .
\]
特别地，$C_c^\infty(\Om)\cap\wt F^s$ 在 $\wt F^s$ 中稠密，即 $C_c^\infty(\Om)$ 在 $\wt F^s_{p,p}(\Om,d^{-aq})$ 中稠密。
\end{proof}

\begin{remark}[直觉]
在情形 (I) 中，$C^\infty$ 的闭包已经填满整个 Hardy 加强空间，而外部定理 Theorem~1.2(i) 又把两个闭包等同起来；在情形 (II) 中，三个空间一起退化为 $\{0\}$。两种情形穷尽了所有可能，所以结论无条件成立。
\end{remark}

%==================================================================
\subsection{非可积权情形：\texorpdfstring{$\wt F^s$}{tilde-F-s} 空间退化为 \texorpdfstring{$\{0\}$}{\{0\}} 与 \texorpdfstring{$\wh F^s$}{hat-F-s} 的结构}

\begin{theorem}[权不可积情形下 $s\in(0,1)$ 的强化]
设 $n\ge 1$，$\Om$ 为 $\R^n$ 中的非空有界域，$d(x)=\dist(x,\partial\Om)$，$0<s<1$，$p,q\in[1,\infty)$，$a>0$，并假设
\[
  \int_\Om d(y)^{-aq}\,dy=\infty
\]
（附录 A 的情形 (II)）。设 $F^s_{p,q}(\Om,d^{-aq})$ 为附录 A 中的加权内蕴空间。对 $f\in F^s$ 定义
\[
  \|f\|_{\wt F^s}:=\|f\|_{F^s}
  +\|f d^{-(s+a)}\|_{L^p(\Om)}
  +\|f d^{-(s+a+n/q)}\|_{L^q(\Om)}
\]
与
\[
  \|f\|_{\wh F^s}:=\|f\|_{F^s}+\|f d^{-(s+a)}\|_{L^p(\Om)} ,
\]
并置
\[
  \wt F^s_{p,q}(\Om,d^{-aq}):=\{f\in F^s : \|f\|_{\wt F^s}<\infty\},
  \qquad
  \wh F^s_{p,q}(\Om,d^{-aq}):=\{f\in F^s : \|f\|_{\wh F^s}<\infty\} .
\]
设 $\wt F^{s,1}$、$\wh F^{s,1}$ 分别为 $C_c^\infty(\Om)\cap\wt F^s$ 在 $\wt F^s$ 中、$C_c^\infty(\Om)\cap\wh F^s$ 在 $\wh F^s$ 中的闭包，并设 $\wt F^{s,2}$、$\wh F^{s,2}$ 分别为 $C^\infty(\Om)\cap\wt F^s$、$C^\infty(\Om)\cap\wh F^s$ 在相应范数下的闭包。则：
\begin{enumerate}[label=(\roman*)]
  \item $\wt F^s_{p,q}(\Om,d^{-aq})=\{0\}$，从而
  \[
    \wt F^{s,1}=\wt F^{s,2}=\wt F^s=\{0\} .
  \]
  \item 若 $\int_\Om d(y)^{-(s+a)p}\,dy=\infty$，则 $\wh F^s_{p,q}(\Om,d^{-aq})=\{0\}$；若 $\int_\Om d(y)^{-(s+a)p}\,dy<\infty$，则 $\wh F^s_{p,q}(\Om,d^{-aq})=\R\cdot 1$（即常数函数）。
  \item 在所有情形下 $\wh F^{s,2}=\wh F^s$ 且 $\wh F^{s,1}=\{0\}$；因此 $\wh F^{s,1}=\wh F^s$ 当且仅当 $\int_\Om d(y)^{-(s+a)p}\,dy=\infty$，而当 $\int_\Om d(y)^{-(s+a)p}\,dy<\infty$ 时强化 $\wh F^{s,1}=\wh F^s$ 不成立。
\end{enumerate}
\end{theorem}

\begin{proof}
设 $n\ge 1$，$\Om$ 为 $\R^n$ 中的非空有界域，$d(x)=\dist(x,\partial\Om)$，$0<s<1$，$p,q\in[1,\infty)$，$a>0$，并假设 $\int_\Om d(y)^{-aq}\,dy=\infty$（情形 (II)）。由附录 A（DEGENERACY THEOREM，取 $k=0$ 与 $\sigma=s\in(0,1)$）可得
\[
  F^s_{p,q}(\Om,d^{-aq})=P_0=\R\cdot 1 ,
\]
即 $\Om$ 上的常数函数，并且
\[
  C_c^\infty(\Om)\cap F^s=\{0\} .
\]

对 $f\in F^s$ 置
\[
  \|f\|_{\wt F^s}:=\|f\|_{F^s}+\|f d^{-(s+a)}\|_{L^p}+\|f d^{-(s+a+n/q)}\|_{L^q}
\]
与
\[
  \|f\|_{\wh F^s}:=\|f\|_{F^s}+\|f d^{-(s+a)}\|_{L^p} ,
\]
并定义
\[
  \wt F^s:=\{f\in F^s : \|f\|_{\wt F^s}<\infty\},
  \qquad
  \wh F^s:=\{f\in F^s : \|f\|_{\wh F^s}<\infty\} .
\]
设 $\wt F^{s,1}$、$\wh F^{s,1}$ 分别为 $C_c^\infty(\Om)\cap\wt F^s$ 在 $\wt F^s$ 中、$C_c^\infty(\Om)\cap\wh F^s$ 在 $\wh F^s$ 中的闭包，并设 $\wt F^{s,2}$、$\wh F^{s,2}$ 分别为 $C^\infty(\Om)\cap\wt F^s$、$C^\infty(\Om)\cap\wh F^s$ 的闭包。

\emph{第一步（$\wt F^s=\{0\}$）}。由于 $F^s=\R\cdot 1$，每个 $f\in F^s$ 都是常数 $c$。对 $c\ne 0$，有
\[
  \|c\,d^{-(s+a+n/q)}\|^q_{L^q(\Om)}
  =|c|^q\int_\Om d(y)^{-[(s+a)q+n]}\,dy .
\]
又 $(s+a)q+n\ge aq$，这是因为
\[
  (s+a)q+n-aq=sq+n>0 .
\]
于是由附录 A（MONOTONICITY OF DIVERGENCE，取 $c_1=aq$ 与 $c_2=(s+a)q+n$）可知，最后的积分发散至无穷。因此对每个 $c\ne 0$ 都有 $\|c\|_{\wt F^s}=\infty$，故
\[
  \wt F^s=\{0\} .
\]
由于零函数同时属于 $C_c^\infty(\Om)$ 与 $C^\infty(\Om)$，故得 $C^\infty(\Om)\cap\wt F^s=\{0\}=\wt F^s$ 且 $C_c^\infty(\Om)\cap\wt F^s=\{0\}$；因此两个闭包都等于 $\{0\}$，即
\[
  \wt F^{s,1}=\wt F^{s,2}=\wt F^s=\{0\} .
\]

\emph{第二步（$\wh F^s$ 的结构）}。仍有 $F^s=\R\cdot 1$。对常数 $c$，有
\[
  \|c\,d^{-(s+a)}\|^p_{L^p(\Om)}
  =|c|^p\int_\Om d(y)^{-(s+a)p}\,dy .
\]
因此，若 $\int_\Om d(y)^{-(s+a)p}\,dy=\infty$，则 $\wh F^s=\{0\}$；若 $\int_\Om d(y)^{-(s+a)p}\,dy<\infty$，则 $\wh F^s=\R\cdot 1$（在后一情形中，常数函数 $1$ 具有有限的 $F^s$ 范数 $|\Om|^{1/p}$ 与有限的 $L^p$ Hardy 项）。

\emph{第三步（$\wh F$ 的两个闭包）}。由于 $\wh F^s$ 包含于 $F^s=\R\cdot 1$，$\wh F^s$ 的每个元素都是常数，从而属于 $C^\infty(\Om)$；因此 $C^\infty(\Om)\cap\wh F^s=\wh F^s$，且闭包 $\wh F^{s,2}$ 等于 $\wh F^s$。另一方面
\[
  C_c^\infty(\Om)\cap\wh F^s\subset C_c^\infty(\Om)\cap F^s=\{0\},
\]
故 $C_c^\infty(\Om)\cap\wh F^s=\{0\}$；因此
\[
  \wh F^{s,1}=\{0\} .
\]

\emph{第四步（结论）}。由第二步与第三步，$\wh F^{s,1}=\{0\}$ 且 $\wh F^{s,2}=\wh F^s$。因此 $\wh F^{s,1}=\wh F^s$ 当且仅当 $\wh F^s=\{0\}$，亦即当且仅当 $\int_\Om d(y)^{-(s+a)p}\,dy=\infty$；当该积分有限时，$\wh F^s=\R\cdot 1$ 而 $\wh F^{s,1}=\{0\}$，故强化不成立。至于 $\wt F$ 空间，第一步给出
\[
  \wt F^{s,1}=\wt F^{s,2}=\wt F^s=\{0\}
\]
在情形 (II) 下无条件成立。
\end{proof}

\begin{remark}[直觉]
在权不可积的情形下，内在空间只剩下常数；而 $L^q$ 的 Hardy 项对应指数 $(s+a)q+n>aq$，奇性太强，连常数都被消灭，于是整个 $\wt F$ 空间退化为 $\{0\}$。$\wh F$ 空间只含较弱的 $L^p$ Hardy 指数项，因此当该项可积时才保住常数函数。
\end{remark}

%==================================================================
\subsection{反例：\texorpdfstring{$C^\infty$}{C-infinity} 稠密而 \texorpdfstring{$C_c^\infty$}{Cc-infinity} 不稠密}

\begin{theorem}[反例：$C^\infty$ 稠密而 $C_c^\infty$ 不稠密]
取 $n=1$、$\Om=(0,1)$、$d(x)=\min(x,1-x)$、$s=1/2$、$p=q=2$ 与 $a=1$。则
\[
  \nuA(\Om)=\int_0^1 d(y)^{-2}\,dy=\infty .
\]
于是由附录 A（取 $k=0$ 与 $\sigma=1/2$）可知，空间 $F^{1/2}_{2,2}(\Om,d^{-2})$ 等于 $P_0(\Om)$，即 $\Om$ 上的常数函数。因此 $C^\infty(\Om)\cap F^{1/2}_{2,2}(\Om,d^{-2})$ 等于 $F^{1/2}_{2,2}(\Om,d^{-2})$，并在其中稠密，而
\[
  C_c^\infty(\Om)\cap F^{1/2}_{2,2}(\Om,d^{-2})=\{0\} .
\]
由于 $F^{1/2}_{2,2}(\Om,d^{-2})$ 非零，$C_c^\infty(\Om)$ 在 $F^{1/2}_{2,2}(\Om,d^{-2})$ 中不稠密；等价地
\[
  F^{1/2,1}_{2,2}(\Om,d^{-2})=\{0\}
  \quad\text{while}\quad
  F^{1/2,2}_{2,2}(\Om,d^{-2})=F^{1/2}_{2,2}(\Om,d^{-2}) .
\]
因此存在 $(\Om,s,p,q,a)$ 使得 $C^\infty$ 稠密而 $C_c^\infty$ 不稠密；修正的稠密性（两个闭包之比较）严格弱于经典的 $C_c^\infty$ 稠密性，且在一般情况下不能加强为 $F^{s,1}=F^s$。
\end{theorem}

\begin{proof}
对 $\Om=(0,1)$，有 $d(y)=\min(y,1-y)$，故当 $y\in(0,1/2)$ 时 $d(y)=y$，且
\[
  \int_0^{1/2} y^{-2}\,dy=\infty .
\]
因此
\[
  \nuA(\Om)=\int_\Om d(y)^{-2}\,dy=\infty .
\]

\emph{第一步（验证退化定理的假设）}。附录 A 的假设在 $n=1$、$\Om=(0,1)$、$a=1$、$q=2$、$s=1/2$、$k=0$、$\sigma=1/2$、$p=2$ 时得到满足。该事实给出
\[
  F^{1/2}_{2,2}(\Om,d^{-2})=P_0(\Om)=\R\cdot 1 ,
\]
即常数函数。

\emph{第二步（两个闭包）}。该事实的 (a) 部分给出
\[
  C^\infty(\Om)\cap F^{1/2}_{2,2}(\Om,d^{-2})=F^{1/2}_{2,2}(\Om,d^{-2}) ,
\]
故闭包 $F^{1/2,2}_{2,2}(\Om,d^{-2})$ 等于 $F^{1/2}_{2,2}(\Om,d^{-2})$。(b) 部分给出
\[
  C_c^\infty(\Om)\cap F^{1/2}_{2,2}(\Om,d^{-2})=\{0\} ,
\]
故闭包 $F^{1/2,1}_{2,2}(\Om,d^{-2})$ 为 $\{0\}$。

\emph{第三步（结论：$C_c^\infty$ 不稠密）}。常数函数 $1$ 属于 $F^{1/2}_{2,2}(\Om,d^{-2})$ 且不为零元，故 $\{0\}$ 是 $F^{1/2}_{2,2}(\Om,d^{-2})$ 的真子集；因此 $C_c^\infty(\Om)$ 在 $F^{1/2}_{2,2}(\Om,d^{-2})$ 中不稠密。这是 $C^\infty$ 稠密成立而 $C_c^\infty$ 稠密失效的一个具体实例。
\end{proof}

\begin{remark}[直觉]
在区间 $(0,1)$ 上权 $d^{-2}$ 不可积，因此分数阶半范数只能对最高阶导数取常值的函数有限；这里 $k=0$，于是只有常数能够存活。常数是光滑的，所以 $C^\infty$ 稠密；但每个紧支撑光滑函数都在边界附近消失，从而具有无穷半范数，于是 $C_c^\infty$ 被压缩成 $\{0\}$。
\end{remark}


%=====================================================================
\section{临界条件：边界余维数，而非权可积性}\label{sec:part5}
%=====================================================================
% ===================================================================
%  第 5 组：临界条件——边界余维数，而非权可积性
% ===================================================================

本节把前面关于稠密性的正面结论反过来使用：在 $\Om=(0,1)$ 这类最简单的有界区域上，权 $d^{-aq}$ 本身完全可积（$aq=1/2<1$，即情形 (I)），$C^\infty(\Om)$ 也确实稠密，然而 $C_c^\infty(\Om)$ 并不稠密。决定性的原因是 Hardy 指数 $sp+aq=3/2$ 超过了边界 $\partial\Om=\{0,1\}$ 的 Assouad 余维数 $1$，而不是权自身的可积性。随后的两条事实把这一定性观察提升为一般原理：$C_c^\infty$ 的闭包恰由 Hardy 条件刻画（Kijaczko 定理 4），临界比较发生在 $sp+b$ 与 $\codimA(\partial\Om)$ 之间而非 $b$ 与 $1$ 之间；最后一条事实则澄清 CLWZ 范数中带权的变量是内层变量 $y$，其指数是 $q$，因此可积性条件用的是 $aq$ 而不是 $ap$。

%==================================================================
\subsection{可积权情形 (I) 的反例：\texorpdfstring{$C^\infty$}{C-infinity} 稠密而 \texorpdfstring{$C_c^\infty$}{Cc-infinity} 不稠密}

\begin{theorem}[第二个反例：在可积权情形 (I) 中 $C^\infty$ 稠密而 $C_c^\infty$ 不稠密]
设 $\Om:=(0,1)$（$\R=\R^1$ 的开子集），并对 $x\in\Om$ 设 $d(x):=\min(x,1-x)=\dist(x,\{0,1\})$。设 $p:=q:=2$，$s:=1/2$，$a:=1/4$。设 $\Fsp$ 为 CLWZ 的加权内蕴 Triebel--Lizorkin 空间，即公式 (1.2)（采用 CLWZ 第 3 页的加权 Gagliardo 范数，它与该文记号下的 $W^{s,p}(\Om,d^{-ap})$ 一致）；由于 $p=q=2$，其范数为
\[
  \|f\|_{F^s}:=\|f\|_{L^2(\Om)}
  +\Big(\int_\Om\int_\Om|f(x)-f(y)|^2\,|x-y|^{-1-2s}\,d(y)^{-aq}\,dy\,dx\Big)^{1/2}
\]
上式对 $\Om$ 上一切使右端两项均有限的可测函数 $f$ 成立。设 $F^{s,1}_{p,q}(\Om,d^{-aq})$ 为 $C_c^\infty(\Om)\cap\Fsp$ 在 $\Fsp$ 中的闭包，$F^{s,2}_{p,q}(\Om,d^{-aq})$ 为 $C^\infty(\Om)\cap F^s$ 在 $F^s$ 中的闭包。则：
\begin{enumerate}[label=(\roman*)]
  \item $\displaystyle\int_\Om d(y)^{-aq}\,dy=\int_0^1 d(y)^{-1/2}\,dy=2\sqrt2<\infty$；从而对边界距离权而言，$(\Om,s,p,q,a)$ 落在二分法的情形 (I) 中。
  \item 常函数 $1$（即 $1_\Om$）属于 $\Fsp$，但 $1$ 不属于 $F^{s,1}_{p,q}(\Om,d^{-aq})$。因此 $F^{s,1}_{p,q}(\Om,d^{-aq})$ 是 $\Fsp$ 的真闭向量子空间，且 $C_c^\infty(\Om)\cap\Fsp$ 在 $\Fsp$ 中不稠密。
  \item $C^\infty(\Om)\cap\Fsp$ 在 $\Fsp$ 中稠密；即 $F^{s,2}_{p,q}(\Om,d^{-aq})=\Fsp$。
\end{enumerate}
因此对这些数据有 $F^{s,1}\neq F^s=F^{s,2}$：$C^\infty(\Om)$ 在 $\Fsp$ 中稠密，而 $C_c^\infty(\Om)$ 不稠密，且这已在可积权情形 (I) 中发生（事实上 $aq=1/2<1$，只是 $(s+a)q=3/2>1=\codimA(\partial\Om)$）。所以该现象并不局限于二分法的退化情形 (II)，并且在情形 (I) 中修正后的稠密性质 $F^{s,1}=F^{s,2}$ 也严格弱于经典稠密性质 $F^{s,1}=F^s$。
\end{theorem}

\begin{proof}
记 $d(x)=\min(x,1-x)$，于是 $d(x)=\dist(x,\partial\Om)$，其中 $\partial\Om=\{0,1\}$；并令
\[
  [f]:=\Big(\int_\Om\int_\Om|f(x)-f(y)|^2\,|x-y|^{-2}\,d(y)^{-1/2}\,dy\,dx\Big)^{1/2},
\]
\[
  F^s:=\big\{f:\Om\to\R\ \text{Lebesgue measurable}:\ \|f\|_{L^2(\Om)}+[f]<\infty\big\}.
\]
注意由 $s=1/2$ 得 $|x-y|^{-1-2s}=|x-y|^{-2}$，由 $a=1/4$ 与 $q=2$ 得 $d(y)^{-aq}=d(y)^{-1/2}$。

\emph{第一步（与 Kijaczko 空间的等同）}. 设 $W:=W^{s,p;\alpha,\beta}(\Om)$ 为 Kijaczko 的加权分数次 Sobolev 空间（arXiv:2110.04012），取 $d:=1$、$s:=1/2$、$p:=2$、$\alpha:=1/2$、$\beta:=0$。按 Kijaczko 的定义（第 1 节定义 1 及记号 $W^{s,p;w_\alpha,v_\beta}=:W^{s,p;\alpha,\beta}$），$W^{s,p;\alpha,\beta}(\Om)$ 的半范数为
\[
  [f]_{W}=\Big(\int_\Om\int_\Om|f(x)-f(y)|^p\,|x-y|^{-d-sp}\,w_\alpha(y)\,v_\beta(x)\,dy\,dx\Big)^{1/p}
\]
其中 $w_\alpha(y)=d(y)^{-\alpha}=d(y)^{-1/2}$，$v_\beta(x)=d(x)^{-\beta}=1$，即 $[f]_{W}=[f]$；$W$ 的范数为 $\|f\|_{L^p(\Om)}+[f]_{W}=\|f\|_{L^2(\Om)}+[f]$。因此作为赋范空间有 $F^s=W^{1/2,2;1/2,0}(\Om)$。按 Kijaczko 的定义 1，$W_0^{1/2,2;1/2,0}(\Om)$ 是 $C_c^\infty(\Om)\cap W^{1/2,2;1/2,0}(\Om)$ 在该范数下的闭包；由于 $C_c^\infty(\Om)\cap F^s=C_c^\infty(\Om)\cap W^{1/2,2;1/2,0}(\Om)$，该闭包恰为 $F^{s,1}$。于是
\begin{equation}\label{eq:sec5f1-ident}
  F^{s,1}_{p,q}(\Om,d^{-aq})=W_0^{1/2,2;1/2,0}(\Om).
\end{equation}

\emph{第二步（$\Om$ 是有界一致域）}. $\Om=(0,1)$ 非空、有界且开。对 $\Om$ 中的 $x<y$，取 $\Gamma:=[x,y]$ 为线段；则 $L(\Gamma)=|x-y|$，故 $c_1:=1$ 可用。设 $z\in\Gamma$。若 $|x-z|\le|y-z|$，则由 $x>0$ 得 $|x-z|=z-x\le z$，且由 $y<1$ 得 $|x-z|\le|y-z|=y-z\le1-z$；故 $|x-z|\le\min(z,1-z)=d(z)$。若 $|y-z|<|x-z|$，则交换 $x$ 与 $y$ 的角色作同样的计算，得 $|y-z|\le\min(z,1-z)=d(z)$。因此对每个 $z\in\Gamma$ 有 $\min\{|x-z|,|y-z|\}\le d(z)$，故 $c_2:=1$ 可用。所以 $\Om$ 是 Kijaczko 第 2.3 节（以及 CLWZ 第 3 页）意义下的一致域。

\emph{第三步（边界的 Assouad 余维数）}. 设 $E:=\partial\Om=\{0,1\}$，并对 $r>0$ 设 $\wt E_r:=\{t\in\R:\dist(t,E)\le r\}=[-r,r]\cup[1-r,1+r]$ 为 $E$ 的管状邻域。这里 $\diam(E)=1$。回忆（Kijaczko 定义 9）：$\codimAl(E)$ 是满足下述条件的 $q\ge0$ 的上确界：存在 $C=C(q)\ge1$，使得
\[
  |\wt E_r\cap B(x,R)|\le C\,|B(x,R)|\,(r/R)^q
\]
对所有 $x\in E$ 与 $0<r<R<\diam(E)$ 成立；$\codimAu(E)$ 是满足下述条件的 $s\ge0$ 的下确界：存在 $c=c(s)>0$，使得
\[
  |\wt E_r\cap B(x,R)|\ge c\,|B(x,R)|\,(r/R)^s
\]
对所有 $x\in E$ 与 $0<r<R<\diam(E)$ 成立。在 $\R^1$ 中 $|B(x,R)|=2R$。

\emph{$\codimAu(E)$ 的计算}. 由于 $\wt E_r$ 包含 $[-r,r]$，对 $x=0$ 与 $0<r<R<1$ 有
\[
  |\wt E_r\cap B(0,R)|\ge|[-r,r]|=2r=2R\,(r/R)=|B(0,R)|\,(r/R);
\]
由对称性，对 $x=1$ 同样成立。故 $c(1):=1$ 与 $s=1$ 可用，从而 $\codimAu(E)\le1$。反之，设 $0\le s<1$，且某个 $c>0$ 对所有 $x\in E$ 与 $0<r<R<1$ 满足 $|\wt E_r\cap B(x,R)|\ge c|B(x,R)|(r/R)^s$。取 $x=0$。由于 $\wt E_r=[-r,r]\cup[1-r,1+r]$，有 $|\wt E_r\cap B(0,R)|\le|\wt E_r|\le4r$。于是所需不等式将迫使
\[
  4r\ge c\cdot2R\,(r/R)^s=2c\,r^sR^{1-s},
\]
即 $2\ge c\,r^{s-1}R^{1-s}=c\,(R/r)^{1-s}$。因 $1-s>0$，当 $R/r$ 趋于 $+\infty$ 时右端 $c(R/r)^{1-s}$ 趋于 $+\infty$，故当 $R/r$ 较大时该不等式不成立；由于 $0<r<R<1$ 任意，这样的 $c$ 不存在。因此 $\codimAu(E)=1$。

\emph{$\codimAl(E)$ 的计算}. 对 $x\in E$ 与 $0<r<R<1$，集合 $\wt E_r\cap B(x,R)$ 包含于 $B(x,r)$ 与 $B(x',r)$ 的并，其中 $x'$ 为 $E$ 的另一点，故
\[
  |\wt E_r\cap B(x,R)|\le2\cdot2r=4r=2\cdot2R\,(r/R)=2\,|B(x,R)|\,(r/R);
\]
故 $C(1):=2$ 与 $q=1$ 可用，且 $\codimAl(E)\ge1$。反之，设 $q>1$，且某个 $C\ge1$ 对所有 $x\in E$ 与 $0<r<R<1$ 满足 $|\wt E_r\cap B(x,R)|\le C|B(x,R)|(r/R)^q$。取 $x=0$。由于 $\wt E_r$ 包含 $[-r,r]$，有 $|\wt E_r\cap B(0,R)|\ge2r$，于是所需不等式将迫使
\[
  2r\le C\cdot2R\,(r/R)^q=2C\,r^qR^{1-q},
\]
即 $1\le C\,r^{q-1}R^{1-q}=C\,(r/R)^{q-1}$。因 $q-1>0$，当 $r/R$ 趋于 $0$ 时右端 $C(r/R)^{q-1}$ 趋于 $0$，故当 $r/R$ 较小时该不等式不成立；由于 $0<r<R<1$ 任意，这样的 $C$ 不存在。因此 $\codimAl(E)=1$。于是
\begin{equation}\label{eq:sec5f1-codim}
  \codimAl(\partial\Om)=\codimAu(\partial\Om)=1.
\end{equation}

\emph{第四步（Kijaczko 定理 4）}. 我们使用 Kijaczko（arXiv:2110.04012）定理 4，其陈述为：设 $\Om$ 为 $\R^d$ 中非空、有界、一致的开集，$0<s<1$、$1\le p<\infty$ 且 $0\le\alpha,\beta<\codimAl(\partial\Om)$；若 $sp+\alpha+\beta>\codimAu(\partial\Om)$，则
\[
  W_0^{s,p;\alpha,\beta}(\Om)
  =\Big\{f\in W^{s,p;\alpha,\beta}(\Om):\int_\Om|f(x)|^p d(x)^{-(sp+\alpha+\beta)}\,dx<\infty\Big\}.
\]
取 $d=1$、$s=1/2$、$p=2$、$\alpha=1/2$、$\beta=0$ 时假设成立：$0<1/2<1$；$1\le2<\infty$；由 \eqref{eq:sec5f1-codim} 得 $0\le\alpha=1/2<1$ 与 $0\le\beta=0<1$；由第二步，$\Om$ 有界、非空、开且一致；并由 \eqref{eq:sec5f1-codim} 得 $sp+\alpha+\beta=(1/2)\cdot2+1/2+0=3/2>1=\codimAu(\partial\Om)$。结论给出
\[
  W_0^{1/2,2;1/2,0}(\Om)=\Big\{f\in W^{1/2,2;1/2,0}(\Om):\int_\Om|f(x)|^2d(x)^{-3/2}\,dx<\infty\Big\},
\]
即由 \eqref{eq:sec5f1-ident} 得
\begin{equation}\label{eq:sec5f1-closure}
  F^{s,1}=\Big\{f\in F^s:\int_\Om|f(x)|^2d(x)^{-3/2}\,dx<\infty\Big\}.
\end{equation}

\emph{第五步（常数函数 $1$）}. 函数 $1$ 属于 $L^2(\Om)$ 且 $\|1\|_{L^2}=1$；因被积函数 $|1-1|^2|x-y|^{-2}d(y)^{-1/2}$ 恒为零，故 $[1]=0$；于是 $1\in F^s$。另一方面，对 $x\in(0,1/2]$ 有 $d(x)^{-3/2}\ge x^{-3/2}$，且对 $0<\varepsilon<1/2$ 有
\[
  \int_\varepsilon^{1/2}x^{-3/2}\,dx=\Big[-2x^{-1/2}\Big]_\varepsilon^{1/2}=2\big(\varepsilon^{-1/2}-\sqrt2\big),
\]
该式当 $\varepsilon\to0$ 时趋于 $+\infty$。因此 $\int_\Om d(x)^{-3/2}\,dx=+\infty$，故由 \eqref{eq:sec5f1-closure}，常函数 $1$ 不属于 $F^{s,1}$。

\emph{第六步（$F^{s,1}$ 是非零真闭子空间）}. 按定义 $F^{s,1}$ 是 $F^s$ 的闭向量子空间；由第五步它不含 $1$，而 $1\in F^s$；因此 $F^{s,1}$ 是 $F^s$ 的真子空间，且 $C_c^\infty(\Om)\cap F^s$ 在 $F^s$ 中不稠密，这正是 (ii)。它非零：若 $g\in C_c^\infty(\Om)$，则 $g'$ 有界且 $\|g'\|_{L^\infty(\Om)}<\infty$；又因 $\Om$ 凸，连接 $x$ 与 $y$ 的线段落在 $\Om$ 内，故由中值定理得 $|g(x)-g(y)|\le\|g'\|_{L^\infty(\Om)}|x-y|$；因此
\[
  [g]^2\le\|g'\|_{L^\infty(\Om)}^2\int_\Om\int_\Om d(y)^{-1/2}\,dy\,dx
  =\|g'\|_{L^\infty(\Om)}^2\,|\Om|\cdot2\sqrt2<\infty,
\]
并且 $\|g\|_{L^2(\Om)}\le\|g\|_{L^\infty(\Om)}<\infty$，故 $g\in C_c^\infty(\Om)\cap F^s$。因此 $C_c^\infty(\Om)\cap F^s$ 是 $F^s$ 的非零子集，其闭包 $F^{s,1}$ 为真子空间。

\emph{第七步（第 (iii) 部分）}. 我们对已核实的附录 A（情形 (I) 中 $p=q$、$s\in(0,1)$ 的 $C^\infty$ 稠密性）取 $n:=1$、$\Om:=(0,1)$、$a:=1/4$、$s:=1/2$、$p:=2$。其假设为 $\int_\Om d(y)^{-ap}\,dy<\infty$，而确实 $ap=(1/4)\cdot2=1/2$，故 $\int_\Om d(y)^{-1/2}\,dy=2\sqrt2<\infty$，与 (i) 中所算一致（在 $(0,1)$ 上，$\int_0^1 d(y)^{-1/2}\,dy=2\int_0^{1/2}y^{-1/2}\,dy=2\cdot2\cdot(1/2)^{1/2}=2\sqrt2$）。因此 $C^\infty(\Om)\cap F^s$ 在 $F^s$ 中稠密，即 $F^{s,2}=F^s$，这正是 (iii)。

最后，(i) 即上面刚作的计算：$\int_\Om d(y)^{-aq}\,dy=\int_0^1 d(y)^{-1/2}\,dy=2\sqrt2<\infty$，故由二分法的附录 A，当前数据落在情形 (I)。结合 (ii) 与 (iii) 得 $F^{s,1}\neq F^s=F^{s,2}$。
\end{proof}

\begin{remark}[直觉]
权的可积性 $aq=1/2<1$ 只保证 $C^\infty$ 稠密；$C_c^\infty$ 的闭包由 Kijaczko 定理 4 精确刻画为 $F^s$ 中满足 Hardy 条件 $\int|f|^2d^{-3/2}<\infty$ 的子空间。这里 Hardy 指数 $sp+aq=3/2$ 大于边界的 Assouad 余维数 $1$，于是该子空间是真子空间；加权半范数在常数上消失的常数函数 $1$ 恰好违反这一 Hardy 条件，因而落在闭包之外。同一个例子于是同时展示了情形 (I) 中 $C^\infty$ 稠密与 $C_c^\infty$ 不稠密并存。
\end{remark}

%==================================================================
\subsection{区间例子：闭包的精确刻画与 Hardy 指数 \texorpdfstring{$sp+aq$}{sp+aq} 的临界性}

\begin{theorem}[区间例子：在可积权情形 (I) 中 $C^\infty$ 稠密而 $C_c^\infty$ 不稠密]
设 $\Om:=(0,1)\subset\R$，对 $x\in\Om$ 设 $d(x):=\min(x,1-x)=\dist(x,\{0,1\})$，并设 $p:=q:=2$，$s:=1/2$，$a:=1/4$。设 $F^s:=F^s_{2,2}(\Om,d^{-aq})$ 为 CLWZ (1.1)--(1.2) 的加权内蕴 Triebel--Lizorkin 空间，即 $\Om$ 上使
\[
  \|f\|_{F^s}:=\|f\|_{L^2(\Om)}
  +\Big(\int_\Om\int_\Om|f(x)-f(y)|^2\,|x-y|^{-2}\,d(y)^{-1/2}\,dy\,dx\Big)^{1/2}
\]
有限的可测函数 $f$ 之集。设 $F^{s,1}$ 为 $C_c^\infty(\Om)\cap F^s$ 在 $F^s$ 中的闭包，$F^{s,2}$ 为 $C^\infty(\Om)\cap F^s$ 在 $F^s$ 中的闭包。则：
\begin{enumerate}[label=(\roman*)]
  \item $\displaystyle\int_\Om d(y)^{-1/2}\,dy=2\sqrt2<\infty$，故对权 $d^{-aq}$ 而言 $(\Om,s,p,q,a)$ 属于情形 (I)（可积权）。
  \item 常函数 $1$ 属于 $F^s$ 且 $\|1\|_{F^s}=1$，但 $1$ 不属于 $F^{s,1}$。更精确地，
  \[
    F^{s,1}=\Big\{f\in F^s:\int_\Om|f(x)|^2d(x)^{-3/2}\,dx<\infty\Big\}.
  \]
  \item $C^\infty(\Om)\cap F^s$ 在 $F^s$ 中稠密，即 $F^{s,2}=F^s$。
  \item 因此 $F^{s,1}\neq F^s=F^{s,2}$：$C^\infty(\Om)$ 在 $F^s$ 中稠密而 $C_c^\infty(\Om)$ 不稠密，尽管权 $d^{-1/2}$ 可积（$aq=1/2<1$）。在此例中，$C_c^\infty$ 的不稠密性是由 Hardy 指数 $sp+aq=3/2$ 超过边界余维数 $1$ 造成的，而不是由可积性指数 $aq=1/2$ 造成的。
\end{enumerate}
\end{theorem}

\begin{proof}
\emph{第零步（等同与外部结果）}. 由于 $p=q=2$ 且 $s=1/2\in(0,1)$，陈述中所写的范数恰是 Kijaczko 加权分数次 Sobolev 空间 $W^{s,p;\alpha,\beta}(\Om)$ 在 $p=2$、$s=1/2$、$\alpha=1/2$（内层变量 $y$ 上的权为 $d(y)^{-1/2}$）、$\beta=0$（$x$ 上无权）时的范数 $\|f\|_{L^2}+[f]$；见 arXiv:2110.04012 第 1 节，其中
\[
  [f]_{W^{s,p;\alpha,\beta}}=\Big(\int_\Om\int_\Om|f(x)-f(y)|^p\,|x-y|^{-d-sp}\,d(y)^{-\alpha}\,d(x)^{-\beta}\,dy\,dx\Big)^{1/p}.
\]
故 $F^s=W^{1/2,2;1/2,0}(\Om)$，且范数相等。

我们使用三条外部结果，并完整引用。以下对 $\R^d$ 中的集合 $E$，$\codimAl(E)$ 与 $\codimAu(E)$ 分别表示 $E$ 的下 Assouad 余维数与上 Assouad 余维数（Kijaczko，arXiv:2110.04012，定义 9）。

(A) Kijaczko，arXiv:2110.04012，定理 4。设 $\Om$ 为 $\R^d$ 中非空、有界、一致的开集，$0<s<1$、$1\le p<\infty$，并设 $\alpha,\beta$ 为满足 $0\le\alpha,\beta<\codimAl(\partial\Om)$ 的实数。若 $sp+\alpha+\beta>\codimAu(\partial\Om)$，则
\[
  W_0^{s,p;\alpha,\beta}(\Om)=\Big\{f\in W^{s,p;\alpha,\beta}(\Om):\int_\Om|f(x)|^p d_\Om(x)^{-sp-\alpha-\beta}\,dx<\infty\Big\},
\]
其中 $d_\Om(x)=\dist(x,\partial\Om)$，$W_0^{s,p;\alpha,\beta}(\Om)$ 是 $C_c^\infty(\Om)\cap W^{s,p;\alpha,\beta}(\Om)$ 在 $W^{s,p;\alpha,\beta}(\Om)$ 中关于范数 $\|f\|_{L^p}+[f]_{W^{s,p;\alpha,\beta}}$ 的闭包（Kijaczko 定义 1）。

(B) Kijaczko，arXiv:2110.04012，定理 19。设 $\Om$ 为 $\R^d$ 中非空开集，$0<s<1$ 且 $p\in[1,\infty)$，并设 $w,v$ 为非负权，它们局部有界且满足
\[
  \int_\Om w(x)(1+|x|)^{-d-sp}\,dx<\infty
  \quad\text{与}\quad
  \int_\Om v(x)(1+|x|)^{-d-sp}\,dx<\infty
\]
（条件 (6)）。则 $C^\infty(\Om)\cap\wt W^{s,p;w,v}(\Om)$ 在半赋范空间 $\wt W^{s,p;w,v}(\Om)$ 中稠密，其中
\[
  \wt W^{s,p;w,v}(\Om)=\Big\{f\ \text{measurable on}\ \Om:\int_\Om\int_\Om|f(x)-f(y)|^p|x-y|^{-d-sp}w(x)v(y)\,dx\,dy<\infty\Big\}
\]
并赋以相应的半范数；此外 $W_0^{s,p;w,v}(\Om)=W_c^{s,p;w,v}(\Om)$，其中 $W_c^{s,p;w,v}(\Om)$ 是 $\Om$ 上紧支可测函数在 $\wt W^{s,p;w,v}(\Om)$ 中的闭包。

(C) Dyda--Kijaczko，arXiv:2009.00077，定理 3。设 $p\in[1,\infty)$ 且 $f\in L^p(\Om)$。若对 $\Om$ 的某个 Whitney 分解 $\{Q_n\}$ 的每个 Whitney 立方体 $Q$ 有 $\eta_k(Q)\to0$，则 $\lim_{k\to\infty}\|P^{\eta_k}f-f\|_{L^p(\Om)}=0$，其中 $P^\eta$ 是该文的 Whitney 局部磨光算子
\[
  P^\eta f=\sum_n (f\psi_n)*h_{\eta(Q_n)}
\]
（$\psi_n$ 为从属于 $\{Q_n\}$ 的单位分解，$h$ 为固定磨光核，$h_t(x)=t^{-d}h(x/t)$）。

\emph{第一步（情形 (I)）}. 由于在 $(0,1/2)$ 上 $d(y)=y$，在 $(1/2,1)$ 上 $d(y)=1-y$，故
\[
  \int_\Om d(y)^{-1/2}\,dy=2\int_0^{1/2}y^{-1/2}\,dy=2\big[2y^{1/2}\big]_0^{1/2}=4\sqrt{1/2}=2\sqrt2<\infty.
\]
于是取 $a=1/4$、$q=2$ 时 $\nuA(\Om)=\int_\Om d^{-aq}=\int_\Om d^{-1/2}=2\sqrt2<\infty$，即二分法的情形 (I) 成立（二分法见附录 A）。

\emph{第二步（边界的余维数等于 $1$）}. 这里 $\partial\Om=\{0,1\}$。对 $\R$ 中的集合 $E$，记 $\wt E_r=\{t\in\R:\dist(t,E)\le r\}$ 为其管状邻域。对 $E=\{0,1\}$ 有 $\wt E_r=[-r,r]\cup[1-r,1+r]$。固定 $x\in E$ 与 $0<r<R<\diam(E)=1$。若 $x=0$，则 $B(x,R)=(-R,R)$，而 $\wt E_r\cap B(x,R)$ 包含 $(-r,r)$，故其 Lebesgue 测度至少为 $2r$；又 $\wt E_r\cap B(x,R)$ 包含于 $[-r,r]\cup[1-r,R)$，因 $R<1$ 其测度至多为 $2r+(R-(1-r))_+\le2r+r=3r$。因此 $2r\le|\wt E_r\cap B(x,R)|\le3r$。由于 $|B(x,R)|=2R$，比值 $|\wt E_r\cap B(x,R)|/|B(x,R)|$ 介于 $(r/R)$ 与 $(3/2)(r/R)$ 之间。由对称性，$x=1$ 时同样成立。

\emph{下余维数}. 由刚证的上界，取 $q=1$、$C=3/2$ 即可用。若某个 $q>1$ 与某个 $C=C(q)$ 可容许，则对所有 $x\in E$ 与 $0<r<R<1$ 将有
\[
  2r\le|\wt E_r\cap B(x,R)|\le C\,|B(x,R)|\,(r/R)^q=2CR\,(r/R)^q=2C\,r^qR^{1-q}.
\]
除以 $2r>0$ 得 $1\le C\,r^{q-1}R^{1-q}=C\,(r/R)^{q-1}$。固定 $R=1/2$，则右端等于 $C(2r)^{q-1}$，因 $q-1>0$，当 $r\to0$ 时它趋于 $0$。故对所有充分小的 $r>0$ 该不等式不成立，矛盾。因此没有 $q>1$ 可容许，$\codimAl(\partial\Om)=1$。

\emph{上余维数}. 使存在 $c=c(s)>0$、对所有 $x\in E$ 与 $0<r<R<1$ 满足 $|\wt E_r\cap B(x,R)|\ge c|B(x,R)|(r/R)^s$ 的 $s\ge0$ 之集包含 $s=1$，这由下界 $2r\ge2R(r/R)/2=|B(x,R)|(r/R)/2$（取 $c=1/2$）可得。若 $s<1$，则 $|\wt E_r\cap B(x,R)|\le3r$，而 $c|B(x,R)|(r/R)^s=2cR(r/R)^s=2c\,r^sR^{1-s}$；令 $R\to1$、$r\to0$ 得 $3r\ge c'\,r^s$，因 $s<1$ 该式对小 $r$ 不成立。因此 $\codimAu(\partial\Om)=1$。

区间 $\Om=(0,1)$ 是有界一致域：其中任意两点 $x<y$ 由线段 $[x,y]$ 相连，其长度为 $|x-y|$（故取 $c_1=1$ 有 $L(\Gamma)\le c_1|x-y|$）；且对 $z\in[x,y]$ 有 $\min\{|x-z|,|y-z|\}\le\min(z-x,y-z)\le\min(z,1-z)=d(z)$（$z\le1/2$ 时用 $\min\{|x-z|,|y-z|\}\le z-x\le z$，$z\ge1/2$ 时用 $\le y-z\le1-z$），故一致域条件以 $c_1=c_2=1$ 成立。

\emph{第三步（闭包的刻画）}. 对定理 (A) 取 $d=1$、$s=1/2$、$p=2$、$\alpha=1/2$、$\beta=0$。假设成立：$\Om$ 有界且一致（第二步）；$0\le\alpha=1/2<1=\codimAl(\partial\Om)$ 且 $0\le\beta=0<1$；以及 $sp+\alpha+\beta=(1/2)(2)+1/2+0=3/2>1=\codimAu(\partial\Om)$。由于 $W^{1/2,2;1/2,0}(\Om)=F^s$ 且范数相等（第零步），并 $W_0^{1/2,2;1/2,0}(\Om)=F^{s,1}$，定理 (A) 给出
\[
  F^{s,1}=\Big\{f\in F^s:\int_\Om|f(x)|^2d(x)^{-(sp+\alpha+\beta)}\,dx<\infty\Big\}
  =\Big\{f\in F^s:\int_\Om|f(x)|^2d(x)^{-3/2}\,dx<\infty\Big\}.
\]
这就是所断言的 $F^{s,1}$ 的精确刻画。注意指数是 $sp+\alpha+\beta=3/2$，而不是权的可积性指数 $aq=1/2$。

\emph{第四步（$1$ 在 $F^s$ 中但不在 $F^{s,1}$ 中）}. 常函数 $1$ 满足 $\|1\|_{L^2(\Om)}=|\Om|^{1/2}=1$，其加权 Gagliardo 半范数为零，因为对所有 $x,y\in\Om$ 有 $|1(x)-1(y)|=0$。故 $1\in F^s$ 且 $\|1\|_{F^s}=1$。另一方面
\[
  \int_\Om|1(x)|^2d(x)^{-3/2}\,dx=\int_\Om d(x)^{-3/2}\,dx=2\int_0^{1/2}y^{-3/2}\,dy=\infty;
\]
等价地，取 $c=3/2\ge1$，附录 A（对有界 Lipschitz 域，对每个实数 $c\ge1$ 有 $\int_\Om d^{-c}=\infty$）给出该积分发散。由第三步的刻画，$1$ 不在 $F^{s,1}$ 中。因此 $F^{s,1}$ 是 $F^s$ 的真闭子空间，且 $C_c^\infty(\Om)\cap F^s$ 在 $F^s$ 中不稠密。

\emph{第五步（$C^\infty$ 稠密性）}. 附录 A 的假设均满足：$n=1$，$\Om=(0,1)$ 为非空有界域，$a=1/4>0$，$s=1/2\in(0,1)$，$p=2\in[1,\infty)$，且 $\int_\Om d(y)^{-ap}\,dy=\int_\Om d(y)^{-1/2}\,dy=2\sqrt2<\infty$（第一步）。该事实（其证明恰如上文定理 (B) 与 (C) 那样引用 Kijaczko 定理 19 与 Dyda--Kijaczko 定理 3）给出 $C^\infty(\Om)\cap F^s$ 在 $F^s$ 中稠密，即 $F^{s,2}=F^s$。为完整起见，其论证为：取 $w:=1$、$v:=d^{-1/2}$，两权在 $\Om$ 上均局部有界，而因 $\Om$ 有界，定理 (B) 的条件 (6) 化为 $w\in L^1(\Om)$ 与 $v\in L^1(\Om)$；前者由 $|\Om|<\infty$ 成立，后者即第一步。于是定理 (B) 适用，其 Whitney 磨光逼近列 $P^{\eta_k}f$（满足 $\eta_k(Q_n)\to0$）在加权 Gagliardo 半范数下收敛于 $f$；由定理 (C) 同一列在 $L^2(\Om)$ 中收敛于 $f$；故它在完整范数 $\|f\|_{F^s}=\|f\|_{L^2}+[f]$ 下收敛，且每个 $P^{\eta_k}f$ 属于 $C^\infty(\Om)\cap F^s$。

\emph{第六步（结论）}. 由第四步，$F^{s,1}$ 是 $F^s$ 的真子空间；由第五步，$F^{s,2}=F^s$。因此 $F^{s,1}\neq F^s=F^{s,2}$。至此 (i)--(iv) 得证。
\end{proof}

\begin{remark}[直觉]
权的可积性（情形 (I)，$aq=1/2<1$）只保证 $C^\infty$ 稠密，并不保证 $C_c^\infty$ 稠密。Kijaczko 定理 4 精确描述了 $C_c^\infty$ 的闭包：它是 $F^s$ 中由 Hardy 条件 $\int|f|^p d^{-(sp+aq)}<\infty$ 切出的子空间。这里 $sp+aq=3/2>\codimA(\partial\Om)=1$，所以它是真子空间，而加权半范数消失的常数函数 $1$ 恰好违反该条件。于是同一个例子同时展示了情形 (I) 中 $C^\infty$ 稠密与 $C_c^\infty$ 不稠密的失败。
\end{remark}

%==================================================================
\subsection{为何 CLWZ 的 Hardy 项不可省，以及临界指数是 \texorpdfstring{$sp+aq$}{sp+aq} 而非 \texorpdfstring{$aq$}{aq}}

\begin{theorem}[为何 CLWZ 的 Hardy 项并非冗余，以及为何临界指数是 $sp+aq$ 而非 $aq$]
设 $n\ge1$，$\Om$ 为 $\R^n$ 中非空有界 Lipschitz 域，设 $d(x)=\dist(x,\partial\Om)$，$0<s<1$、$p\in[1,\infty)$ 且 $b\ge0$。定义
\[
  W_b^{s,p}(\Om):=\Big\{f\in L^p(\Om):
  [f]_{s,p,b}:=\Big(\int_\Om\int_\Om|f(x)-f(y)|^p|x-y|^{-n-sp}d(y)^{-b}\,dy\,dx\Big)^{1/p}<\infty\Big\},
\]
并以 $\|f\|_{W_b^{s,p}}:=\|f\|_{L^p(\Om)}+[f]_{s,p,b}$ 赋范；设 $W_{b,0}^{s,p}(\Om)$ 为 $C_c^\infty(\Om)\cap W_b^{s,p}(\Om)$ 在 $W_b^{s,p}(\Om)$ 中的闭包。设 $\codimAl(\partial\Om)$ 与 $\codimAu(\partial\Om)$ 为 $\partial\Om$ 的下、上 Assouad 余维数，$\dimM(\partial\Om)$ 为其上 Minkowski 维数。则：
\begin{enumerate}[label=(\roman*)]
  \item （闭包刻画，Kijaczko 定理 4）若 $0\le b<\codimAl(\partial\Om)$ 且 $sp+b>\codimAu(\partial\Om)$，则
  \[
    W_{b,0}^{s,p}(\Om)=\Big\{f\in W_b^{s,p}(\Om):\int_\Om|f(x)|^p d(x)^{-(sp+b)}\,dx<\infty\Big\}.
  \]
  在这些假设下该子空间为真子空间：$\Om$ 是 $\kappa$-plump 的（有界 Lipschitz 域是一致域，而一致域是 $\kappa$-plump 的），故 Kijaczko 定理 3(III) 在 $sp+b>\codimAu(\partial\Om)$ 下适用，并给出 $W_{b,0}^{s,p}(\Om)\neq W_b^{s,p}(\Om)$。
  \item （哪个指数是临界的）若 $sp+b<n-\dimM(\partial\Om)$，则 $W_{b,0}^{s,p}(\Om)=W_b^{s,p}(\Om)$（Kijaczko 定理 3(I)）。可积性条件 $\int_\Om d^{-b}<\infty$（在远离临界情形时等价于 $b<n-\dimM(\partial\Om)$）是 Kijaczko $C^\infty$ 稠密性定理（定理 19）的假设，它只保证 $C^\infty(\Om)\cap W_b^{s,p}(\Om)$ 在 $W_b^{s,p}(\Om)$ 中稠密。在子区间 $0\le b<\codimAl(\partial\Om)$ 且 $sp+b>\codimAu(\partial\Om)$ 内（该子区间非空：有界 Lipschitz 域的边界是 $(n-1)$ 维 Ahlfors 正则集，故
$\codimAl(\partial\Om)=\codimAu(\partial\Om)=n-(n-1)=1$；于是所论区间即 $\big(\max\{0,1-sp\},\,1\big)$，而 $s>0$、$p\ge1$ 给出 $sp>0$，故它非空）有 $b<n-\dimM(\partial\Om)<sp+b$，故 $C^\infty$ 稠密而 $C_c^\infty$ 不稠密。因此决定 $C_c^\infty$ 稠密性的指数是 $sp+b$，而不是权的可积性指数 $b$：临界比较发生在 $sp+b$ 与边界余维数之间，而不是 $\int_\Om d^{-b}$ 的有限性。
  \item （区间见证）对 $\Om=(0,1)$、$n=1$、$s=1/2$、$p=2$、$b=1/2$ 有 $\int_\Om d^{-b}=2\sqrt2<\infty$ 且
  \[
    b=1/2<1=n-\dimM(\partial\Om)=\codimAl(\partial\Om)=\codimAu(\partial\Om),
  \]
  但 $sp+b=3/2>1$。故 (i) 适用，并得
  \[
    W_{b,0}^{s,p}(\Om)=\Big\{f\in W_b^{s,p}(\Om):\int_\Om|f(x)|^2d(x)^{-3/2}\,dx<\infty\Big\},
  \]
  该子空间为真子空间且不含常数 $1$。因此情形 (I)（$\int d^{-aq}<\infty$）并不蕴含 $C_c^\infty$ 稠密性。
  \item （CLWZ 的 Hardy 项）设 $a>0$、$q\in[1,\infty)$，并设 $F^s:=\Fsp$ 为 CLWZ 加权内蕴空间 (1.1)--(1.2)，其中 $s\in(0,\infty)\setminus\N$。
  \begin{enumerate}[label=(\alph*)]
    \item 若 $0<s<1$、$p=q$ 且 $b=aq$，则 $sp+b=sp+aq=(s+a)p$；此外在 $0\le aq<\codimAl(\partial\Om)$ 且 $sp+aq>\codimAu(\partial\Om)$ 的情形下，部分 (i) 给出
    \begin{align*}
      F^{s,1}&=W_{b,0}^{s,p}(\Om)
      =\Big\{f\in F^s:\int_\Om|f(x)|^p d(x)^{-(s+a)p}\,dx<\infty\Big\}\\
      &=\Big\{f\in F^s:\|f\,d^{-(s+a)}\|_{L^p(\Om)}<\infty\Big\}.
    \end{align*}
    在此情形下，第一个 CLWZ Hardy 范数恰是定义 $C_c^\infty$ 闭包所缺的那个条件，故它并非冗余。
    \item 不论处于何种情形，若 $(s+a)p\ge1$，则常函数 $1$ 属于 $F^s$，但 $\int_\Om d^{-(s+a)p}\,dx=\infty$（附录 A，取 $c=(s+a)p\ge1$），故 $\|1\cdot d^{-(s+a)}\|_{L^p(\Om)}=\infty$；因此在该范围内第一个 Hardy 项不是内蕴范数的推论。
    \item 对每个 $s\in(0,\infty)\setminus\N$、每个 $a>0$ 与每个 $q\in[1,\infty)$，常数 $1$ 属于 $F^s$，且 $\|1\cdot d^{-(s+a+n/q)}\|_{L^q(\Om)}=\infty$：其 $q$ 次幂为 $\int_\Om d^{-(s+a)q-n}\,dx$，而 $c:=(s+a)q+n>n\ge1$，故由附录 A 得 $\int_\Om d^{-c}\,dx=\infty$。因此第二个 Hardy 项也绝不冗余。
  \end{enumerate}
\end{enumerate}
\end{theorem}

\begin{proof}
我们完整引用 Kijaczko（arXiv:2110.04012）的定理 4 与定理 3。

\emph{Kijaczko 定理 4}. 设 $\Om$ 为 $\R^d$ 中非空、有界、一致的开集，$0<s<1$、$1\le p<\infty$，并设 $\alpha,\beta$ 为满足 $0\le\alpha,\beta<\codimAl(\partial\Om)$ 的实数。若 $sp+\alpha+\beta>\codimAu(\partial\Om)$，则
\[
  W_0^{s,p;\alpha,\beta}(\Om)=\Big\{f\in W^{s,p;\alpha,\beta}(\Om):\int_\Om|f(x)|^p d_\Om(x)^{-sp-\alpha-\beta}\,dx<\infty\Big\},
\]
其中 $d_\Om(x)=\dist(x,\partial\Om)$，$W_0^{s,p;\alpha,\beta}(\Om)$ 为 $C_c^\infty(\Om)\cap W^{s,p;\alpha,\beta}(\Om)$ 在 $W^{s,p;\alpha,\beta}(\Om)$ 中的闭包（同上文定义 1）。

\emph{Kijaczko 定理 3}. 设 $\Om$ 为 $\R^d$ 中非空、有界开集，$0<s<1$、$1\le p<\infty$、$\alpha,\beta\ge0$。
\begin{enumerate}[label=(\Roman*)]
  \item 若 $sp+\alpha+\beta<d-\dimM(\partial\Om)$，则 $W_0^{s,p;\alpha,\beta}(\Om)=W^{s,p;\alpha,\beta}(\Om)$。
  \item 若 $\Om$ 是 $(d-sp-\alpha-\beta)$-齐次的，$p>1$ 且 $sp+\alpha+\beta=\codimAl(\partial\Om)$，则 $W_0^{s,p;\alpha,\beta}(\Om)=W^{s,p;\alpha,\beta}(\Om)$。
  \item 若 $\Om$ 是 $\kappa$-plump 的且 $sp+\alpha+\beta>\codimAu(\partial\Om)$，则 $W_0^{s,p;\alpha,\beta}(\Om)\neq W^{s,p;\alpha,\beta}(\Om)$。
\end{enumerate}

这里 $W^{s,p;\alpha,\beta}(\Om)$ 是 $L^p(\Om)$ 中使下述半范数有限的 $f$ 所成空间：
\[
  \Big(\int_\Om\int_\Om|f(x)-f(y)|^p|x-y|^{-d-sp}d(y)^{-\alpha}d(x)^{-\beta}\,dy\,dx\Big)^{1/p},
\]
并以 $\|f\|_{L^p}+[f]$ 赋范，$W_0^{s,p;\alpha,\beta}(\Om)$ 是 $C_c^\infty(\Om)\cap W^{s,p;\alpha,\beta}(\Om)$ 在其中的闭包。陈述中的单权空间 $W_b^{s,p}(\Om)$ 是 $\alpha,\beta$ 之一等于 $b$、另一个等于 $0$ 的特殊情形：两种取法下权均为 $d(y)^{-b}$（至多交换变量的名称），$sp+\alpha+\beta=sp+b$，且假设 $0\le\alpha,\beta<\codimAl(\partial\Om)$ 等价于 $0\le b<\codimAl(\partial\Om)$。

\emph{所用的结构事实}. (1) 每个有界 Lipschitz 域都是一致域（有界 Lipschitz 域是 $(\varepsilon,\delta)$-域，而作为 $(\varepsilon,\delta)$-域的有界域是一致域；参见 CLWZ 第 3--4 页关于 $(\varepsilon,\delta)$-域与一致域的讨论）。(2) 每个一致域对某个 $\kappa\in(0,1)$ 是 $\kappa$-plump 的（Kijaczko，定理 7 的证明）。(3) Kijaczko 第 2.1 节（定义 11 之后）记载：对 $\R^d$ 中的 $E$ 有 $\codimAl(E)\le d-\dimM(E)$；对偶地有 $d-\dimM(E)\le\codimAu(E)$。因此 $d-\dimM(\partial\Om)$ 介于下 Assouad 余维数与上 Assouad 余维数之间。(4) 距离 zeta 函数
\[
  \zeta_\Om(q)=\int_\Om d_\Om(x)^{-q}\,dx
\]
当 $q<n-\dimM(\partial\Om)$ 时有限，当 $q>n-\dimM(\partial\Om)$ 时无穷（Kijaczko，第 2.1 节，定义 11 及其后数行，依据 [14, Lemma 3.3 and Lemma 3.5]）。

\emph{(i) 的证明}. 对定理 4 取 $d=n$ 与 $\{\alpha,\beta\}=\{b,0\}$。假设 $0\le b<\codimAl(\partial\Om)$ 与 $sp+b>\codimAu(\partial\Om)$ 正是 (i) 的假设，故定理 4 给出 $W_{b,0}^{s,p}(\Om)$ 的所述刻画。关于真性：由 (1)--(2)，有界 Lipschitz 域 $\Om$ 是一致域，从而是 $\kappa$-plump 的；以相同的 $\{\alpha,\beta\}=\{b,0\}$ 与假设 $sp+b>\codimAu(\partial\Om)$ 应用定理 3(III)，得 $W_{b,0}^{s,p}(\Om)\neq W_b^{s,p}(\Om)$。因此所刻画的子空间为真子空间。

\emph{(ii) 的证明}. 对定理 3(I) 取 $d=n$ 与 $\{\alpha,\beta\}=\{b,0\}$：由 $sp+b<n-\dimM(\partial\Om)$ 得 $W_{b,0}^{s,p}(\Om)=W_b^{s,p}(\Om)$。由 (4)，$\int_\Om d^{-b}<\infty$ 等价于（在 $b=n-\dimM(\partial\Om)$ 之外）$b<n-\dimM(\partial\Om)$；这恰是 Kijaczko 定理 19 的假设，该定理给出 $C^\infty(\Om)\cap W_b^{s,p}(\Om)$ 在半赋范空间中稠密（与 Dyda--Kijaczko，arXiv:2009.00077，定理 3 结合则得完整范数下的稠密性）。因 $s>0$，有 $sp+b>b$，故对子区间 $0\le b<\codimAl(\partial\Om)$ 且 $sp+b>\codimAu(\partial\Om)$ 中的 $b$，利用 (3) 得
\[
  b<\codimAl(\partial\Om)\le n-\dimM(\partial\Om)\le\codimAu(\partial\Om)<sp+b;
\]
因此由部分 (i)，在该子区间内 $C^\infty$ 稠密而 $C_c^\infty$ 不稠密。该子区间的非空性并不来自 (iii)：有界 Lipschitz 域的边界是 $(n-1)$ 维 Ahlfors 正则集，故 $\codimAl(\partial\Om)=\codimAu(\partial\Om)=1$，所论区间即 $\big(\max\{0,1-sp\},\,1\big)$，而 $sp>0$ 保证其非空（(iii) 只是 $\Om=(0,1)$ 的一个具体见证）。这就证明了临界指数是 $sp+b$ 而非 $b$。

\emph{(iii) 的证明}. 对 $\Om=(0,1)$ 有 $\partial\Om=\{0,1\}$、$n=1$、$\dimM(\partial\Om)=0$ 且 $\codimAl(\partial\Om)=\codimAu(\partial\Om)=1$（见附录 A 第二步中的计算），且 $(0,1)$ 是有界 Lipschitz 域。取 $s=1/2$、$p=2$、$b=1/2$：
\[
  \int_\Om d^{-b}=\int_0^1 d^{-1/2}=2\sqrt2<\infty
\]
（附录 A 第一步），故 $b=1/2<1=\codimAl(\partial\Om)$；且 $sp+b=(1/2)(2)+1/2=3/2>1=\codimAu(\partial\Om)$。于是 (i) 适用，并给出
\[
  W_{b,0}^{s,p}(\Om)=\Big\{f\in W_b^{s,p}(\Om):\int_\Om|f(x)|^2d(x)^{-3/2}\,dx<\infty\Big\}.
\]
该子空间为真子空间且不含 $1$，因为 $\int_\Om d^{-3/2}=\infty$（附录 A，取 $c=3/2\ge1$；完整计算见附录 A 第四步）。因此 $\int d^{-aq}<\infty$ 并不蕴含 $C_c^\infty$ 稠密性。

\emph{(iv) 的证明}. 部分 (a)：若 $0<s<1$ 且 $p=q$，则 CLWZ 带单权 $d(y)^{-aq}$ 的加权内蕴空间 $F^s_{p,p}(\Om,d^{-aq})$ 恰是 $b=aq$ 时的 $W_b^{s,p}(\Om)$（半范数相同），且 $F^{s,1}=W_{b,0}^{s,p}(\Om)$。又 $sp+b=sp+aq=sp+ap=(s+a)p$。在所设情形下部分 (i) 适用，给出
\begin{align*}
  F^{s,1}&=\Big\{f\in F^s:\int_\Om|f(x)|^p d(x)^{-(sp+b)}\,dx<\infty\Big\}\\
  &=\Big\{f\in F^s:\int_\Om|f|^p d(x)^{-(s+a)p}\,dx<\infty\Big\}\\
  &=\Big\{f\in F^s:\|f\,d^{-(s+a)}\|_{L^p(\Om)}<\infty\Big\}.
\end{align*}
因此在该情形下第一个 CLWZ Hardy 范数恰是定义 $C_c^\infty$ 闭包的条件，故在该情形下它不能省略。

部分 (b)：对 $s\in(0,\infty)\setminus\N$、$a>0$ 与 $p\in[1,\infty)$，常数 $1$ 属于 $F^s$ 且 $\|1\|_{F^s}=|\Om|^{1/p}$（论证见附录 A 第一步）。若 $(s+a)p\ge1$，则由附录 A（取 $c=(s+a)p\ge1$）得 $\int_\Om d^{-(s+a)p}\,dx=\infty$，故 $\|1\cdot d^{-(s+a)}\|_{L^p(\Om)}=\infty$，而 $1$ 属于 $F^s$。因此在该范围内第一个 Hardy 范数不受内蕴范数控制。

部分 (c)：$\|1\cdot d^{-(s+a+n/q)}\|_{L^q(\Om)}$ 的 $q$ 次幂为 $\int_\Om d(x)^{-(s+a)q-n}\,dx$。这里 $c:=(s+a)q+n>n\ge1$，故由附录 A 该积分为无穷。因此对一切 $s\in(0,\infty)\setminus\N$、一切 $a>0$ 与一切 $q\in[1,\infty)$，常数 $1$ 的第二个 Hardy 范数为无穷，第二个 Hardy 项绝不冗余。
\end{proof}

\begin{remark}[直觉]
加权内在半范数在常数上消失，因此它看不到边界迹；而 Hardy 项恰恰能看见边界迹，所以它们是与内在范数独立的附加条件，不可省略。对 $C_c^\infty$ 稠密性而言，决定性的比较发生在 Hardy 指数 $sp+aq$ 与边界余维数之间，而不是在 $aq$ 与 $1$ 之间：权的可积性只给出 $C^\infty$ 稠密（Kijaczko 定理 19），而 $C_c^\infty$ 的闭包是由定理 4 描述的 Hardy 子空间。在 $(0,1)$ 上取 $s=1/2$、$p=2$、$aq=1/2$ 时 $aq<1$ 但 $sp+aq=3/2>1$，这正是常数 $1$ 属于内在空间却不属于 $C_c^\infty$ 闭包的原因。第二个 Hardy 项的奇性更强：其指数 $(s+a)q+n$ 恒大于 $n$，因此在任意参数下都在边界处发散，绝不可能冗余。
\end{remark}

%==================================================================
\subsection{可积性指数是 \texorpdfstring{$aq$}{aq} 而非 \texorpdfstring{$ap$}{ap}}

\begin{theorem}[Kijaczko 关于距离权 $d^{-aq}$ 的稠密性条件：可积性指数是 $aq$ 而非 $ap$]
设 $n\ge1$，$\Om$ 为 $\R^n$ 中非空有界域，$0<s<1$ 且 $1\le p<\infty$，并设 $w,v:\Om\to[0,\infty)$ 为 Borel 权。Kijaczko（arXiv:2110.04012，定理 19）考虑单指数单侧加权空间
\[
  \wt W^{s,p;w,v}(\Om):=\Big\{f\ \text{measurable on}\ \Om:
  \int_\Om\int_\Om|f(x)-f(y)|^p|x-y|^{-n-sp}w(x)v(y)\,dx\,dy<\infty\Big\},
\]
以相应的半范数赋范，并证明：当 $w$ 与 $v$ 局部有界且满足积分条件
\[
  (6)\qquad
  \int_\Om w(x)(1+|x|)^{-n-sp}\,dx<\infty
  \quad\text{与}\quad
  \int_\Om v(x)(1+|x|)^{-n-sp}\,dx<\infty .
\]
时，$C^\infty(\Om)\cap\wt W^{s,p;w,v}(\Om)$ 在 $\wt W^{s,p;w,v}(\Om)$ 中稠密。则：
\begin{enumerate}[label=(\roman*)]
  \item 因 $\Om$ 有界，对 $x\in\Om$ 有 $1\le(1+|x|)^{n+sp}\le(1+\sup_\Om|x|)^{n+sp}<\infty$，故条件 (6) 等价于 $w\in L^1(\Om)$ 与 $v\in L^1(\Om)$，即等价于 $\int_\Om w<\infty$ 与 $\int_\Om v<\infty$（此即 Kijaczko 的注 20）。
  \item 取 $\Om$ 上 $w:=1$，$v(y):=d(y)^{-aq}$，其中 $a>0$、$q\in[1,\infty)$ 为 CLWZ 范数 (1.2) 的内层指数。则条件 (6) 等价于
  \[
    \nuA(\Om)=\int_\Om d(y)^{-aq}\,dy<\infty,
  \]
  其中 $\nuA(E):=\int_E d(y)^{-aq}\,dy$。因此需要可积的指数是 $aq$，即混合范数 (1.2) 中附着于内层积分变量 $y$ 的指数，而不是 $ap$。当 $p=q$ 时两指数 $aq$ 与 $ap$ 相同；当 $p\neq q$ 时它们不同，且 $ap<1$ 与 $aq<1$ 两个条件互不蕴含：对 $\Om=(0,1)$ 与 $d(y)=\min(y,1-y)$，$\int_\Om d^{-b}<\infty$ 当且仅当 $b<1$，故若 $p<q$ 且 $1/q\le a<1/p$，则 $ap<1$ 但 $aq\ge1$；而若 $p>q$ 且 $1/p\le a<1/q$，则 $aq<1$ 但 $ap\ge1$。
  \item 因此，在情形 (I)（$\nuA(\Om)<\infty$）下，Kijaczko 定理 19 对每个 $p\in[1,\infty)$ 给出 $C^\infty(\Om)\cap\wt W^{s,p;1,d^{-aq}}(\Om)$ 在半赋范空间 $\wt W^{s,p;1,d^{-aq}}(\Om)$ 中稠密，即在加权 Gagliardo 半范数下逼近成立。当 $p=q$ 时该半范数正是 CLWZ 范数 (1.2) 中的加权 Gagliardo 半范数部分；定理 19 只控制 CLWZ 范数的这一部分，而不控制额外的 $L^q(\Om)$ 项。
\end{enumerate}
\end{theorem}

\begin{proof}
\emph{第一部分}. 部分 (i) 是直接的：当 $\Om$ 有界时，因子 $(1+|x|)^{-n-sp}$ 被仅依赖 $\Om,n,s,p$ 的正常数从上下方控制（因 $0<(1+|x|)^{-n-sp}\le1$ 且 $(1+|x|)^{n+sp}\le(1+\sup_{x\in\Om}|x|)^{n+sp}<\infty$），故 (6) 中每个积分有限当且仅当该权在 $\Om$ 上的相应积分有限，即当且仅当该权属于 $L^1(\Om)$。Kijaczko 在 arXiv:2110.04012 的注 20 中对有界 $\Om$ 陈述了同样的等价性。

\emph{第二部分}. 部分 (ii)：对 (i) 取 $w:=1$（此时 $\int_\Om w=|\Om|<\infty$）与 $v:=d^{-aq}$，后者在开集 $\Om$ 上连续且为正，故局部有界。由 (i)，(6) 中第二个条件等价于 $\int_\Om d(y)^{-aq}\,dy<\infty$，即 $\nuA(\Om)<\infty$。对 $\Om=(0,1)$ 与 $d(y)=\min(y,1-y)$，直接计算得
\[
  \int_0^1 d(y)^{-b}\,dy=2\int_0^{1/2}y^{-b}\,dy=\frac{2^b}{1-b}
  \quad\text{for } b<1,
  \qquad\text{与}\qquad =+\infty \quad\text{for } b\ge1 .
\]
因此两条件分别读作 $ap<1$ 与 $aq<1$。若 $p<q$ 且 $1/q\le a<1/p$，则 $ap<1$ 且 $aq\ge1$；若 $p>q$ 且 $1/p\le a<1/q$，则 $aq<1$ 且 $ap\ge1$。两种情形下恰有一个条件成立，故两条件互不蕴含。

\emph{第三部分}. 部分 (iii)：在情形 (I) 下，以单指数 $p$、$w:=1$ 与 $v:=d^{-aq}$（二者局部有界，且由 (ii) 满足 (6)）应用 Kijaczko 定理 19，得 $C^\infty(\Om)\cap\wt W^{s,p;1,d^{-aq}}(\Om)$ 在半赋范空间 $\wt W^{s,p;1,d^{-aq}}(\Om)$ 中稠密：逼近函数在加权 Gagliardo 半范数下收敛，该半范数为
\[
  \Big(\int_\Om\int_\Om|f(x)-f(y)|^p|x-y|^{-n-sp}d(y)^{-aq}\,dy\,dx\Big)^{1/p}.
\]当 $p=q$ 时该半范数正是 CLWZ 范数 (1.2) 在 $p=q$ 时出现的加权 Gagliardo 半范数；$p=q$ 时的完整 CLWZ 范数是该半范数加上 $L^q(\Om)$ 范数 $\|f\|_{L^q(\Om)}$，故定理 19 只控制半范数部分，$L^q$ 部分需要单独的逼近论证。
\end{proof}

\begin{remark}[直觉]
权 $d^{-aq}$ 落在内层（$y$-）积分上，而该积分对应的指数是 $q$；Kijaczko 的可积性条件对"哪个指数乘上权"并不敏感，因为在他的空间里单个指数 $p$ 同时乘在两个变量上。因此情形 (I) 的阈值是 $\int d^{-aq}<\infty$，而当 $p\neq q$ 时必须把它与任何 $ap$ 条件分开验证。
\end{remark}


%=====================================================================
\section{Hardy 项的非冗余性、延拓路线的终结与技术工具}\label{sec:part6}
%=====================================================================
% ===================================================================
%  第 6 组：Hardy 项的非冗余性、延拓路线的终结与技术工具
% ===================================================================

本节处理论证中的两个转折点。前两条事实说明 Hardy 增广项绝不是内蕴范数的推论：常数函数 $1$ 属于 $\Fsp$ 而其第二个 Hardy 范数 $\|f d^{-(s+a+n/q)}\|_{L^q(\Om)}$ 恒为 $+\infty$，因此 $\Fsp$ 永远不含于 $\tiF$，关于稠密性的正确断言只能在 $\tiF$ 内部表述。第三条事实进一步表明，当 $s>1$ 时连经典意义下的 $C_c^\infty$ 稠密性都不成立（Poincaré 障碍），而第四条事实则终结了权化延拓算子路线：在无限质量情形下 $\R^n$ 上的加权空间退化为 $\{0\}$，CLWZ 命题 1.12 的结论不可能成立。最后两条事实是为上述全部论证提供支撑的技术工具，分别是向量值 $L^p$ 与无权重混合 Gagliardo 半范数中的平移连续性，以及混合范数加权空间在局部与情形 (I) 下的可积性。

%==================================================================
\subsection{Hardy 项的非冗余性：常数函数属于 \texorpdfstring{$F^{s}_{p,q}$}{Fsp} 但不属于 \texorpdfstring{$\tiF$}{tilde F}}
%==================================================================

\begin{theorem}[常数属于 $F^{s}_{p,q}$ 但不属于 Hardy 加强空间：内蕴空间不包含于 $\wt F^{s}_{p,q}$]
设 $n\ge1$，$\Om$ 是 $\R^n$ 中的非空有界 Lipschitz 区域，$d(x)=\dist(x,\partial\Om)$（$x\in\Om$），$s\in(0,\infty)\setminus\N$，$a>0$，$p,q\in[1,\infty)$。记 $s=k+\sigma$，其中 $k\in\N_0$，$\sigma\in(0,1)$；设 $\Fsp$ 为 CLWZ (1.1)--(1.3) 的加权内蕴 Triebel--Lizorkin 空间，其范数定义为
\begin{equation}\label{eq:sec6-f1-norm}
  \begin{split}
    \|f\|_{F^s}
    &:=\sum_{|\alpha|\le k}\|D^\alpha f\|_{L^p(\Om)}\\
    &\quad+\sum_{|\beta|=k}\Big(\int_\Om\Big[\int_\Om|D^\beta f(x)-D^\beta f(y)|^q\,|x-y|^{-n-\sigma q}\,d(y)^{-aq}\,dy\Big]^{p/q}dx\Big)^{1/p},
  \end{split}
\end{equation}
其中每一项在其不为有限时约定取 $+\infty$。再令
\begin{equation}\label{eq:sec6-f1-tildenorm}
  \|f\|_{\wt F^s}:=\|f\|_{F^s}+\|f d^{-(s+a)}\|_{L^p(\Om)}+\|f d^{-(s+a+n/q)}\|_{L^q(\Om)},
\end{equation}
\begin{equation}\label{eq:sec6-f1-hatnorm}
  \|f\|_{\wh F^s}:=\|f\|_{F^s}+\|f d^{-(s+a)}\|_{L^p(\Om)},
\end{equation}
以及 $\tiF:=\{f\in\Fsp:\|f\|_{\wt F^s}<\infty\}$，$\hoF:=\{f\in\Fsp:\|f\|_{\wh F^s}<\infty\}$。则：
\begin{enumerate}[label=(\roman*)]
  \item 常数函数 $1$（在 $\Om$ 上恒等于 $1$）属于 $\Fsp$，且 $\|1\|_{F^s}=|\Om|^{1/p}$；
  \item $\int_\Om d(y)^{-(s+a)q-n}\,dy=+\infty$；等价地，函数 $1\cdot d^{-(s+a+n/q)}$ 的 $L^q$ 范数为 $+\infty$；
  \item 因此 $1$ 不属于 $\tiF$，且对每个 $s\in(0,\infty)\setminus\N$、每个 $a>0$ 与所有 $p,q\in[1,\infty)$，$\Fsp$ \emph{不}包含于 $\tiF$——特别地，对 $p=q$ 亦然；
  \item $1$ 属于 $\hoF$ 当且仅当 $\int_\Om d(y)^{-(s+a)p}\,dy<\infty$；特别地，只要 $(s+a)p\ge1$，$1$ 就不属于 $\hoF$。
\end{enumerate}
(iii) 的要点在于：第二个 Hardy 范数，即 $f d^{-(s+a+n/q)}$ 的 $L^q$ 范数，永远不可能是内蕴范数的推论，因为它含有额外的因子 $d^{-n}$，而该因子使得对\emph{每一个}在边界上不消失的 $f$，相应的权都不可积（对常数 $1$ 已经如此）。因此，作为关于原始内蕴空间的陈述，$\Fsp$ 包含于 $\tiF$ 是错误的；而 Hardy 加强空间的稠密性断言必须（在已验证的 $p=q$ 结果 附录 A 中确实如此）在 $\tiF$ \emph{内部}表述。
\end{theorem}

\begin{proof}
\emph{第 0 步（常数函数）}. 对每个满足 $|\alpha|\ge1$ 的重指标 $\alpha$ 都有 $D^\alpha 1=0$（弱导数）。若 $k\ge1$，则对每个满足 $|\beta|=k$ 的 $\beta$ 有 $D^\beta 1=0$，因为 $|\beta|=k\ge1$。若 $k=0$，第二个和只含单项 $\beta=0$，且 $\sigma=s$；该项消失，因为 $D^0 1=1$ 为常数，故对一切 $x,y\in\Om$ 有 $|D^0 1(x)-D^0 1(y)|=|1-1|=0$。于是，无论 $k=0$ 还是 $k\ge1$，\eqref{eq:sec6-f1-norm} 中 $\|1\|_{F^s}$ 的每一项都消失，仅剩第一个和中 $\alpha=0$ 的项，而该项为
\[
  \|1\|_{L^p(\Om)}=|\Om|^{1/p}<\infty
\]
（因为 $\Om$ 有界，故 $|\Om|<\infty$）。这就证明了 (i)。

\emph{第 1 步（$q$-Hardy 权的发散性）}. 取 $c:=(s+a)q+n$，应用附录 A（Lipschitz 区域边界层的线性下界，以及所有幂 $d^{-c}$（$c\ge1$）的不可积性）。由于 $s>0$、$a>0$、$q\ge1$、$n\ge1$，有 $c>n\ge1$，故 $c$ 是满足 $c\ge1$ 的实数；所引事实给出
\[
  \int_\Om d(y)^{-c}\,dy=+\infty,
\]
即 $\int_\Om d(y)^{-(s+a)q-n}\,dy=+\infty$。又 $(s+a+n/q)q=(s+a)q+n$，故 $\|1\cdot d^{-(s+a+n/q)}\|_{L^q(\Om)}$ 的 $q$ 次幂等于
\[
  \int_\Om\big(d(y)^{-(s+a+n/q)}\big)^q\,dy=\int_\Om d(y)^{-(s+a)q-n}\,dy=+\infty.
\]
这就证明了 (ii)。

\emph{第 2 步（不属于增广空间）}. 由 (i)，$1$ 是 $\Fsp$ 的元素。由 (ii)，$1$ 的 $L^q$-Hardy 项为无穷，故按 \eqref{eq:sec6-f1-tildenorm} 有 $\|1\|_{\wt F^s}=+\infty$，从而常数函数 $1$ 不是
\[
  \tiF=\{f\in\Fsp:\|f\|_{\wt F^s}<\infty\}
\]
的元素。由于 $1$ 属于原始空间而不属于带 tilde 的空间，原始空间 $\Fsp$ 不是 $\tiF$ 的子集。这就证明了 (iii)。注意论证中任何地方都没有用到 $p$ 与 $q$ 之间的关系，也没有用到 $d^{-aq}$ 的可积性（没有用到情形 (I)/(II) 的任何假设）：发散来自第二个 Hardy 权内部的因子 $d^{-n}$。

\emph{第 3 步（hat 空间）}. 由 (i) 以及 hat 范数的定义 \eqref{eq:sec6-f1-hatnorm}，$1$ 属于 $\hoF$ 当且仅当 $\|1\cdot d^{-(s+a)}\|_{L^p(\Om)}<\infty$，亦即当且仅当
\[
  \int_\Om d(y)^{-(s+a)p}\,dy<\infty
\]
（常数的 $F^s$ 项由 (i) 有限）。若 $(s+a)p\ge1$，则 $c':=(s+a)p$ 是 $\ge1$ 的实数，以 $c=c'$ 应用附录 A 得 $\int_\Om d(y)^{-(s+a)p}\,dy=+\infty$，故 $1$ 不在 hat 空间中。这就证明了 (iv)。
\end{proof}

%==================================================================
\subsection{同一非包含的独立证明：边界层的层饼计算}
%==================================================================

\begin{proposition}[对每个非整数光滑度与每个权，Hardy 加强空间都严格小于内蕴空间]
设 $n\ge1$，$\Om$ 是 $\R^n$ 中的非空有界 Lipschitz 区域，$d(x)=\dist(x,\partial\Om)$（$x\in\Om$），$s\in(0,\infty)\setminus\N$，记 $k:=\lfloor s\rfloor\in\N_0$，$\sigma:=s-k\in(0,1)$，并设 $a>0$，$p,q\in[1,\infty)$。设 $\Fsp$ 为 CLWZ (1.1)--(1.2) 的加权内蕴 Triebel--Lizorkin 空间，即使得
\begin{equation}\label{eq:sec6-f2-norm}
  \begin{split}
    \|f\|_{F^s}
    &:=\sum_{|\alpha|\le k}\|D^\alpha f\|_{L^p(\Om)}\\
    &\quad+\sum_{|\beta|=k}\Big(\int_\Om\Big[\int_\Om|D^\beta f(x)-D^\beta f(y)|^q\,|x-y|^{-n-\sigma q}\,d(y)^{-aq}\,dy\Big]^{p/q}dx\Big)^{1/p}
  \end{split}
\end{equation}
有限的 $f\in L^p(\Om)$ 之集（每一项在其不为有限时取 $+\infty$）。令 $\|f\|_{\wt F^s}:=\|f\|_{F^s}+\|f d^{-(s+a)}\|_{L^p(\Om)}+\|f d^{-(s+a+n/q)}\|_{L^q(\Om)}$，$\tiF:=\{f\in\Fsp:\|f\|_{\wt F^s}<\infty\}$。则：
\begin{enumerate}[label=(\roman*)]
  \item 常数函数 $1$（在 $\Om$ 上恒等于 $1$）属于 $\Fsp$，且 $\|1\|_{F^s}=|\Om|^{1/p}$；
  \item $\int_\Om d(y)^{-(s+a)q-n}\,dy=+\infty$，等价地 $\|1\cdot d^{-(s+a+n/q)}\|_{L^q(\Om)}=+\infty$；
  \item 从而 $1\notin\tiF$，且 $\Fsp$ \emph{不}包含于 $\tiF$。这对每个 $s\in(0,\infty)\setminus\N$、每个 $a>0$ 与所有 $p,q\in[1,\infty)$ 成立；特别地，对 $p=q$ 也成立。
\end{enumerate}
\end{proposition}

\begin{proof}
\emph{第 1 步（常数 $1$ 的内蕴范数）}. 令 $f:=1$（在 $\Om$ 上）。每个满足 $|\alpha|\ge1$ 的弱导数 $D^\alpha f$ 都是零函数，且 $\|D^0 f\|_{L^p(\Om)}=\|1\|_{L^p(\Om)}=|\Om|^{1/p}$。若 $k\ge1$，则 \eqref{eq:sec6-f2-norm} 第二个和中的每一项都有 $|\beta|=k\ge1$，故 $D^\beta f=0$，相应的加权 Gagliardo 半范数消失，因为对一切 $x,y\in\Om$ 有 $|D^\beta f(x)-D^\beta f(y)|=|0-0|=0$。若 $k=0$（即 $0<s<1$），第二个和只含唯一的重指标 $\beta=0$，其对应项为
\[
  \Big(\int_\Om\Big[\int_\Om|1-1|^q\,|x-y|^{-n-\sigma q}\,d(y)^{-aq}\,dy\Big]^{p/q}dx\Big)^{1/p}=0.
\]
两种情形下整个第二个和都为零，而第一个和等于 $\|1\|_{L^p(\Om)}=|\Om|^{1/p}$。由于 $\Om$ 有界，$|\Om|<\infty$，故 $1\in\Fsp$ 且 $\|1\|_{F^s}=|\Om|^{1/p}$。这就证明了 (i)。

\emph{第 2 步（Hardy 指数）}. 令 $c:=(s+a)q+n$。由于 $n\ge1$、$s>0$、$a>0$、$q\ge1$，有 $c>n\ge1$，特别地 $c\ge1$。对这样的 $c$ 应用附录 A（对 $\R^n$ 中的有界 Lipschitz 区域 $\Om$ 与每个实数 $c\ge1$，$\int_\Om d(x)^{-c}\,dx=\infty$），得
\[
  \int_\Om d(y)^{-(s+a)q-n}\,dy=\int_\Om d(y)^{-c}\,dy=+\infty.
\]
为完整起见，下面给出这一发散性的直接计算。由层饼公式，对 $c>0$ 有
\[
  \int_\Om d(y)^{-c}\,dy=c\int_0^\infty t^{-c-1}|\Omt|\,dt,
  \qquad \Omt=\{y\in\Om:d(y)<t\};
\]
这由代换 $t=d(y)$ 与 Fubini 定理得到。由于 $\Om$ 是有界 Lipschitz 区域，存在 $\kappa>0$ 与 $t_0>0$，使得对一切 $t\in(0,t_0)$ 有 $|\Omt|\ge\kappa t$（这正是附录 A 中所证的边界层线性下界）。于是
\[
  \int_\Om d^{-c}\,dy\ge c\kappa\int_0^{t_0}t^{-c-1}t\,dt=c\kappa\int_0^{t_0}t^{-c}\,dt=+\infty,
\]
这是因为 $c\ge1$：当 $c>1$ 时原函数 $t^{1-c}/(1-c)$ 在 $t\to0^+$ 时趋于 $+\infty$，而当 $c=1$ 时截断积分等于 $\log(t_0/\varepsilon)\to+\infty$（$\varepsilon\to0^+$）。这在当前记号下重新证明了所引的发散性。

\emph{第 3 步（$1$ 的 $L^q$ Hardy 范数）}. 由于 $q\ge1$，函数 $1\cdot d^{-(s+a+n/q)}$ 的 $L^q(\Om)$ 范数的 $q$ 次幂为
\[
  \int_\Om|1\cdot d(x)^{-(s+a+n/q)}|^q\,dx=\int_\Om d(x)^{-(s+a)q-n}\,dx,
\]
由第 2 步知其为 $+\infty$。故 $\|1\cdot d^{-(s+a+n/q)}\|_{L^q(\Om)}=+\infty$，这正是 (ii) 中所断言的等价性。

\emph{第 4 步（$1$ 不在增广空间中）}. 常数函数 $1$ 的 $\wt F^s$ 范数为
\[
  \|1\|_{\wt F^s}=\|1\|_{F^s}+\|d^{-(s+a)}\|_{L^p(\Om)}+\|d^{-(s+a+n/q)}\|_{L^q(\Om)}.
\]
由第 1 步，第一项等于 $|\Om|^{1/p}<\infty$；第二项是非负的广义实数；由第 3 步，第三项为 $+\infty$。因此 $\|1\|_{\wt F^s}=+\infty$，故 $1\notin\tiF$。

\emph{第 5 步（内蕴空间不含于增广空间）}. 由第 1 步，$1$ 属于 $\Fsp$；由第 4 步，$1$ 不属于 $\tiF$。因此 $\Fsp$ 不是 $\tiF$ 的子集。由于 $s\in(0,\infty)\setminus\N$、$a>0$、$p,q\in[1,\infty)$ 都是任意的，这就证明了 (iii)，其中包含 $p=q$ 的情形。
\end{proof}

\begin{remark}[直觉]
内蕴的加权半范数对常数是盲的：常数 $1$ 的一切差分都消失。Hardy 项则相反，它测量边界迹。第二个 Hardy 指数为 $(s+a+n/q)q=(s+a)q+n>n$，而在有界 Lipschitz 区域上，由于边界层 $\{d<t\}$ 的测度至少是常数乘以 $t$，每个幂 $d^{-c}$（$c\ge1$）都不可积。于是对每个允许的参数，常数 $1$ 的 $L^q$ Hardy 范数都是无穷，Hardy 增广空间因而永远不包含内蕴空间；对 $s,a,p,q$ 的任何参数限制都无法使这一项成为冗余的。
\end{remark}

%==================================================================
\subsection{Poincar\'e 障碍：非整数 \texorpdfstring{$s>1$}{s>1} 时 \texorpdfstring{$C_c^\infty$}{Cc-infinity} 从不稠密}
%==================================================================

\begin{theorem}[Poincar\'e 障碍：非整数 $s>1$ 时，紧支光滑函数在 $F^{s}$ 中从不稠密]
设 $n\ge1$，$\Om$ 是 $\R^n$ 中的非空有界区域，$d(x)=\dist(x,\partial\Om)$（$x\in\Om$），$k\in\N$，$\sigma\in(0,1)$，$s=k+\sigma$，$q\in[1,\infty)$，$a\in\R$。令 $w(y):=d(y)^{-aq}$，并设 $F^{s}_{q,q}(\Om,d^{-aq})$ 为按 (1.2) 取 $p=q$ 赋予范数的加权内蕴空间，即
\begin{equation}\label{eq:sec6-f3-norm}
  \begin{split}
    \|f\|_{F^s}
    &:=\sum_{|\alpha|\le k}\|D^\alpha f\|_{L^q(\Om)}\\
    &\quad+\sum_{|\beta|=k}\Big(\int_\Om\Big[\int_\Om|D^\beta f(x)-D^\beta f(y)|^q\,|x-y|^{-n-\sigma q}\,w(y)\,dy\Big]dx\Big)^{1/q}.
  \end{split}
\end{equation}
则：
\begin{enumerate}[label=(\roman*)]
  \item 常数函数 $1$（在 $\Om$ 上恒等于 $1$）属于 $F^{s}_{q,q}(\Om,d^{-aq})$，且 $\|1\|_{F^s}=|\Om|^{1/q}$（事实上对 $|\beta|=k\ge1$ 有 $D^\beta 1=0$，故所有分数次项都消失）；
  \item 若 $Q$ 是包含 $\Om$ 的开立方体，$L$ 为其边长，则每个 $u\in C_c^\infty(\Om)$（零延拓到 $\R^n$）满足
  \[
    \|u\|_{L^q(\Om)}\le L\,\|\nabla u\|_{L^q(\Om)};
  \]
  \item 从而，对每个满足 $u\in F^{s}_{q,q}(\Om,d^{-aq})$ 的 $u\in C_c^\infty(\Om)$ 有
  \[
    \|1-u\|_{F^s}\ge|\Om|^{1/q}/(1+L)>0.
  \]
\end{enumerate}
于是 $C_c^\infty(\Om)\cap F^{s}_{q,q}(\Om,d^{-aq})$ 在 $F^{s}_{q,q}(\Om,d^{-aq})$ 中的闭包 $F^{s,1}_{q,q}(\Om,d^{-aq})$ 不含常数函数 $1$，从而是 $F^{s}_{q,q}(\Om,d^{-aq})$ 的真闭子空间：
\[
  F^{s,1}_{q,q}(\Om,d^{-aq})\ne F^{s}_{q,q}(\Om,d^{-aq})
\]
对每个 $a\in\R$ 与每个 $q\in[1,\infty)$ 成立。特别地，在高阶情形 $s>1$ 下，经典稠密性陈述（问题 I）对每个权指数 $a$ 都不成立，只有关于 $C^\infty$（而非 $C_c^\infty$）的较弱陈述才有可能成立。
\end{theorem}

\begin{proof}
\emph{第 0 步（设定）}. $\Om$ 是 $\R^n$ 的非空有界开子集；令 $D:=\diam(\Om)\in(0,\infty)$。对每个实数 $a$，权 $w(y)=d(y)^{-aq}$ 是 $\Om$ 上的正 Borel 函数。\eqref{eq:sec6-f3-norm} 中所写的范数正是 CLWZ 文稿的 (1.2) 在 $p=q$ 时的形式。

\emph{第 1 步（常数函数）}. 对函数 $1$ 有 $D^0 1=1$，且对每个满足 $|\alpha|\ge1$ 的重指标 $\alpha$ 有 $D^\alpha 1=0$。由于 $k\ge1$（因为 $s>1$），对每个 $|\beta|=k$ 函数 $D^\beta 1$ 恒为零，故定义第 $\beta$ 个分数次项的内层积分为零。因此
\[
  \|1\|_{F^s}=\|1\|_{L^q(\Om)}=|\Om|^{1/q}<\infty,
\]
这就证明了 (i)。

\emph{第 2 步（经由立方体的 Poincar\'e 不等式）}. 取 $\R^n$ 中的开立方体 $Q=I_1\times\cdots\times I_n$，其中每个 $I_i$ 是长度为 $L$ 的开区间，使得 $\Om$ 的闭包含于 $Q$（这样的立方体存在，因为 $\Om$ 有界；例如取 $L>2\diam(\Om)$，并令 $Q$ 以 $\Om$ 中某点为中心）。设 $u\in C_c^\infty(\Om)$，$U$ 为其零延拓到 $\R^n$。则 $U$ 除可能在 $\partial\Om$ 上之外是光滑的，$U$ 是 Lebesgue 可测且有界的，且 $\supp U$ 含于 $\Om$，从而含于 $Q$。固定 $x=(x_1,\dots,x_n)\in Q$。当 $t<\inf I_1$ 时，点 $(t,x_2,\dots,x_n)$ 落在 $Q$ 之外，从而落在 $\supp U$ 之外，故 $U(t,x_2,\dots,x_n)=0$。因此由微积分基本定理，
\[
  U(x)=\int_{\inf I_1}^{x_1}(\partial_1U)(t,x_2,\dots,x_n)\,dt.
\]
由 H\"older 不等式，对 $q\in(1,\infty)$ 及其共轭指数 $q'=q/(q-1)$ 有
\[
  |U(x)|^q\le\Big(\int_{\inf I_1}^{x_1}1^{q'}\,dt\Big)^{q/q'}\int_{I_1}|(\partial_1U)(t,x_2,\dots,x_n)|^q\,dt
  \le L^{q-1}\int_{I_1}|(\partial_1U)(t,x_2,\dots,x_n)|^q\,dt,
\]
而对 $q=1$，同一不等式以因子 $L^0=1$ 成立（直接由上述恒等式与积分的三角不等式）。将此不等式对 $x\in Q$ 积分并用 Fubini 定理，
\[
  \int_Q|U|^q\,dx\le L^{q-1}\cdot L\int_{Q'}\int_{I_1}|(\partial_1U)(t,x')|^q\,dt\,dx'
  =L^q\int_Q|\partial_1U|^q\,dx\le L^q\int_Q|\nabla U|^q\,dx,
\]
其中 $Q'=I_2\times\cdots\times I_n$，前面的因子 $L$ 来自将 $x_1$ 变量在 $I_1$（长度为 $L$）上积分；对 $q=1$，同样的计算给出 $\int_Q|U|\le L\int_Q|\partial_1U|$。由于在 $\Om$ 上 $U=u$、在 $\Om$ 外 $U=0$，
\[
  \int_Q|U|^q=\|u\|^q_{L^q(\Om)},\qquad \int_Q|\nabla U|^q=\|\nabla u\|^q_{L^q(\Om)}.
\]
因此 $\|u\|_{L^q(\Om)}\le L\|\nabla u\|_{L^q(\Om)}$，这正是 (ii)，其中显式常数 $C_P=L$。

\emph{第 3 步（下界）}. 设 $u\in C_c^\infty(\Om)$ 且 $u\in F^{s}_{q,q}(\Om,d^{-aq})$，令 $g:=1-u$。由范数的定义以及 \eqref{eq:sec6-f3-norm} 中第一个和对一切 $|\alpha|\le k$ 求和（而 $k\ge1$）可知
\[
  \|g\|_{F^s}\ge\sum_{|\alpha|\le k}\|D^\alpha g\|_{L^q(\Om)}
  \ge\sum_{|\alpha|=1}\|D^\alpha g\|_{L^q(\Om)}
  =\sum_{i=1}^n\|\partial_i u\|_{L^q(\Om)}\ge\|\nabla u\|_{L^q(\Om)},
\]
因为对 $|\alpha|=1$ 有 $D^\alpha 1=0$，且逐点有 $|\nabla u|=(\sum_i|\partial_i u|^2)^{1/2}\le\sum_i|\partial_i u|$。另一方面，由第 2 步，
\[
  \|\nabla u\|_{L^q(\Om)}\ge L^{-1}\|u\|_{L^q(\Om)}
  \ge L^{-1}\big(\|1\|_{L^q(\Om)}-\|1-u\|_{L^q(\Om)}\big)
  =L^{-1}\big(|\Om|^{1/q}-\|g\|_{L^q(\Om)}\big),
\]
其中中间的估计是三角不等式 $\|1\|\le\|u\|+\|1-u\|$。由于 $\|g\|_{L^q}\le\|g\|_{F^s}$，合并上面两式得
\[
  \|g\|_{F^s}\ge L^{-1}\big(|\Om|^{1/q}-\|g\|_{F^s}\big),
\]
故 $(1+L^{-1})\|g\|_{F^s}\ge L^{-1}|\Om|^{1/q}$，即 $\|g\|_{F^s}\ge|\Om|^{1/q}/(1+L)>0$。这就是 (iii)。

\emph{第 4 步（结论）}. 由第 3 步，$F^s$ 中固定元素 $1$ 与 $C_c^\infty(\Om)\cap F^{s}_{q,q}(\Om,d^{-aq})$ 的每个元素之间的 $F^s$-距离至少为 $|\Om|^{1/q}/(1+L)>0$。因此 $1$ 不属于该集合的闭包 $F^{s,1}_{q,q}(\Om,d^{-aq})$，而由第 1 步 $1$ 确实属于 $F^{s}_{q,q}(\Om,d^{-aq})$。所以对每个实数 $a$ 与每个 $q\in[1,\infty)$，$F^{s,1}_{q,q}(\Om,d^{-aq})$ 都真包含于 $F^{s}_{q,q}(\Om,d^{-aq})$。（论证没有用到关于权的任何假设；特别地，该结论在 $\nuA=d^{-aq}dy$ 的二择情形 (I) 与情形 (II) 中都成立。）
\end{proof}

\begin{remark}[直觉]
$F^s$ 范数控制 $W^{k,q}$ 范数（$k\ge1$），从而控制任何紧支光滑函数 $u$ 的梯度的 $L^q$ 范数。若 $u$ 在 $F^s$ 中接近常数 $1$，则 $u$ 在 $L^q$ 中接近 $1$，而其梯度必须很小；但 $W^{1,q}_0$ 的 Poincar\'e 不等式迫使紧支 $u$ 的 $L^q$ 范数被其梯度的 $L^q$ 范数控制。因此没有任何紧支光滑函数能在 $F^s$ 中逼近常数 $1$，于是对 $s>1$，$C_c^\infty$ 从不稠密。
\end{remark}

%==================================================================
\subsection{无限质量情形：\texorpdfstring{$\R^n$}{Rn} 上的加权空间退化为 \texorpdfstring{$\{0\}$}{0}，加权延拓算子不存在}
%==================================================================

\begin{theorem}[$\R^n$ 上加权空间的平凡性，以及无限质量情形下加权延拓的不存在性]
设 $\Om$ 是 $\R^n$ 中的非空有界区域，$d(y)=\dist(y,\partial\Om)$（$y\in\R^n$），于是 $d$ 在整个 $\R^n$ 上有定义且恰在 $\partial\Om$ 上为零；设 $a>0$，$p,q\in[1,\infty)$，$s\in(0,\infty)\setminus\N$，$k=\lfloor s\rfloor$，$\sigma=s-k\in(0,1)$。令 $w(y)=d(y)^{-aq}$，并对 $r>0$ 令 $N_r=\{y\in\R^n:d(y)<r\}$，且假设 $\mathrm{(H_{\infty})}$：对每个 $r>0$ 有 $\int_{N_r}w(y)\,dy=+\infty$。（只要 $aq\ge1$ 且 $\partial\Om$ 满足下方 Ahlfors 正则性 $|\{d<t\}|\ge ct$（对一切充分小的 $t>0$），该假设就成立；特别地，对每个有界 Lipschitz 区域都成立：Lipschitz 域的边界是 $(n-1)$ 维 Ahlfors 正则集，故存在 $c>0$ 使 $|\{d<t\}|\ge ct$ 对一切充分小的 $t>0$ 成立。）\defterm{一般一致区域不具备此处所需的该下界论证}，本文只在有界 Lipschitz 区域的情形使用 (H$_{\infty}$)。设 $F^{s}_{p,q}(\R^n,d^{-aq})$ 为 $\R^n$ 上由 CLWZ 范数 (1.2)（两个积分都在 $\R^n$ 上取，权为 $d(y)^{-aq}$）定义的加权内蕴 Triebel--Lizorkin 空间，即
\begin{equation}\label{eq:sec6-f4-norm}
  \begin{split}
    \|g\|_{F^s_{p,q}(\R^n,d^{-aq})}
    &:=\sum_{|\alpha|\le k}\|D^\alpha g\|_{L^p(\R^n)}\\
    &\quad+\sum_{|\beta|=k}\Big(\int_{\R^n}\Big(\int_{\R^n}|D^\beta g(x)-D^\beta g(y)|^q\,|x-y|^{-n-\sigma q}\,w(y)\,dy\Big)^{p/q}dx\Big)^{1/p}.
  \end{split}
\end{equation}
则 $F^{s}_{p,q}(\R^n,d^{-aq})=\{0\}$。从而，若此外还有 $\nuA(\Om)=\infty$（由前述事实，这等价于源空间 $\Fsp$ 等于 $P_k(\Om)$，而后者含常数函数 $1$），则\emph{不存在}有界线性算子
\[
  E\colon\Fsp\longrightarrow F^{s}_{p,q}(\R^n,d^{-aq})
\]
使得对源空间中的一切 $f$ 有 $(Ef)|_\Om=f$。特别地，CLWZ 命题 1.12 的结论（$\Lambda_0$ 的加权有界性）在该情形下不成立，因此区域 $aq\ge1$（带 $(n-1)$-Ahlfors 正则边界）被排除在该命题的任何正确版本之外。
\end{theorem}

\begin{proof}
\emph{第 0 步（球上关于一般权的刚性引理）}. 设 $B$ 是 $\R^n$ 中半径为 $R$ 的开球，$w$ 是 $B$ 上的非负 Lebesgue 可测函数且 $\int_B w=\infty$，$\sigma\in(0,1)$，$p,q\in[1,\infty)$，并设 $g$ 在 $B$ 上 Lebesgue 可测且满足
\[
  \int_B\Big(\int_B|g(x)-g(y)|^q\,|x-y|^{-n-\sigma q}\,w(y)\,dy\Big)^{p/q}dx<\infty.
\]
则 $g$ 在 $B$ 上几乎处处为常数。

\emph{证明}. 令 $J(x):=\int_B|g(x)-g(y)|^q|x-y|^{-n-\sigma q}w(y)\,dy$；则 $J$ 可测且 $\int_B J^{p/q}<\infty$，故对几乎每个 $x$ 有 $J(x)<\infty$。令 $E:=\{x\in B:J(x)<\infty\}$，则 $|B\setminus E|=0$。固定 $x\in E$。对一切 $y\in B$ 有 $|x-y|\le2R$，故 $|x-y|^{-n-\sigma q}\ge(2R)^{-n-\sigma q}$；因此
\[
  \int_B|g(x)-g(y)|^q\,w(y)\,dy\le(2R)^{n+\sigma q}J(x)<\infty.
\]
对 $m\in\N$ 令 $A_m(x):=\{y\in B:|g(x)-g(y)|>1/m\}$。则
\[
  m^{-q}\int_{A_m(x)}w\le\int_B|g(x)-g(y)|^q\,w(y)\,dy<\infty,
\]
故 $\int_{A_m(x)}w<\infty$，从而记 $B_m(x):=B\setminus A_m(x)$ 便有
\[
  \int_{B_m(x)}w=\int_B w-\int_{A_m(x)}w=\infty.
\]
现在取 $x_1,x_2\in E$。对每个 $m$，集合 $B_m(x_1)\cap B_m(x_2)$ 具有无限的 $w$-测度（它在 $B$ 中的补集含于 $A_m(x_1)\cup A_m(x_2)$，而后者 $w$-测度有限），故非空；取其中一点 $y$。则
\[
  |g(x_1)-g(x_2)|\le|g(x_1)-g(y)|+|g(y)-g(x_2)|\le2/m.
\]
由 $m$ 的任意性得 $g(x_1)=g(x_2)$。故 $g$ 在 $E$ 上为常数，从而在 $B$ 上几乎处处为常数。

\emph{第 1 步（$\R^n$ 上的空间是 $\{0\}$）}. 设 $g\in F^{s}_{p,q}(\R^n,d^{-aq})$。则 $g\in L^p(\R^n)$，且对每个满足 $|\beta|=k$ 的重指标 $\beta$ 有 $D^\beta g\in L^p(\R^n)$，并且
\[
  \int_{\R^n}\Big(\int_{\R^n}|D^\beta g(x)-D^\beta g(y)|^q\,|x-y|^{-n-\sigma q}\,w(y)\,dy\Big)^{p/q}dx<\infty.
\]
固定一个包含 $\partial\Om$ 之邻域的球 $B_R$（例如 $\partial\Om$ 含于半径为 $R/2$ 的球中）；则按 $\mathrm{(H_{\infty})}$ 有 $\int_{B_R}w=\infty$，而限制在 $B_R\times B_R$ 上的二重积分有限。在 $B_R$ 上以 $D^\beta g$ 代替 $g$ 应用第 0 步，可知对每个 $|\beta|=k$，$D^\beta g$ 在 $B_R$ 上几乎处处为常数。由于这对一切充分大的 $R$ 成立且球 $B_R$ 穷尽 $\R^n$，所有 $D^\beta g$（$|\beta|=k$）都在 $\R^n$ 上几乎处处为常数。在 $B_R$（有界区域）上应用已验证的多项式引理 附录 A，知 $g$ 在 $B_R$ 上几乎处处等于一个次数不超过 $k$ 的多项式；令 $R$ 沿一列数趋于无穷，便得到同一个次数不超过 $k$ 的多项式 $P$ 使得 $g=P$ 在 $\R^n$ 上几乎处处成立（两个这样的多项式在一个非空开集上相等）。由于 $g\in L^p(\R^n)$ 而非零多项式不属于 $L^p(\R^n)$（\defterm{非零常数多项式}因 $\R^n$ 的测度无穷而不属于 $L^p(\R^n)$；\defterm{次数至少为 $1$ 的非零多项式}在某个正测度的锥形区域内按多项式增长，故亦不属于 $L^p(\R^n)$），故 $P=0$，即 $g=0$ 几乎处处。因此 $F^{s}_{p,q}(\R^n,d^{-aq})=\{0\}$。

\emph{第 2 步（不存在有界的加权延拓）}. 再假设 $\nuA(\Om)=\infty$。由已验证的前述附录 A，$\Fsp=P_k(\Om)$，其中含常数函数 $1$。设 $E\colon\Fsp\to F^{s}_{p,q}(\R^n,d^{-aq})$ 线性有界，且对源空间中的每个 $f$ 满足 $(Ef)|_\Om=f$。则由第 1 步，目标空间中的 $E(1)$ 等于 $\{0\}$，但对每个 $x\in\Om$ 有 $((E(1))|_\Om)(x)=1$，故 $E(1)$ 不是目标空间的零元素（它在 $\R^n$ 上不几乎处处为零，在 $L^p(\R^n)$ 中也不为零）。此矛盾表明不存在这样的有界线性延拓算子。特别地，只要 $\nuA(\Om)=\infty$ 且 $\mathrm{(H_{\infty})}$ 成立，CLWZ (6.2) 中的算子 $\Lambda_0$（按构造满足在 $\Om$ 上 $\Lambda_0f=f$）就不可能从 $\Fsp$ 有界地映到 $F^{s}_{p,q}(\R^n,d^{-aq})$；因此任何能推出 CLWZ 命题 1.12 之结论的假设都必须排除这一情形，特别地，只要 $\partial\Om$ 是 $(n-1)$-Ahlfors 正则的，它就必须迫使 $aq<1$，因为此时一旦 $aq\ge1$，$\nuA(\Om)=\infty$ 与 $\mathrm{(H_{\infty})}$ 便同时成立（对 $\mathrm{(H_{\infty})}$：$|\{d<t\}|\ge ct$（$t>0$ 小）给出 $aq\ge1$ 时 $\int_{N_r}d^{-aq}\ge c\int_0^rt^{-aq}\,dt=\infty$）。
\end{proof}

\begin{remark}[直觉]
若距离权在边界的每个邻域内都有无限质量，则分数次半范数迫使空间中的每个函数在任意球上为常数（在有界球上权的无限测度压倒了有界的核），从而整体上是多项式；而 $L^p(\R^n)$ 又消灭一切非零多项式（分常数与非常数两种情形，见第 1 步）。因此整个 $\R^n$ 上的加权空间是 $\{0\}$。由于 $\Om$ 上的源空间退化为多项式（其中含常数 $1$），而任何延拓算子都必须在 $\Om$ 上还原 $1$，故不可能存在进入这个平凡空间的有界延拓。于是 $\Lambda_0$ 的加权有界性在无限质量情形下不可能成立。
\end{remark}

%==================================================================
\subsection{技术工具 I：向量值 \texorpdfstring{$L^p$}{Lp} 与无权重混合 Gagliardo 半范数中的平移连续性}
%==================================================================

\begin{lemma}[向量值 $L^p$ 与无权重混合 Gagliardo 半范数中的平移连续性]
设 $n\ge1$。

\emph{(a)} 设 $1\le p<\infty$，$(E,\|\cdot\|_E)$ 是 Banach 空间。设 $L^p(\R^n;E)$ 为 Bochner 可测映射 $g\colon\R^n\to E$ 且 $\int_{\R^n}\|g(x)\|_E^p\,dx<\infty$ 的 Bochner 空间，其范数为
\[
  \|g\|_{L^p(\R^n;E)}^p:=\int_{\R^n}\|g(x)\|_E^p\,dx.
\]
对 $z\in\R^n$ 令 $(\tau_zg)(x):=g(x-z)$。则对每个 $g\in L^p(\R^n;E)$ 有
\[
  \|\tau_zg-g\|_{L^p(\R^n;E)}\longrightarrow0\qquad(z\to0).
\]

\emph{(b)} 设 $0<s<1$，$1\le p,q<\infty$。对 Lebesgue 可测的 $h\colon\R^n\to\R$ 令
\[
  [h]_{\mathrm{unw},p,q}:=\Big(\int_{\R^n}\Big(\int_{\R^n}|h(x)-h(y)|^q\,|x-y|^{-n-sq}\,dy\Big)^{p/q}dx\Big)^{1/p}\in[0,\infty].
\]
若 $[h]_{\mathrm{unw},p,q}<\infty$，则当 $z\to0$ 时 $[\tau_zh-h]_{\mathrm{unw},p,q}\to0$。
\end{lemma}

\begin{proof}
\emph{(a) 第 1 步}. 设 $\varphi\in C_c(\R^n;E)$（连续、紧支、$E$-值）。则 $\varphi$ 一致连续：对 $\varepsilon>0$ 存在 $\delta>0$，使得只要 $|x-y|\le\delta$ 就有 $\|\varphi(x)-\varphi(y)\|_E\le\varepsilon$。设 $\supp\varphi$ 含于 $\{|x|\le R\}$。若 $|z|\le1$，则 $\supp(\tau_z\varphi)$ 含于 $\{|x|\le R+1\}$，且
\[
  \|\tau_z\varphi-\varphi\|_{L^p(\R^n;E)}^p=\int_{|x|\le R+1}\|\varphi(x-z)-\varphi(x)\|_E^p\,dx\le(2R+2)^n\varepsilon^p
\]
对 $|z|\le\min(\delta,1)$ 成立。

\emph{(a) 第 2 步}. 当 $1\le p<\infty$ 时 $C_c(\R^n;E)$ 在 $L^p(\R^n;E)$ 中稠密。给定 $g\in L^p(\R^n;E)$ 与 $\varepsilon>0$，取 $\varphi\in C_c(\R^n;E)$ 使 $\|g-\varphi\|\le\varepsilon/4$。由于范数在平移下不变，
\[
  \|\tau_zg-g\|\le\|\tau_z(g-\varphi)\|+\|\tau_z\varphi-\varphi\|+\|\varphi-g\|=2\|g-\varphi\|+\|\tau_z\varphi-\varphi\|;
\]
取 $|z|$ 充分小使 $\|\tau_z\varphi-\varphi\|\le\varepsilon/2$，便得 $\|\tau_zg-g\|\le\varepsilon$。

\emph{(b) 设定}. 对 $\varepsilon>0$ 令 $\mu_\varepsilon(du):=|u|^{-n-sq}\mathbf{1}_{\{|u|\ge\varepsilon\}}\,du$。这是一个有限测度：
\[
  \mu_\varepsilon(\R^n)=\omega_{n-1}\int_\varepsilon^\infty r^{-1-sq}\,dr=\omega_{n-1}\varepsilon^{-sq}/(sq)<\infty .
\]
令 $E_\varepsilon:=L^q(\R^n,\mu_\varepsilon)$；由于 $1\le q<\infty$，这是可分 Banach 空间。定义
\[
  G(x)(u):=h(x)-h(x-u)\quad(x,u\in\R^n),\qquad G^\varepsilon(x)(u):=G(x)(u)\mathbf{1}_{\{|u|\ge\varepsilon\}},
\]
以及
\[
  F(x):=\Big(\int_{\R^n}|G(x)(u)|^q\,|u|^{-n-sq}\,du\Big)^{1/q},
\]
于是 $\int_{\R^n}F(x)^p\,dx=[h]_{\mathrm{unw},p,q}^p<\infty$；故对几乎每个 $x$ 有 $F(x)<\infty$，即对几乎每个 $x$ 与每个 $\varepsilon>0$ 有 $G(x)\in L^q(\mu_\varepsilon)$。此外，由单调收敛（先对每个 $x$ 关于 $u$，再关于 $x$），$\varepsilon\mapsto\|G^\varepsilon\|_{L^p(\R^n;E_\varepsilon)}^p$ 单调不减，且当 $\varepsilon\to0$ 时极限为 $[h]_{\mathrm{unw},p,q}^p$。

\emph{(b) 第 1 步（$G^\varepsilon$ 属于 $L^p$ 且 Bochner 可测）}. 对 $M>0$ 令 $T_M(t):=\min(M,\max(-M,t))$ 为 $1$-Lipschitz 截断，并令
\[
  G^\varepsilon_M(x)(u):=\big(T_M(h)(x)-T_M(h)(x-u)\big)\mathbf{1}_{\{|u|\ge\varepsilon\}}.
\]
由于 $|T_M(a)-T_M(b)|\le|a-b|$ 且 $|T_M(t)|\le\min(|t|,M)$，对一切 $x,u$ 有 $|G^\varepsilon_M(x)(u)|\le|G^\varepsilon(x)(u)|$，且 $G^\varepsilon_M(x)\in E_\varepsilon$，并有
\[
  \|G^\varepsilon_M(x)\|_{E_\varepsilon}\le2M\,\mu_\varepsilon(\R^n)^{1/q}<\infty .
\]
我们证明 $x\mapsto G^\varepsilon_M(x)$ 是 Bochner 可测的。由于 $E_\varepsilon$ 可分，由 Pettis 定理，只需对每个 $\phi\in L^{q'}(\mu_\varepsilon)$（$1/q+1/q'=1$）证明 $x\mapsto\langle G^\varepsilon_M(x),\phi\rangle$ 可测。对这样的 $\phi$，
\[
  \langle G^\varepsilon_M(x),\phi\rangle=T_M(h)(x)\int\phi\,d\mu_\varepsilon-\int T_M(h)(x-u)\phi(u)\,\mu_\varepsilon(du).
\]
第一项可测。对第二项，固定 $R>0$：函数 $(x,u)\mapsto T_M(h)(x-u)\phi(u)$ 联合可测，且
\[
  \int_{|x|\le R}\int_{\R^n}|T_M(h)(x-u)|\,|\phi(u)|\,\mu_\varepsilon(du)\,dx
  \le M\,|B_R|\,\|\phi\|_{L^1(\mu_\varepsilon)}
  \le M\,|B_R|\,\mu_\varepsilon(\R^n)^{1/q}\|\phi\|_{L^{q'}(\mu_\varepsilon)}<\infty,
\]
其中对 $(\phi,1)$ 以指数 $q'$ 与 $q$ 应用了 H\"older 不等式，故 $\mu_\varepsilon(\R^n)$ 上的指数为 $1/q$。于是 Fubini 定理给出 $x\mapsto\int T_M(h)(x-u)\phi(u)\,\mu_\varepsilon(du)$ 在每个球上可测，从而在 $\R^n$ 上可测。因此 $G^\varepsilon_M$ 弱可测，且其像落在可分空间 $E_\varepsilon$ 的某个球内，故本质上是可分值的，由 Pettis 定理得 Bochner 可测。

其次，由构造，当 $M\to\infty$ 时对每个 $t$ 有 $T_M(h)(t)\to h(t)$，故对一切 $(x,u)$ 有 $G^\varepsilon_M(x)(u)\to G^\varepsilon(x)(u)$。对几乎每个 $x$（即满足 $F(x)<\infty$ 的那些 $x$）有 $|G^\varepsilon_M(x)(u)-G^\varepsilon(x)(u)|\le2|G^\varepsilon(x)(u)|$，且 $u\mapsto G^\varepsilon(x)(u)$ 属于 $L^q(\mu_\varepsilon)$，故关于 $\mu_\varepsilon$ 的控制收敛给出 $\|G^\varepsilon_M(x)-G^\varepsilon(x)\|_{E_\varepsilon}\to0$。此外
\[
  \|G^\varepsilon_M(x)-G^\varepsilon(x)\|_{E_\varepsilon}^p\le2^pF(x)^p,
\]
且 $x\mapsto F(x)^p=(F(x)^q)^{p/q}$ 属于 $L^1(dx)$，因为 $\int_{\R^n}(F(x)^q)^{p/q}\,dx=[h]_{\mathrm{unw},p,q}^p=\int_{\R^n}F(x)^p\,dx<\infty$；故关于 $x$ 的控制收敛给出 $\|G^\varepsilon_M-G^\varepsilon\|_{L^p(\R^n;E_\varepsilon)}\to0$。因此 $G^\varepsilon$ 是 Bochner 可测函数的 $L^p(\R^n;E_\varepsilon)$-极限，从而 Bochner 可测，且
\[
  \|G^\varepsilon\|_{L^p(\R^n;E_\varepsilon)}^p\le[h]_{\mathrm{unw},p,q}^p<\infty .
\]

\emph{(b) 第 2 步（识别与小 $z$ 估计）}. 对 $z\in\R^n$ 与 $x,y\in\R^n$ 记 $u:=x-y$，则 $dy=du$ 且 $|x-y|=|u|$。于是
\begin{align*}
  (\tau_zh-h)(x)-(\tau_zh-h)(y)
  &=\big[h(x-z)-h(x)\big]-\big[h(x-u-z)-h(x-u)\big]\\
  &=\big[h(x-z)-h(x-z-u)\big]-\big[h(x)-h(x-u)\big]\\
  &=G(x-z)(u)-G(x)(u),
\end{align*}
且以 $G^\varepsilon$ 代替 $G$ 的同一恒等式也成立。因此
\[
  [\tau_zh-h]_{\mathrm{unw},p,q}^p=\int_{\R^n}\Big(\int_{\R^n}|G(x-z)(u)-G(x)(u)|^q\,|u|^{-n-sq}\,du\Big)^{p/q}dx.
\]

现在固定 $\varepsilon>0$，把内层积分按 $|u|\ge\varepsilon$ 与 $|u|<\varepsilon$ 拆开。

(i) $|u|\ge\varepsilon$ 的部分等于 $\|\tau_zG^\varepsilon-G^\varepsilon\|_{L^p(\R^n;E_\varepsilon)}^p$；对 $G^\varepsilon\in L^p(\R^n;E_\varepsilon)$ 应用部分 (a) 可知当 $z\to0$ 时它趋于 $0$。

(ii) 对 $|u|<\varepsilon$ 的部分，
\[
  |G(x-z)(u)-G(x)(u)|^q\le2^{q-1}\big(|G(x-z)(u)|^q+|G(x)(u)|^q\big);
\]
将其取 $p/q$ 次幂并对 $x$ 积分，第一项经变量替换 $x\mapsto x+z$ 后与第二项取值相同；故该部分不超过 $C\,b(\varepsilon)$，其中 $C$ 只依赖于 $p,q$，而
\[
  b(\varepsilon):=\int_{\R^n}\Big(\int_{|u|<\varepsilon}|G(x)(u)|^q\,|u|^{-n-sq}\,du\Big)^{p/q}dx .
\]
这里须区分两类截断的方向。$\{|u|<\varepsilon\}$ 随 $\varepsilon\downarrow0$ \emph{收缩}，故 $b(\varepsilon)$ 关于 $\varepsilon$ 单调\emph{不增}；因此不能引用递增型单调收敛定理，而应使用\textbf{控制收敛定理}。具体地，令
\[
  \Phi_\varepsilon(x):=\Big(\int_{|u|<\varepsilon}|G(x)(u)|^q\,|u|^{-n-sq}\,du\Big)^{p/q},
\]
则 $0\le\Phi_\varepsilon\le\Phi:=F^{\,p}$，而 $\int_{\R^n}\Phi\,dx=[h]_{\mathrm{unw},p,q}^p<\infty$ 给出支配函数 $\Phi\in L^1(\R^n)$；又对每个满足 $F(x)<\infty$ 的 $x$（几乎处处的 $x$），$\Phi_\varepsilon(x)\downarrow0$ 当 $\varepsilon\to0$。故由控制收敛定理 $b(\varepsilon)=\int_{\R^n}\Phi_\varepsilon\,dx\to0$。因此对每个 $\varepsilon>0$ 有
\[
  \limsup_{z\to0}\,[\tau_zh-h]_{\mathrm{unw},p,q}^p\le C\,b(\varepsilon);
\]
令 $\varepsilon\to0$ 即得结论。
\end{proof}

\begin{remark}[直觉]
(a) 是 Bochner $L^p$ 空间上平移连续性的标准事实。(b) 把混合 Gagliardo 半范数重写为 $x\mapsto\big(u\mapsto h(x)-h(x-u)\big)$ 的 $L^p$ 范数。正确的结构不是单个 $L^q$ 空间，因为 $|u|^{-n-sq}du$ 只是 $\sigma$-有限的（在 $u=0$ 处发散）；办法是先截断奇点，取 $\mu_\varepsilon(du)=|u|^{-n-sq}\mathbf{1}_{\{|u|\ge\varepsilon\}}du$，在 $L^q(\mu_\varepsilon)$ 中施用 Bochner 理论，再令 $\varepsilon\to0$ 并用单调收敛。两个技术要点：逼近序列必须取 $1$-Lipschitz 的截断 $T_M(h)$ 而不是阈值截断，以保证 $|G^\varepsilon_M|\le|G^\varepsilon|$；而 $G(x)\in L^q$ 的几乎处处成员性来自 $F(x)<\infty$ a.e.，后者由 $[h]^p<\infty$ 推出，而不是来自某种乘积形式的 Tonelli 积分。
\end{remark}

%==================================================================
\subsection{技术工具 II：混合范数加权空间的可积性（局部 \texorpdfstring{$L^q$}{Lq} 与情形 (I)）}
%==================================================================

\begin{lemma}[混合范数加权空间的可积性：局部 $L^q$ 恒成立，情形 (I) 下关于加权测度的 $L^q$]
设 $n\ge1$，$\Om$ 是 $\R^n$ 中的非空有界区域，$d(x)=\dist(x,\partial\Om)$（$x\in\Om$），$s\in(0,1)$，$p,q\in[1,\infty)$，$a>0$，并对 Borel 集 $E\subset\Om$ 令 $\nuA(E):=\int_E d(y)^{-aq}\,dy$。设 $\Fsp$ 为 CLWZ 公式 (1.1)--(1.2) 的加权内蕴 Triebel--Lizorkin 空间，即使得
\begin{equation}\label{eq:sec6-f6-seminorm}
  [f]^p:=\int_\Om\Big(\int_\Om|f(x)-f(y)|^q\,|x-y|^{-n-sq}\,\nuA(dy)\Big)^{p/q}dx<\infty
\end{equation}
的 $f\in L^p(\Om)$ 之集（内层积分是关于 $\nuA$ 的 $L^q$ 范数）。则：
\begin{enumerate}[label=(\roman*)]
  \item 对每个 $a>0$，$\Fsp$ 含于 $L^q_{\mathrm{loc}}(\Om)$（不需要关于权的任何可积性假设）。
  \item 若此外还有 $\nuA(\Om)=\int_\Om d(y)^{-aq}\,dy<\infty$（情形 (I)），则 $\Fsp$ 含于 $L^q(\Om,\nuA)$：$\int_\Om|f|^q\,d^{-aq}\,dy<\infty$。更确切地说，令 $D:=\diam(\Om)\in(0,\infty)$，$J(x):=\int_\Om|f(x)-f(y)|^q|x-y|^{-n-sq}\,\nuA(dy)$；则对每个 $f\in\Fsp$ 与每个满足 $|f(x_0)|<\infty$ 且 $J(x_0)<\infty$ 的 $x_0\in\Om$（这样的 $x_0$ 之集在 $\Om$ 中具有全 Lebesgue 测度），
  \[
    \int_\Om|f(y)|^q\,\nuA(dy)\le2^{q-1}\big(|f(x_0)|^q\nuA(\Om)+D^{n+sq}J(x_0)\big)<\infty .
  \]
\end{enumerate}
\end{lemma}

\begin{proof}
令 $D:=\diam(\Om)$，故 $0<D<\infty$，并令 $c_0:=D^{-aq}>0$。由于对 $y\in\Om$ 有 $0<d(y)\le D$，故 $\nuA(dy)=d(y)^{-aq}dy\ge c_0\,dy$，权 $d^{-aq}$ 在开集 $\Om$ 上连续且为正，且 $\nuA$ 是 $\Om$ 上的非负 Borel 测度。

设 $f\in\Fsp$。则 $f\in L^p(\Om)$，映射 $(x,y)\mapsto|f(x)-f(y)|^q|x-y|^{-n-sq}$ 在对角线外联合可测，且由 Tonelli 定理 $J$ 非负可测。由于 $\int_\Om J^{p/q}dx=[f]^p<\infty$ 且 $p/q>0$，对 Lebesgue 几乎每个 $x\in\Om$ 有 $J(x)<\infty$。又因 $f\in L^p(\Om)$，几乎处处有 $|f(x)|<\infty$。故集合
\[
  G:=\{x\in\Om:J(x)<\infty\ \text{与}\ |f(x)|<\infty\}
\]
具有全 Lebesgue 测度；特别地，由于 $\Om$ 是开集，$G$ 在 $\Om$ 中稠密。

\emph{(i) 的证明}. 固定 $x_0\in G$ 与 $0<r<d(x_0)/2$，令 $B:=B(x_0,r)\cap\Om$。对 $y\in B$ 有 $|x_0-y|\le r$，故 $|x_0-y|^{-n-sq}\ge r^{-n-sq}$。利用 $\nuA\ge c_0\,dy$、$|a+b|^q\le2^{q-1}(|a|^q+|b|^q)$ 以及 $f(x_0)$ 的有限性，得
\begin{align*}
  \int_B|f(y)|^q\,dy
  &\le2^{q-1}\Big(\int_B|f(x_0)-f(y)|^q\,dy+r^n|f(x_0)|^q\Big)\\
  &\le2^{q-1}\Big(r^{n+sq}c_0^{-1}\int_B|f(x_0)-f(y)|^q\,|x_0-y|^{-n-sq}\,\nuA(dy)+r^n|f(x_0)|^q\Big)<\infty,
\end{align*}
因为 $r^{n+sq}c_0^{-1}J(x_0)<\infty$。因此对每个 $x_0\in G$ 与每个 $0<r<d(x_0)/2$ 有 $f\in L^q(B(x_0,r)\cap\Om)$。

现在设 $K$ 是 $\Om$ 的紧子集，固定 $z\in K$。令 $\delta_z:=d(z)>0$。由于 $G$ 稠密，可取 $x_z\in G$ 使 $|x_z-z|<\delta_z/8$。则 $d(x_z)\ge d(z)-|x_z-z|>7\delta_z/8$，故 $r_z:=\delta_z/4$ 满足 $0<r_z<d(x_z)/2$，且 $|x_z-z|<r_z$；因此 $z\in B(x_z,r_z)$。由紧性，$K$ 可被有限多个这样的球 $B(x_j,r_j)$（$x_j\in G$，$0<r_j<d(x_j)/2$）覆盖；于是
\[
  \int_K|f|^q\,dy\le\sum_j\int_{B(x_j,r_j)\cap\Om}|f|^q\,dy<\infty .
\]
因此 $f\in L^q_{\mathrm{loc}}(\Om)$。这就证明了 (i)，其中关于 $\nuA$ 没有任何假设。

\emph{(ii) 的证明}. 假设 $\nuA(\Om)<\infty$。对 $x_0\in G$ 与每个 $y\in\Om$ 有 $|x_0-y|\le D$，故 $|x_0-y|^{-n-sq}\ge D^{-n-sq}$。乘以 $|f(x_0)-f(y)|^q$ 并关于 $\nuA$ 积分得
\[
  J(x_0)\ge D^{-n-sq}\int_\Om|f(x_0)-f(y)|^q\,\nuA(dy),
\]
即 $\int_\Om|f(x_0)-f(y)|^q\,\nuA(dy)\le D^{n+sq}J(x_0)<\infty$。由 $L^q(\Om,\nuA)$ 范数的三角不等式，
\[
  \Big(\int_\Om|f(y)|^q\,\nuA(dy)\Big)^{1/q}
  \le\Big(\int_\Om|f(x_0)-f(y)|^q\,\nuA(dy)\Big)^{1/q}+|f(x_0)|\,\nuA(\Om)^{1/q}.
\]
两边取 $q$ 次幂并用 $|a+b|^q\le2^{q-1}(|a|^q+|b|^q)$ 得
\[
  \int_\Om|f(y)|^q\,\nuA(dy)\le2^{q-1}\big(D^{n+sq}J(x_0)+|f(x_0)|^q\nuA(\Om)\big)<\infty,
\]
这正是所写的估计，并证明 $f\in L^q(\Om,\nuA)$。注意两项都用到了 $\nuA(\Om)<\infty$；情形 (I) 正是在此处进入。
\end{proof}

\begin{remark}[直觉]
加权的内层积分 $J$ 对几乎每个 $x$ 有限，而在有界区域上核有下界 $D^{-n-sq}$；于是在一个「好点」$x_0$ 处 $J$ 的有限性就控制了 $f$ 的 $L^q(\nuA)$ 范数，因为 $\nuA(\Om)$ 有限。若没有 $\nuA(\Om)<\infty$，常数项 $|f(x_0)|^q\nuA(\Om)$ 就是无穷。局部地，同一个论证用 $\nuA\ge c_0\,dy$ 便免费给出 $L^q_{\mathrm{loc}}$（不需要任何情形假设）。
\end{remark}


%=====================================================================
\section{综合结论}\label{sec:concl}
%=====================================================================

\subsection{主结论汇总}

将上述 32 条结论归并，可得如下四条主结论。

\begin{theorem}[边界距离测度的二分；\S\ref{sec:part1}]\label{thm:main-dichotomy}
对任意有界域 $\Om$，$a>0$，$q\in[1,\infty)$，以下二者恰有一个成立：
\begin{enumerate}[label=\textup{(\Roman*)}]
\item $\nuA(\Om)<\infty$；
\item $\nuA(\Om)=\infty$，且此时自动地 $\nuA(\Omr)=\infty$ 对每个 $r>0$ 成立。
\end{enumerate}
若 $\Om$ 为有界 Lipschitz 域且 $aq\ge1$，则必属情形 (II)。
\end{theorem}

\begin{theorem}[情形 (II) 的完全退化；\S\ref{sec:part2}]\label{thm:main-degen}
设 $\nuA(\Om)=\infty$，$s\in(0,\infty)\setminus\N$，$k=\lfloor s\rfloor$，$p,q\in[1,\infty)$。则
\[
  \Fsp=P_k(\Om),
\]
其中 $P_k(\Om)$ 为次数不超过 $k$ 的多项式在 $\Om$ 上的限制。从而
$C^\infty(\Om)\cap\Fsp=\Fsp$（即 $F^{s,2}=\Fsp$）而
$C_c^\infty(\Om)\cap\Fsp=\{0\}$（即 $F^{s,1}=\{0\}$）。
\end{theorem}

\begin{theorem}[情形 (I) 的 Hardy 强化；\S\ref{sec:part3}--\ref{sec:part4}]\label{thm:main-strength}
设 $\Om$ 为任意非空有界域，$a>0$，$s\in(0,1)$，$p\in[1,\infty)$。则
\[
  \wt F^{s,1}_{p,p}(\Om,d^{-ap})=\wt F^{s,2}_{p,p}(\Om,d^{-ap})
  =\wt F^s_{p,p}(\Om,d^{-ap}).
\]
特别地 $C_c^\infty(\Om)$ 在 $\wt F^s_{p,p}(\Om,d^{-ap})$ 中稠密。
\end{theorem}

\begin{theorem}[临界条件与强化的失效；\S\ref{sec:part4}--\ref{sec:part5}]\label{thm:main-crit}
取 $n=1$、$\Om=(0,1)$、$d(x)=\min(x,1-x)$、$s=\tfrac12$、$p=q=2$、$a=\tfrac14$。则
$\nuA(\Om)=2\sqrt2<\infty$（属情形 (I)），且
\[
  F^{1/2,1}_{2,2}\big((0,1),d^{-1/2}\big)\ \neq\ F^{1/2}_{2,2}\big((0,1),d^{-1/2}\big)
  \ =\ F^{1/2,2}_{2,2}\big((0,1),d^{-1/2}\big).
\]
更一般地，对 $\Om=(0,1)$、权的指数 $b$，决定 $C_c^\infty$ 稠密性的临界比较是
$sp+b$ 与 $\codimA(\partial\Om)$，而\defterm{不是} $b$ 与 $1$（即 $\int_\Om d^{-b}<\infty$）的比较。
\end{theorem}

\begin{remark}[为何需要 Hardy 加强]
定理 \ref{thm:main-strength} 与定理 \ref{thm:main-crit} 并不矛盾：
后者针对未加强的 $F^s_{p,q}$，而前者针对 Hardy 加强空间。
加强之所以必要，原因是内蕴加权半范数在常数上为零、\defterm{看不见边界迹}，
而 Hardy 项看得见：常数 $1\in\Fsp$ 但 $1\notin\tiF$，
即 $\Fsp\not\subseteq\tiF$ 对一切非整数 $s$、一切 $a>0$、一切 $p,q$ 成立
（\S\ref{sec:part6}）。
\end{remark}

\subsection{事实依赖图}

圆括号内的数字为\defterm{附录 A 的序号}。
（箭头由后置指向前置；仅列出有前置者）：

\begin{center}
\footnotesize
\begin{longtable}{@{}lll@{}}
\toprule
后置事实 & 主题 & 前置事实 \\
\midrule
\endfirsthead
\toprule
后置事实 & 主题 & 前置事实 \\
\midrule
\endhead
(12) & $\sigma$-引理 & (20) \\
(27) & $\sigma$-引理（复证） & (20) \\
(6) & 情形 (II) 退化 & (20) \\
(11) & 情形 (II) 退化（复证） & (20), (12), (21) \\
(25) & 情形 (II) 反例 & (6), (20) \\
(28) & 高阶修正空间退化 & (20), (27) \\
(31) & 情形 (II) 强化与失效 & (11), (8) \\
(30) & $\R^n$ 上空间平凡 & (20), (6), (21) \\
(7) & 整数阶非退化 & (6) \\
(16) & 情形 (I) $C^\infty$ 稠密 & (20) \\
(24) & Hardy 加强稠密 & (16) \\
(18) & 修正高阶空间非退化 & (28), (14) \\
(10) & 主强化定理 & (20), (31), (24) \\
(23) & 区间高阶稠密 & (16), (6) \\
(4) & 情形 (I) 反例 & (20), (16) \\
(5) & 区间例（复核） & (20), (9), (16) \\
(22) & Hardy 空间严格更小（复核） & (9) \\
\bottomrule
\end{longtable}
\end{center}

无前置（根）事实为
(20)（二分）、(21)（多项式引理）、
(8)（单调性）、(9)（边界层下界）、
(14)（内部半范数估计）、(15)（平移连续性）、
(3)（Poincaré 障碍）、(13)、
(17)、(19)、(26)。

\subsection{尚未解决的问题}

\begin{enumerate}
\item \textbf{$p\neq q$ 的 Hardy 加强稠密性。}定理 \ref{thm:main-strength} 要求 $p=q$。
情形 (I) 的 $C^\infty$ 稠密性（$p=q$）依赖 Kijaczko 定理 19 与 Dyda--Kijaczko 定理 3。
对 $p\neq q$，事实 (26) 已确定：所需可积指数是 $aq$（内层变量 $y$ 的指数），
\defterm{不是} $ap$；且 $p\neq q$ 时二者互不蕴含。
但 Kijaczko 定理 19 只控制加权 Gagliardo 半范数部分，
\defterm{不}控制 CLWZ 范数 \eqref{eq:intrinsic} 中额外的 $L^q(\Om)$ 项。
故 $p\neq q$ 时 $\wt F^s$ 上的稠密性\textbf{仍未解决}。

\item \textbf{高阶 $s$ 的情形。}事实 (3) 证明：
对\defterm{非整数} $s>1$（即 $s\in(1,\infty)\setminus\N$）与任意 $a\in\R$，
$C_c^\infty(\Om)$ 在 $F^s$ 中\textbf{绝不}稠密。
故对未加强空间，\defterm{非整数} $s>1$ 时强化必然失败；
\defterm{整数阶 $s\in\N$ 与端点 $s=1$ 的情形本文未作证明}，不能由上述结果推出。
对 Hardy 加强空间 $\wt F^s$ 在 $s\ge1$ 时是否成立，尚无结论。
目前唯一覆盖 $s\ge1$ 的正面结果是区间上的
(23)（$\Om=(0,1)$、$p=q$），其证明用一维的
「最高阶导数积分还原」绕过 Leibniz 障碍，不能推广到高维。

\item \textbf{$a<0$ 的情形。}本文所有结论均设 $a>0$。
$a<0$ 时权在边界附近衰减，二分与退化现象的结构如何变化，尚未处理。

\item \textbf{常数显式化。}除个别估计外，本文所有 $\lesssim$ 均未给出常数，
亦未追踪对 $\Om$、$s$、$p$、$q$、$a$ 的依赖。
若需定量结果，须另行处理。

\item \textbf{更弱的边界正则性。}除 (9)（Lipschitz）与
(4)（区间，$\codimA(\partial\Om)=1$）外，
本文未处理 $(\varepsilon,\delta)$-domain 等更弱的边界假设。
\end{enumerate}

%=====================================================================
\section{验证状态与局限}\label{sec:limits}
%=====================================================================

\subsection{产出与验证流程}

本文 32 条结论由 4 个自主 worker（\texttt{high}、\texttt{high2}、\texttt{xhigh}、\texttt{xhigh2}）
在 DeepSeek Harness 上并行研究产出。每次提交须经一个独立的验证会话审核，
给出 \texttt{correct} 或 \texttt{needs\_repair}；仅前者写入事实图。
本项目共消耗 46 次以上验证会话，最终得到 32 条事实、230 条术语、35 条依赖边。

\subsection{独立复核}

由于事实 (4)（情形 (I) 反例）与 (13)
（常数不属于 Hardy 加强空间）是整套结论的支点，二者另经一次\defterm{独立重推}：
复核者被禁止读取这两条事实的正文，仅依据结论陈述与原始文献独立证明。
复核产出为 (5)、(22) 与 (17)，
其前置事实与被复核的两条\defterm{无交集}，结论与原事实一致。

\subsection{第三方审核与修订记录}\label{sec:revision}

在本稿完成之后，我们另请一个\defterm{独立的大语言模型}（非生产本稿的模型家族）对全文的 32 条结论做了逐条审核。
审核共提出三类意见，我们逐条判断后的处置如下。

\begin{enumerate}
\item \textbf{确认为错误并已修订（1 条）。}
  审核指出 \S\ref{sec:part2} 中整数阶修正空间非退化的原证明，把
  $\int_{|h|<cd(x)}|h|^{-n-(k-r)q}dh=C_2d(x)^{-(k-r)q}$ 这一\emph{发散}积分当作有限量。
  该指摘\defterm{成立}：当 $\alpha=(k-r)q>0$ 时该积分在 $|h|\to0$ 处为 $+\infty$。
  （i）\defterm{原证明的积分估计有误}，（ii）\defterm{witness 的适用范围需精确化}，两点处置如下。
  其一，原证明把发散的 $\int_{|h|<cd(x)}|h|^{-n-(k-r)q}dh$ 写成有限值，该指摘成立，其论证已删除。
  其二，我们最初进一步断言 witness $f_0=d^k$ 在 $n=1,k=1$ 时亦不属于该空间；该断言\defterm{有误}，已被随后的机器独立复核（事实 (2)）推翻：$k=1$ 时 $d\in F^{1}_{p,q}((0,1))$，因为 $d$ 在中点两侧分别为仿射映射，二阶差分在小 $|h|$ 上恒为零。正确的结论是：$d^k$ 仅在 $k=1$ 时可用，$k\ge2$ 时确实不属于（事实 (1)）。
  无论哪种情形，本文统一改用 $C_c^\infty$ 作 witness，该证明已由独立复核（事实 (29)）确认，结论（非退化）保持不变。

\item \textbf{因外部文献到位而消除（1 条）。}
  审核最初指出主强化定理依赖的 ``CLWZ 定理 1.2(i)'' 缺少可核验出处。
  在提供该论文原文并核对定理 1.2(i)、整数阶定理与相关定义后，此项疑点消除，
  外加强闭包等式可作为有明确出处的外部定理使用。

\item \textbf{判断不完整，已据此修订（1 条）。}
  首次审核即指出：平移连续性证明中把某尾项收敛归因于``单调收敛''不准确。
  我们初次复核时认为该指摘有误，理由是被截断的集合 $\{|u|\ge\varepsilon\}$ 在 $\varepsilon\downarrow0$ 时\emph{增大}、
  相应的 $E_\varepsilon$ 范数单调不减且极限 $[h]^p_{unw,p,q}<\infty$ 已知，故递增型单调收敛定理适用。

  \textbf{后续更细致的审阅表明前述判断并不完整。}该证明同时涉及\defterm{两类方向相反}的截断：
  \begin{enumerate}[label=(\alph*)]
    \item 截断空间 $L^q(\mu_\varepsilon)$ 与范数 $\|G^\varepsilon\|_{E_\varepsilon}$：集合 $\{|u|\ge\varepsilon\}$
          随 $\varepsilon\downarrow0$ 增大，范数单调不减且极限有限，此处单调收敛定理\emph{确实}适用；
    \item 小尺度尾项 $b(\varepsilon)=\int_{\R^n}\big(\int_{|u|<\varepsilon}|G(x)(u)|^q|u|^{-n-sq}du\big)^{p/q}dx$：
          集合 $\{|u|<\varepsilon\}$ 随 $\varepsilon\downarrow0$ \emph{收缩}，$b(\varepsilon)$ 单调\emph{不增}，
          此处不能用递增型单调收敛定理；原稿又把它误写成``单调不减''并援引单调收敛。
  \end{enumerate}
  我们最初的反驳只针对 (a)，却误用它为 (b) 作保。现已在 \S\ref{sec:part6} 相应处改用
  \defterm{控制收敛定理}（支配函数 $\Phi=F^p\in L^1$ 现成存在），并分别说明两类截断的不同单调方向。
  结论本身不变，但论证依据已更正。

\item \textbf{第三次审阅：四处可修补缺口的订正。}
  最近一次逐页审核（2026 年 10 月 26 日）在确认主要结论不变的前提下，指出四处需要订正之处，均已落实：
  \begin{enumerate}[label=(\arabic*)]
    \item 评注 \ref{rem:int-counter} 中 $k\ge2$ 情形的下界指数误写为 $\lambda^{1-kq}$；
          正确值为 $\lambda^{-(k-1)q}$（因 $\int_\lambda^{2\lambda}h^{-1-kq}dh\asymp\lambda^{-kq}$）。结论不变。
    \item 整数阶非退化定理的证明中，原稿直接使用了 $d(x)\ge\delta$，但 $\Delta_h^m\varphi(x)\ne0$ 时 $x$ 未必在 $\supp\varphi$ 内；
          现补入\defterm{支集传播}论证（$d(x)\ge\tfrac23\delta$）并把 $h$ 积分按 $|h|\le1$ 与 $|h|>1$ 两段分别估计。
    \item 主强化定理的陈述中 $q$ 未被量化，现明确声明 $q:=p$。
    \item 临界条件定理中子区间非空性的理由原为``由区间例子可知''，现改为边界余维数性质：
          Lipschitz 域的 $\codimAl=\codimAu=1$，故所论区间为 $(\max\{0,1-sp\},1)$，非空。
  \end{enumerate}
\item \textbf{第四次复核：一处实质缺口与两处澄清。}
  最近一次复核（2026 年 10 月 26 日）在确认其余各项已修复的同时，指出：
  \begin{enumerate}[label=(\arabic*)]
    \item 主强化定理所依赖的尺度选择原先只处理了两个 Hardy 项与 $L^p$ 项，
          \defterm{漏掉了原加权分数阶半范数}这一项；而“存在一组尺度使半范数收敛”不能推出
          “任意更小的逐立方体尺度仍收敛”（该半范数关于各尺度坐标并不单调）。
          现已在统一选择中补入第四项阈值 $\theta^{(1)}_{n,k}$，并以 (H$_{\mathrm{quant}}$) 的形式
          明确声明所用的是所引定理证明中的\defterm{定量}内容；
    \item 外部定理的结论原写作“对对角选取 $\eta_k(Q_n)<1/k$ 有收敛”，已改为
          \defterm{存在}一组满足该上界的尺度使其收敛；
    \item $(H_{\infty})$ 原称对有界 Lipschitz 域“或一致区域”成立，已收窄为仅对有界 Lipschitz 域，
          并注明一般一致区域缺少此处所需的边界层下界论证。
  \end{enumerate}

\item \textbf{第五次复核：把定量借用降为可核查的推导。}
  再次复核确认核心缺口已原则上修复，并指出两处可改进：
  \begin{enumerate}[label=(\arabic*)]
    \item \emph{原加权分数阶半范数的阈值应给出推导，而不只是声明引自外部定理。}
          现已在正文写出所引定理证明中的估计
          $[P^{\eta}f-f]_{F^s}^p\le M\sum_n\int\|\tau_{\eta(Q_n)u}\phi_n-\phi_n\|_{L^p}^p\,h(u)\,du$，
          并据此\defterm{逐项分配误差预算} $2^{-n}/(Mk)$、由平移在 $L^p$ 中的强连续性导出阈值
          $\theta^{(1)}_{n,k}$，从而使 $[P^{\eta}f-f]_{F^s}^p\le1/k$ 成为\defterm{结论}而不再是一条外加假设。
          本文只使用该估计的形状，不复述所引定理的证明。
    \item \emph{定理 7.6 证明中``非零多项式因为模在无穷远处增长``的措辞有误。}
          该理由对非零\defterm{常数}多项式不成立。现已改为分两种情形：
          常数多项式因 $\R^n$ 的测度无穷而不属于 $L^p(\R^n)$；
          次数至少为 $1$ 的非零多项式在某个正测度锥形区域内按多项式增长，故亦不属于 $L^p(\R^n)$。
          结论（$P=0$ 从而空间为 $\{0\}$）不变。
  \end{enumerate}

\item \textbf{第六次复核：改正上一轮引入的估计简化与指数记账。}
  本轮复核对照所引文献后指出，上一轮写下的半范数估计被\defterm{错误简化}为单变量形式。已作三处订正：
  \begin{enumerate}[label=(\arabic*)]
    \item 恢复所引定理证明中的\defterm{二变量加权估计}：以
          $g_n(x,y)=\dfrac{f(x)\psi_n(x)-f(y)\psi_n(y)}{|x-y|^{n/p+s}}\mathbf 1_{\Om\times\Om}(x,y)$
          为对象，估计中含奇异核、指示函数与权 $w\times v$；并明确指出该范数不是普通 $L^p(\Om)$ 范数，
          普通平移连续性无法控制它。
    \item 加权平移连续性原拟用“$C_c$ 在 $L^p(\omega)$ 中稠密”论证，但该论证对\defterm{非平移不变}的
          $\omega=w\times v$ 无效（$\|\tau_{tu}(G-G')\|_{L^p(\omega)}\ne\|G-G'\|_{L^p(\omega)}$）。
          现改用\defterm{紧支性 + 权的局部可比性}：在 Whitney 立方体附近的紧集上有
          $\omega\asymp dx$，从而把加权平移模控制为 Lebesgue 平移模，再对 $|u|\le1$ 取一致估计。
    \item 误差预算由 $2^{-n}/(Mk)$ 改为 $2^{-n}/(Mk^{p})$，使开 $p$ 次方后仍得 $1/k$
          （原先由 $[\,\cdot\,]^p\le1/k$ 直接写 $[\,\cdot\,]\le1/k$，在 $p>1$ 时不成立）。
  \end{enumerate}

\item \textbf{第七次复核：对照文献原文逐句核对并定稿。}
  复核指出上一轮的半范数估计被简化为单变量形式。我们\defterm{取得并核对了 arXiv:2110.04012v2 原文}，
  其定理 19 的证明确实使用二变量差商
  $g_n(x,y)=\dfrac{f(x)\psi_n(x)-f(y)\psi_n(y)}{|x-y|^{n/p+s}}\mathbf 1_{\Om\times\Om}(x,y)$
  与二变量加权估计，常数 $M=12^{n(p-1)}$。据此：
  \begin{enumerate}[label=(\arabic*)]
    \item 恢复二变量估计并\defterm{逐字对齐}原文（含 $M=12^{n(p-1)}$、$\supp h=B(0,1)$ 与算子约束
          $\eta(Q)<2\varepsilon l(Q)$）；
    \item 补充说明原文对 $\|\tau_tg_n-g_n\|^p$ 的 $I_1+I_2+I_3$ 三分解及其归结；
          本文改用\defterm{紧支 + 权的局部可比性}给出自足的平移连续性论证，并注明两种路线的一致性；
    \item 加权平移连续性原拟用稠密性论证，该论证对非平移不变权\defterm{无效}，已改为可比性论证；
    \item 误差预算改为 $2^{-n}/(Mk^{p})$，使开 $p$ 次方后仍得 $1/k$；
    \item 补入\defterm{参考文献}，给出两篇文献的完整书目信息与逐条引用位置。
  \end{enumerate}

\item \textbf{第八次复核：订正 $g_n$ 的支集并恢复三分解。}
  复核指出上一轮对二变量差商支集的判断有误。该指摘\defterm{成立}：取 $x\in\supp\psi_n$、
  $y\in\Om\setminus Q_n^{**}$ 时 $\psi_n(y)=0$ 但 $g_n(x,y)\ne0$，正确支集为
  $(Q_n^{**}\times\Om)\cup(\Om\times Q_n^{**})$，且 $g_n$ 也不是全局差商的限制。
  据此：
  \begin{enumerate}[label=(\arabic*)]
    \item 删除“只在 $Q_n^{**}\times Q_n^{**}$ 上比较权”的论证（该论证不成立，因为交叉区域里
          $v=d^{-ap}$ 可无界）；
    \item 依 [\ref{ref:K}] 定理 19 与 [\ref{ref:DK}] 定理 12 恢复\defterm{$I_1+I_2+I_3$ 三分解}：
          局部项用 $C_n=\sup_{Q_n^{**}}w\cdot\sup_{Q_n^{**}}v$ 控制，
          交叉项用引理 11 的 $|x-y|\ge c|y-x_n|$ 因子化，并把权的作用归结为两个
          \defterm{有限}常数 $D_n^{(v)},D_n^{(w)}$（其有限性用本节假设 $\int_\Om d^{-ap}<\infty$ 与 $w=1$ 验证）；
    \item 阈值改由\defterm{Lebesgue} 平移连续性给出（不涉及非平移不变的 $w\times v$）；
    \item 统一尺度中补入算子约束 $2\varepsilon l(Q_n)$，使 $\eta_k(Q_n)\le\varepsilon l(Q_n)<2\varepsilon l(Q_n)$。
  \end{enumerate}

\item \textbf{第九次复核：订正一处估计与三处符号冲突。}
  复核指出两处问题，均\defterm{成立}，且都源于新加入的证明文本：
  \begin{enumerate}[label=(\arabic*)]
    \item 估计 $A_{\mathrm{II}}$ 时把因子 $|Q_n^{*}|$ 同时加到了两项上。由
          $|f(x)-f(y)|^p\le2^{p-1}(|f(x)|^p+|f(y)|^p)$ 只能得
          $\int_{Q_n^*}\!\int_\Om|f(x)-f(y)|^p\le2^{p-1}(|\Om|+|Q_n^{*}|)\|f\|^p_{L^p}$；
          已按此订正。所得界依赖于 $n$，但每个固定 $n$ 均有限，这已足够（阈值本就允许依赖 $n,k$）。
    \item 符号 $C_n$ 被\defterm{重复定义}：第 30 页的 Hardy 权常数 $\asymp l(Q_n)^{-(s+a)p}$ 与
          第 32 页控制 $I_1$ 的局部常数 $\asymp l(Q_n)^{-ap}$ 同名，导致所选阈值在文字上
          控制不了式 (3.8) 的系数。现已拆分为
          $H_n,H'_n$（Hardy 权常数）、$A_n$（局部常数）、$D_n$（交叉项常数），
          并把标量 $\int|u_n(\eta)|^p$ 由 $g_n(\eta)$ 改名为 $e_n(\eta)$（原 $g_n$ 已是二变量差商）。
          两个 Hardy 阈值现明确写为 $H_n e_n(\eta)\le 2^{-n}/(C(N)k^{p})$ 与
          $H'_n e_n(\eta)\le 2^{-n}/(C(N)k^{p})$，直接控制式 (3.8) 与 $w_Q$ 项的系数。
  \end{enumerate}

\end{enumerate}

\noindent
由此，本稿相对于初次成稿的实质更动是：\textbf{删除一处错误积分论证，更换其 witness 并重新证明}；
其余 27 条结论的数学内容未因审核而变动。

\subsection{必须声明的局限}

\begin{enumerate}
\item \textbf{验证器是 LLM，不是形式化证明助手。}
它检查论证的合理性、跳步与隐含假设，但不做 Lean/Coq 式机器证明。
\textbf{通过验证不等于数学真理。}
本文全部结论\textbf{未经人类同行评审}。

\item \textbf{常数未显式化。}见前节第 4 条。

\item \textbf{独立复核只覆盖两条。}其余 26 条仅经过一次验证器审核。

\item \textbf{$p\neq q$ 与 $s\ge1$ 的关键区间尚缺。}
不得把本文结论外推到这些范围。

\item \textbf{部分事实的独立性有限。}
事实 (6) 与 (11)、
(12) 与 (27) 分别为同一结论的两份独立证明，
但后者可能已通过共享的全局内存了解前者的结论。
\end{enumerate}

\subsection*{给引用者的建议}

把本文视为\defterm{研究草稿与路线图}。若需引用，应优先人工复核以下四条：
(10)（主强化定理）、(4)（情形 (I) 反例）、
(17)（临界条件机制）、(13)（Hardy 项非冗余）。
其中 (10) 的正确性单向依赖于 CLWZ 定理 1.2(i) 的逐字表述，
该表述需与原文核对。

\bigskip
\begin{center}
\small
数据来源：Danus 事实图 \texttt{cx-density-weighted-tl}\\
32 条事实 \textbullet\ 230 条术语 \textbullet\ 35 条依赖边\\
生产者：\texttt{high}（13）、\texttt{xhigh}（8）、\texttt{high2}（7）、\texttt{xhigh2}（4）\\
\textbf{全部结论由 LLM 验证器判定，未经形式化验证或人类同行评审}
\end{center}

\newpage
\section*{参考文献}
\addcontentsline{toc}{section}{参考文献}

\begin{enumerate}[label={[\arabic*]},leftmargin=2.2em,itemsep=2pt]
  \item\label{ref:K} M.~Kijaczko, \emph{Fractional Sobolev spaces with power weights},
        arXiv:2110.04012v2 [math.AP] (2021).
        （期刊版：\emph{Annali dell'Universit\`a di Ferrara}, DOI 10.2422/2036-2145.202112\_002。）
        本文引用其定理 2（全半范数与截断半范数可比）、定理 3(I)(III)、定理 4、定理 19（$C^\infty\cap\wt W$ 的密度）
        及其证明中的二变量估计；算子 $P^\eta$ 的定义见该文 \S 2.2，常数 $M=12^{n(p-1)}$ 见该文定理 19 的证明。

  \item\label{ref:DK} B.~Dyda, M.~Kijaczko, \emph{On density of smooth functions in weighted fractional Sobolev spaces},
        arXiv:2009.00077 [math.FA]。
        本文引用其定理 3（Whitney 磨光在 $L^p$ 中的收敛）；Kijaczko 定理 19 证明中的 $I_1,I_2,I_3$ 估计亦归结于该文定理 12 的证明。

  \item\label{ref:CLWZ} CLWZ（待引）, \emph{加权 Triebel--Lizorkin 空间中的稠密性}（研究手稿）。
        本文引用其定义 1.1--1.6（加权内蕴范数、Hardy 加强空间、修正整数阶空间）、定理 1.2(i)（低阶 Hardy 加强闭包等式）
        与命题 1.12。
\end{enumerate}

\newpage
\section*{附录 A：结论与事实图标识对照}
\addcontentsline{toc}{section}{附录 A：结论与事实图标识对照}

\noindent 本文每条结论都对应事实图 \texttt{cx-density-weighted-tl} 中一个内容寻址（哈希命名）的记录。
为不干扰正文的数学阅读，正文一律使用规范的定理、引理编号，并以“附录 A”标注对应关系；
\S 8.2 的依赖图与 \S 9.3 中用圆括号给出的数字，即为本附录的序号。
下表按\defterm{事实标识的字典序}编号。

\begin{footnotesize}
\begin{longtable}{@{}p{1.0cm}p{4.3cm}p{8.1cm}@{}}
\toprule
序号 & 事实标识 & 主题 \\
\midrule
\endfirsthead
\toprule
序号 & 事实标识 & 主题 \\
\midrule
\endhead
\bottomrule
\endfoot
1 & \texttt{0ce5988234f5dee1} & THE WITNESS f 0(x) d(x) k IS NOT IN F k p,q (Omega) WHEN k 2 EXPLICIT D \\
2 & \texttt{0ebddc640a94b546} & FOR k 1 THE DISTANCE FUNCTION d ITSELF BELONGS TO F 1 p,q ((0,1)), SO THE k \\
3 & \texttt{10798b65b8d848a0} & THE POINCARE OBSTRUCTION FOR s 1, COMPACTLY SUPPORTED SMOOTH FUNCTIONS ARE \\
4 & \texttt{1304371856329ed6} & 情形 (I) 反例 \\
5 & \texttt{1ec46500fa263850} & 区间例 复核 \\
6 & \texttt{237af5d532395487} & 情形 (II) 退化 \\
7 & \texttt{2fc46552e3c35887} & 整数阶非退化 \\
8 & \texttt{3f558528dee6c786} & MONOTONICITY OF DIVERGENCE OF DISTANCE-WEIGHT INTEGRALS \\
9 & \texttt{44341d37c4dd3b94} & LINEAR LOWER BOUND FOR BOUNDARY LAYERS OF A LIPSCHITZ DOMAIN, AND NON-INTEGRAB \\
10 & \texttt{4c2d28e56f14f811} & 主强化定理 \\
11 & \texttt{5854a0ed976146d3} & 情形 (II) 退化 复证 \\
12 & \texttt{6da1d5ac382868e7} & SIGMA-LEMMA (top derivatives are constant when the weight is non-integrable) \\
13 & \texttt{7432f881124464bf} & CONSTANTS LIE IN F s p,q BUT NOT IN THE HARDY-AUGMENTED SPACE THE INTRINSIC \\
14 & \texttt{770e5e88a6a0a3f4} & INTERIOR WEIGHTED-SEMINORM ESTIMATE AND MOLLIFICATION CONVERGENCE IN THE INTEG \\
15 & \texttt{776da31b10113934} & TRANSLATION CONTINUITY IN VECTOR-VALUED L p AND IN THE UNWEIGHTED MIXED GAGLIA \\
16 & \texttt{818859edc858ed7d} & CASE (I) C-INFINITY DENSITY FOR p q AND s IN (0,1) \\
17 & \texttt{83625cdf3c7334ee} & WHY THE CLWZ HARDY TERMS ARE NOT REDUNDANT, AND WHY THE CRITICAL EXPONENT IS s \\
18 & \texttt{858348dbbd70be5d} & 修正高阶空间非退化 \\
19 & \texttt{94aa8a0dd29acb5e} & INTEGRABILITY OF THE MIXED-NORM WEIGHTED SPACE LOCAL L q ALWAYS, AND L q OF T \\
20 & \texttt{9764dcb121fb7ece} & DICHOTOMY FOR THE BOUNDARY-DISTANCE WEIGHT \\
21 & \texttt{a1ba0794b85a24ab} & POLYNOMIAL LEMMA CONSTANT TOP WEAK DERIVATIVES IMPLY f IS A POLYNOMIAL \\
22 & \texttt{a397aeabf43b52b2} & Hardy 空间严格更小 复核 \\
23 & \texttt{a663fb5f0a706ebe} & 区间高阶稠密 \\
24 & \texttt{aeb82d868d91bb82} & Hardy 加强稠密 \\
25 & \texttt{b42857020ae1b5e8} & 情形 (II) 反例 \\
26 & \texttt{ca4e36cccd4cba14} & KIJACZKO'S DENSITY CONDITION FOR THE DISTANCE WEIGHT d -aq THE INTEGRABILIT \\
27 & \texttt{ce71440005148eec} & NON-INTEGRABILITY RIGIDITY (sigma-lemma) \\
28 & \texttt{d1d044d3a9fe58ce} & 高阶修正空间退化 \\
29 & \texttt{da04613bf09ed7de} & C c infinity(Omega) IS CONTAINED IN THE INTEGER-ORDER MODIFIED SPACE F k p,q \\
30 & \texttt{eb711496ca489785} & TRIVIALITY OF THE WEIGHTED R n SPACE AND NON-EXISTENCE OF A WEIGHTED EXTENSION \\
31 & \texttt{ee7412a064416f68} & 情形 (II) 强化与失效 \\
32 & \texttt{f845b7305ca6dece} & TRUNCATION LEMMA COMPACTLY SUPPORTED ELEMENTS ARE DENSE IN THE INTEGER-ORDER \\
\end{longtable}
\end{footnotesize}


\end{document}
